\chapter{Foundations of Angular Electromagnetic Structure}
\label{ch:01}
\setcounter{section}{-1}

\section{Introduction and Scope}
\label{sec:ch01.0}

Suppose a narrow circularly polarized beam falls on a torsion-mounted wave plate: the plate can rotate even when the intensity pattern on a transverse screen looks featureless about the beam axis \cite{beth1936}. The mechanical response is not read from a photograph of phase contours; it is the integrated Maxwell stress on the material boundary. That experiment motivates the distinction pursued in this chapter: what the field \emph{looks like} in angle, what angular momentum it \emph{stores} locally, and what signed angular momentum it \emph{transports} outward through a closed surface are three different questions \cite{stratton1941,beth1936}.

Maxwell's equations determine $(\E,\H)$; the same pair carries energy through the Poynting vector, linear momentum through the Maxwell stress tensor, and angular momentum as the moment of that momentum about a chosen origin \cite{maxwell1873,poynting1884}. \emph{Angular structure} is how $(\E,\H)$ is organized under rotations of physical space. A spiral phase factor $e^{\jj m\phi}$ labels azimuthal organization of a transverse pattern; net torque on matter follows only from $\Tavg$, traction $\mathbf{t}=\Tavg\cdot\hat{\mathbf{n}}$, and the stress-moment integrals of Subsection~\ref{sec:ch01.1.5} \cite{allen1992,stratton1941}.

Consider next a harmonic field on a material interface, a transverse mode with $e^{\jj m\phi}$, or a near-zone shell where $\Im(\mathbf{S}_\mathrm{c}\cdot\hat{\mathbf{n}})$ dominates local normal exchange \cite{stratton1941}. Active power crossing a surface element is fixed by $\Re(\mathbf{S}_\mathrm{c}\cdot\hat{\mathbf{n}})\,dS$; mechanical exchange is fixed by $\Tavg$ and $\Mavg=\mathbf{r}\times\Tavg$. Two fields can share the same spiral isophases yet differ in $\oint \Mavg\cdot\hat{\mathbf{n}}\,dS$ because the stress integrand depends on both amplitude and the Maxwell-linked phase of $\H$, not on $\arg E$ alone \cite{stratton1941,allen1992}.

\begin{figure}[t]
\centering
\VolIPhysFig{\VolIPhysThreeRadiationConsequences}
\caption{Three Maxwellian consequences of radiation: linear momentum delivered to an absorber, angular momentum delivered through circular polarization, and equal counter-rotating orbital components whose signed transport cancels despite a structured field \cite{poynting1884,beth1936,allen1992}.}
\label{fig:ch1.three-radiation-consequences}
\end{figure}

Throughout Volume~I the time dependence is $\exp(+\jj\omega t)$, with angular frequency
$\omega$ in rad/s. A phasor $\E(\mathbf{r})$ is the complex amplitude whose real part
$\Re\{\E e^{\jj\omega t}\}$ is the instantaneous electric field. Cycle-averaged bilinear products use $(\cdot)^*$: for example $\langle\E\times\H\rangle=\tfrac12\Re\{\E\times\H^*\}$
\cite{stratton1941,harrington2001}.

For harmonic fields the \emph{complex Poynting vector} is
\begin{equation}
\label{eq:ch1.intro-poynting}
\mathbf{S}_\mathrm{c}=\tfrac{1}{2}\E\times\H^*,
\end{equation}
where $(\cdot)^*$ denotes complex conjugation \cite{poynting1884,stratton1941}.
Time-averaged real power flux density is $\langle\mathbf{S}\rangle=\Re(\mathbf{S}_\mathrm{c})$;
$\mathbf{S}_\mathrm{c}$ itself is \emph{not} the cycle-averaged Poynting vector.
The divergence $\divg\mathbf{S}_\mathrm{c}$ is net outward complex power per unit volume:
its real part balances impressed work on sources, and its imaginary part tracks harmonic
exchange between electric and magnetic stored energy (section~\ref{sec:ch01.1.4})
\cite{stratton1941}. Both $\mathbf{S}_\mathrm{c}$ and $\langle\mathbf{S}\rangle$ have
dimensions W\,m$^{-2}$; $\langle\mathbf{S}\rangle\cdot\hat{\mathbf{n}}\,dS$ is active
power per unit area crossing a surface element with outward normal $\hat{\mathbf{n}}$.

Linear momentum is carried by the Maxwell stress tensor \cite{maxwell1873,stratton1941}. In real time, the instantaneous tensor $\Tinst$ enters
\begin{equation}
\label{eq:ch1.intro-momentum-balance}
\frac{\partial\mathbf{g}_{\mathrm{em}}}{\partial t}+\mathbf{f}_{\mathrm{src}}
=\divg\Tinst,
\end{equation}
with $\mathbf{g}_{\mathrm{em}}=\varepsilon_0\,\E\times\B$ and $\mathbf{f}_{\mathrm{src}}=\rho\E+\Jimp\times\B$ using the source pair $(\rho,\Jimp)$ of Eqs.~\eqref{eq:ch1.maxwell-local} \cite{maxwell1873,stratton1941}. On phasor fields, cycle-averaged traction uses $\Tavg$ of Eq.~\eqref{eq:ch1.maxwell-stress} and $\mathbf{t}=\Tavg\cdot\hat{\mathbf{n}}$ \cite{stratton1941}. Let $V_m$ enclose the material whose force and torque are sought; $\hat{\mathbf{n}}$ is outward from $V_m$. At harmonic steady state, $\mathbf{F}_{\mathrm{matter}}=\oint_{\partial V_m}\Tavg\cdot\hat{\mathbf{n}}\,dS$.

The stress-moment density is $\Mavg=\mathbf{r}\times\Tavg$ (instantaneous counterpart $\Minst=\mathbf{r}\times\Tinst$). Local balances use source-exchange torque $\boldsymbol{\tau}_{\mathrm{src}}=\mathbf{r}\times\mathbf{f}_{\mathrm{src}}$; mechanical $\boldsymbol{\tau}_{\mathrm{matter}}$ equals $\boldsymbol{\tau}_{\mathrm{src}}$ only under Route~A (Subsection~\ref{sec:ch01.1.4}). The stress moment $\boldsymbol{\mathcal{T}}_J=\oint_{\partial V_m}\mathbf{r}\times(\Tavg\cdot\hat{\mathbf{n}})\,dS$ enters the torque balance; outward electromagnetic angular-momentum transport is $\boldsymbol{\Phi}_{J,\mathrm{out}}=-\boldsymbol{\mathcal{T}}_J$ (Eqs.~\eqref{eq:ch1.stress-moment-integral} and~\eqref{eq:ch1.outward-am-flux}). At steady state on a source-free enclosing surface, $\mathbf{T}_{\mathrm{matter}}=\boldsymbol{\mathcal{T}}_J=-\boldsymbol{\Phi}_{J,\mathrm{out}}$ when field angular-momentum storage in $V_m$ has zero cycle-averaged rate of change \cite{beth1936,stratton1941}.

In vacuum take
$\E=E_0 e^{-\jj k z}\hat{\mathbf{x}}$ and, from $\curl\E=-\jj\omega\mu_0\H$,
$\H=(E_0/\eta_0)\,e^{-\jj k z}\hat{\mathbf{y}}$ with
$\eta_0=\sqrt{\mu_0/\varepsilon_0}$ and $k=\omega\sqrt{\varepsilon_0\mu_0}$
\cite{stratton1941}. Then
\begin{equation}
\label{eq:ch1.intro-plane-wave-preview}
\mathbf{S}_\mathrm{c}=\tfrac12\E\times\H^*=\tfrac{|E_0|^2}{2\eta_0}\,\hat{\mathbf{z}},
\end{equation}
so the power through a disk of radius $a$ normal to $\hat{\mathbf{z}}$ is
$P=\pi a^2 |E_0|^2/(2\eta_0)$. Let the field control volume be the vacuum region
$z<0$, so that at the plane $z=0$ the outward unit normal is $\hat{\mathbf{n}}=+\hat{\mathbf{z}}$.
Momentum conservation gives radiation pressure $p_{\mathrm{abs}}=\langle S_z\rangle/c=
|E_0|^2/(2\eta_0 c)$ on a perfect absorber, directed along $+\hat{\mathbf{z}}$ into the
material \cite{stratton1941}. A perfect reflector at normal incidence receives
$p_{\mathrm{refl}}=2\langle S_z\rangle/c$. The factor two is not an empirical coefficient: a perfect reflector reverses the normal component of the momentum flux carried by the wave, so the magnitude of the momentum \emph{delivered} to the mirror is the incident flux plus the reversed outgoing flux, $|S|/c+|{-S}|/c=2|S|/c$ \cite{maxwell1873,stratton1941}. The vacuum field region $V_f$ ($z<0$) and the matter-enclosing volume $V_m$ (absorber plus any region whose mechanical reaction is required) play different roles: $V_f$ fixes the outward normal $\hat{\mathbf{n}}=+\hat{\mathbf{z}}$ at $z=0$, while $p_{\mathrm{abs}}$ is the mechanical pressure on the absorber \cite{maxwell1873,stratton1941}. With Eq.~\eqref{eq:ch1.maxwell-stress}, $\langle T_{zz}\rangle=-\varepsilon_0|E_0|^2/2$ on the vacuum side. Because the field is transversely symmetric
about $\hat{\mathbf{z}}$, $\dot J_{z,\mathrm{out}}=\mathbf{0}$ even
when $P>0$ \cite{beth1936}. Adding $e^{\jj m\phi}$ to the transverse factor alters the
phase portrait without, by itself, fixing the sign of the stress-flux integral
\cite{allen1992,stratton1941}.

A compatible impressed current $\Jimp$ on finite support, together with constitutive
response, boundary conditions, and outgoing selection, determines a Maxwell field
$(\E,\H)$ \cite{maxwell1873,stratton1941}. On an enclosing surface $\Gamma_R$ in
a predominantly radiative region, outward power is
$\oint \Re(\mathbf{S}_\mathrm{c}\cdot\hat{\mathbf{n}})\,dS$ and outward electromagnetic angular-momentum transport is
$-\oint \Mavg\cdot\hat{\mathbf{n}}\,dS=\boldsymbol{\Phi}_{J,\mathrm{out}}$ on the same shell under the stress convention of Eq.~\eqref{eq:ch1.outward-am-flux} \cite{beth1936}.
Near-zone fields can exchange substantial reactive energy through
$\Im(\mathbf{S}_\mathrm{c}\cdot\hat{\mathbf{n}})$ while the signed angular-momentum flux on
$\Gamma_R$ vanishes by symmetry---as the equal counter-rotating pair of
Subsection~\ref{sec:ch01.6.2} demonstrates \cite{allen1992}.

\begin{figure}[t]
\centering
\VolIPhysFig{\VolIPhysReactiveVsRadiative}
\caption{Bounded source on finite support, large net reactive power in the near zone (dashed shell), and outward power flux through enclosing surface $\Gamma_R$ (solid shell). Reactive shells can carry large $|\Im(\mathbf{S}_\mathrm{c}\cdot\hat{\mathbf{n}})|$ while radiative shells carry the conserved outward power \cite{poynting1884,chu1948}.}
\label{fig:ch1.master-philosophy}
\end{figure}

Figure~\ref{fig:ch1.master-philosophy} separates reactive from radiative shells.
The impressed current occupies finite support; the dashed inner shell lies where
$\Im(\mathbf{S}_\mathrm{c}\cdot\hat{\mathbf{n}})$ can dominate local normal
flux; the solid shell $\Gamma_R$ lies where
$\Re(\mathbf{S}_\mathrm{c}\cdot\hat{\mathbf{n}})$ carries the conserved outward power of the
source. The angular-momentum balance holds on any admissible closed surface; to read outward radiated angular momentum, a predominantly radiative enclosing surface is preferred because reactive near-zone exchange and storage then cease to obscure the transport interpretation \cite{chu1948,stratton1941}.

Source realizability, aperture traces, and harmonic charts are developed in Subsection~\ref{sec:ch01.0.1}; counterexamples appear in Subsections~\ref{sec:ch01.6.2} and~\ref{sec:ch01.6.4} \cite{stratton1941,harrington2001}.
\phantomsection\label{thm:ch1.realizability-hierarchy}\label{prin:ch1.rep-physics}\label{def:ch1.hierarchy-classes}\label{hier:ch1.angular-status}\label{prin:ch1.finite-geometry-angular}\label{prin:ch1.accessibility}

\begin{figure}[t]
\centering
\VolIPhysFig{\VolIPhysThreeLevelHierarchy}
\caption{Impressed current on finite support $\Omega_s$, Maxwell fields $(\E,\H)$ on enclosing surface $\Gamma_R$, and angular variation of the tangential field in a harmonic chart on $S^2$. The chart lists angular organization of an already specified field pair; it is not a substitute for impressed sources and boundary closure \cite{stratton1941}.}
\label{fig:ch1.three-level-hierarchy}
\end{figure}

\subsection{Historical lineage of field energy, momentum, and angular structure}
\label{sec:ch01.0.historical-lineage}

Faraday's experiments on induction and magnetism forced a shift from action at a
distance to \emph{fields} as the primary carriers of electric and magnetic action in
space \cite{maxwell1873}. Maxwell's 1865 paper and the 1873 \emph{Treatise} supplied the
mathematical field pair $(\E,\H)$, the displacement current that closes Amp\`ere's law in
time-dependent form, and the wave equation that predicted electromagnetic propagation at
the speed of light \cite{maxwell1873}. Once the field itself became dynamical, two
questions followed immediately: where does electromagnetic energy reside, and how does
the field exchange linear and rotational mechanical action with matter?

Maxwell answered the mechanical side within the same field picture. His stress tensor
relates bilinear products of $(\E,\H)$ to force and momentum density in the
electromagnetic medium \cite{maxwell1873,stratton1941}. Poynting \cite{poynting1884}
then identified energy transport with the flux $\mathbf{s}=\E\times\H$, giving a local
account of power crossing a surface. Hertz's experiments (1887--1888) showed that
Maxwell's predicted waves could be generated and detected in the laboratory, converting
the wave equation from theory into measurable radiation \cite{stratton1941}. Radiation-pressure
measurements on absorbing and reflecting surfaces confirmed that the same field carries
directed momentum as well as energy; on a perfect absorber the pressure is $p_{\mathrm{abs}}=|\langle\mathbf{S}\rangle|/c$, on a
perfect reflector $p_{\mathrm{refl}}=2|\langle\mathbf{S}\rangle|/c$ in vacuum \cite{stratton1941}.

The early twentieth century did not \emph{introduce} electromagnetic stress---Maxwell
had already done so---but reopened the question of momentum \emph{inside} material
media. Abraham \cite{abraham1910} and Minkowski \cite{minkowski1908} proposed different
symmetric stress--energy forms for polarizable matter; their disagreement concerns how
much momentum is stored in a dielectric, not whether vacuum radiation exerts pressure.
In source-free exterior regions both reduce to the vacuum Maxwell stress used throughout
this chapter \cite{stratton1941}.

Stratton \cite{stratton1941}, Sommerfeld \cite{sommerfeld1949}, Harrington
\cite{harrington2001}, and Jackson \cite{jackson1999} consolidated the mathematical
apparatus of boundary-value radiation problems: impressed sources, outgoing conditions,
dyadic Green functions, and surface flux integrals for power and momentum. Chu
\cite{chu1948} emphasized that finite source size and stored reactive energy bound what
a compact antenna can radiate, so far-zone patterns cannot be separated from near-zone
storage. Tai \cite{tai1994} systematized dyadic Green-function and boundary-value methods in monochromatic media; gauge-invariant helicity language and dual rotations on $(\E,\H)$ are treated separately in the field-theoretic literature \cite{belinfante1940,tai1994}.

Rotational action entered experimentally with Beth \cite{beth1936}, who measured the
mechanical torque on a quartz wave plate in circularly polarized light. The observation
confirmed that angular momentum is transported by the Maxwell field and delivered to
matter through stress at an interface, not inferred from phase portrait alone. Belinfante
\cite{belinfante1940} clarified that canonical and kinetic momentum densities differ by
the divergence of a spin current; correspondingly, local ``orbital'' and ``spin''
labels partition total angular momentum in a convention-dependent way, while closed-surface
flux integrals remain anchored to $(\E,\H)$ \cite{belinfante1940,stratton1941}.

Allen \emph{et al.} \cite{allen1992} later gave paraxial beams a countable azimuthal
index $m$ through the phase $\exp(\jj m\phi)$. That vocabulary made helical phase
structure experimentally accessible, but the measurable quantities of this chapter remain
the Poynting flux, Maxwell traction, and angular flux $\Mavg=\mathbf{r}\times\Tavg$
on a radiative shell---quantities whose history runs from Maxwell and Poynting through
Beth to the present chapter \cite{poynting1884,beth1936,allen1992}.
\subsection{Source data, boundaries, and Maxwell consistency}
\label{sec:ch01.0.1}

Every flux integral in this chapter refers to a Maxwell field pair $(\E,\H)$ fixed
by impressed sources, constitutive law, boundary data, and---when radiation is
claimed---an outgoing condition \cite{stratton1941,maxwell1873}. A tangential field
prescribed tangential pattern on an aperture is not, by itself, such a specification
\cite{harrington2001}.

On a PEC-backed circular aperture, a charge-continuous impressed current $\Jimp$ closes Maxwell's equations and defines radiated flux on $\Gamma_R$; prescribing $\E_t$ alone on the rim without a matching source may satisfy a local boundary relation yet fail to tie Poynting and torque integrals to any synthesizable excitation \cite{stratton1941,harrington2001}. Figure~\ref{fig:ch1.wp-matched-aperture} contrasts the two configurations. Stratton \cite{stratton1941} fixed the physical content of a radiation problem: spatial domain \(\Omega\), boundary \(\Gamma\), constitutive data \(\chi\), impressed sources \((\Jimp,\rho)\), and harmonic frequency under the \(\exp(+\jj\omega t)\) convention adopted throughout this volume.

\begin{figure}[t]
\centering
\VolIPhysFig{\VolIPhysMatchedAperture}
\caption{Same circular aperture and observation shell $\Gamma_R$: left, a charge-continuous impressed current $\Jimp$ launches a Maxwell field whose outward flux on $\Gamma_R$ is defined; right, a tangential $\E_t$ pattern is prescribed on the rim without a matching source, so flux integrals need not correspond to any realizable excitation \cite{stratton1941,harrington2001}.}
\label{fig:ch1.wp-matched-aperture}
\end{figure}

Given impressed current and charge, together with the constitutive medium, boundary conditions on $\Gamma=\partial\Omega$, and angular frequency $\omega>0$ in the $\exp(+\jj\omega t)$ convention, Maxwell's equations determine the electromagnetic field once the appropriate radiation condition is imposed \cite{maxwell1873,stratton1941}. An angular claim is meaningful only after that boundary-value problem is closed.

Inconsistent boundary traces cannot be repaired by deeper harmonic listing. The source and boundary conditions precede every transport statement \cite{stratton1941}.
\subsection{Rotation and Maxwell compatibility}
\label{sec:ch01.0.2}

Rotate a vector by angle $\alpha$ about axis $\hat{\mathbf{u}}$:
$\mathbf{v}'=R_\alpha\mathbf{v}$. A second-order tensor transforms as
$T'_{ij}=R_{ik}T_{kl}R_{jl}$, or $\mathbf{T}'=R\mathbf{T}R^\top$ in dyadic notation
\cite{stratton1941}. A vector field transforms actively as
$(T_R\mathbf{F})(\mathbf{r})=R\,\mathbf{F}(R^{-1}\mathbf{r})$: the source is rotated in
space, not merely relabelled \cite{hansen1988}.

If the impressed current on a feed is rotated through $\alpha$ while the horn, ground
plane, and medium share the same axisymmetry, does the radiated field rotate with it?
When source, medium, and boundaries share rotational symmetry, solving Maxwell's
equations after rotating $\Jimp$ yields the same physical field as rotating the
solution \cite{stratton1941,hansen1988}. An anisotropic radome or azimuthally varying $\varepsilon$ breaks that equivalence only when the full constitutive and boundary configuration fails to share the same rotation symmetry about the axis in question \cite{stratton1941}.

\subsection{Assumptions, finite energy, and domain of validity}
\label{sec:ch01.0.3}

A coax-fed pyramidal horn over a ground plane illustrates how an impressed current on a finite feed support, together with conducting boundaries and an outgoing radiation condition, defines a Maxwell boundary-value problem \cite{stratton1941}. Before Poynting, stress, or angular-flux integrals are interpreted, one must know whether the underlying field pair is regular enough on the chosen control volume for those quadratic flux densities to have finite surface integrals \cite{maxwell1873,poynting1884,stratton1941}.

Physical stored energy and mathematical field integrability are related but not identical. Under the harmonic convention \(\exp(+\jj\omega t)\), the cycle-averaged electromagnetic energy density in a \emph{linear, isotropic, lossless} medium with real, frequency-independent \(\varepsilon\) and \(\mu\) at the operating frequency is \(u=\tfrac14(\varepsilon|\E|^{2}+\mu|\H|^{2})\) \cite{stratton1941}. In dispersive or dissipative media the same phasor magnitudes do not, by themselves, define a unique local stored-energy density; dissipation appears in \(\Im(\varepsilon)\) and \(\Im(\mu)\) through the complex Poynting balance of Subsection~\ref{sec:ch01.1.4} \cite{stratton1941,harrington2001}. In vacuum, \(\varepsilon=\varepsilon_0\) and \(\mu=\mu_0\), so \(u\) has SI units J\,m\(^{-3}\). By contrast, \(|\E|^{2}\) alone carries units V\(^2\)m\(^{-2}\) and \(|\H|^{2}\) carries A\(^2\)m\(^{-2}\); adding them without constitutive weights is not a single physical density. The admissibility condition used here is therefore a \emph{weighted} field norm on a specified \emph{bounded} volume \(V\): Eq.~\eqref{eq:ch1.energy-class} certifies integrability of \(\varepsilon_0|\E|^{2}+\mu_0|\H|^{2}\) over \(V\). It does not, without further trace or smoothness hypotheses, certify classical boundary values on \(\partial V\) or the finiteness of every quadratic Poynting or stress flux on \(\partial V\) \cite{stratton1941}.

Let \(\Omega\subset\mathbb{R}^{3}\) be the analysis region with boundary \(\Gamma=\partial\Omega\). For any finite control volume \(V\subset\Omega\) that excludes the support of ideal point or filament singularities, require
\begin{equation}
\label{eq:ch1.energy-class}
\int_{V}\!\left(\varepsilon_0\,\E^*\!\cdot\E+\mu_0\,\H^*\!\cdot\H\right)\,dV<\infty.
\end{equation}
The integrand is the unweighted harmonic energy density \(4u\) in source-free vacuum; Eq.~\eqref{eq:ch1.energy-class} is therefore simultaneously a finite stored-energy test on \(V\) and an \(L^2\)-type integrability condition on the weighted field pair \cite{stratton1941}. It is imposed on the bounded volume \(V\) actually used in the calculation, not as a claim of global square-integrability over all space \cite{stratton1941}.

Outgoing \(1/r\) radiation fields are generally \emph{not} globally square-integrable on an unbounded exterior domain in the ordinary \(L^2\) sense, because their energy density decays only as \(r^{-2}\) \cite{stratton1941}. Nevertheless, time-averaged radiated power \(P_{\mathrm{rad}}\), integrated Maxwell traction \(\boldsymbol{\mathcal{T}}_J\), and angular flux \(\boldsymbol{\Phi}_{J,\mathrm{out}}\) on a finite enclosing shell \(S_R\) can remain well defined when the outgoing radiation condition selects the physical branch \cite{stratton1941}. Ideal point dipoles and filaments are retained only as explicit distributional limits of finite-width, charge-continuous sources; otherwise the regular impressed-current class of Eq.~\eqref{eq:ch1.source-charge-continuous} applies. Square-integrable \(\Jimp\) and charge continuity \(\divg\Jimp+\jj\omega\rho=0\) do not by themselves guarantee Eq.~\eqref{eq:ch1.energy-class} in the exterior region, nor do they guarantee classical normal traces of \(\Jimp\) on \(\Gamma\) or finiteness of every quadratic stress or Poynting surface integral invoked later---boundary data, outgoing selection, and the field regularity of the actual solution control those conclusions \cite{stratton1941}. Boundary data on \(\Gamma\) complete the physical problem together with the source specification \cite{maxwell1873,stratton1941}.
\phantomsection\label{eq:ch1.admissible-space}

\begin{figure}[t]
\centering
\VolIPhysFig{\VolIPhysGaussianWaist}
\caption{Schematic only: a Gaussian envelope illustrates why a finite analysis domain and enclosing shell $S$ must be chosen before interpreting flux integrals; it is not, by itself, a Maxwell solution with specified sources, boundaries, and outgoing selection \cite{allen1992,stratton1941}.}
\label{fig:ch1.wp-gaussian-waist}
\end{figure}

Finite analysis domains must enclose the radiated power before flux integrals are interpreted; waist parameters alone do not define transport on a shell \cite{allen1992,stratton1941}. Figure~\ref{fig:ch1.wp-gaussian-waist} illustrates that enclosure requirement \cite{stratton1941}.

\section{Maxwell Theory and Conservation Laws}
\label{sec:ch01.1}

Charge continuity follows from the divergence of Amp\`ere's law. Energy and momentum
follow from the curl structure of $(\E,\H)$. Angular momentum is the moment of
electromagnetic momentum carried by the Maxwell stress tensor
\cite{maxwell1873,poynting1884,beth1936}. Maxwell's equations in harmonic form are
\begin{equation}
\label{eq:ch1.maxwell-centre}
\curl \E=-\jj\omega\mu_0\H,\qquad
\curl \H=\Jimp+\jj\omega\varepsilon_0\E,
\end{equation}
together with the divergence constraints and boundary completion of section~\ref{sec:ch01.0} \cite{stratton1941}.

Helicity sectors, modal coefficients, and orbital labels are coordinate listings on radiative fields, not independent degrees of freedom that override curl--divergence constraints \cite{maxwell1873,jackson1999}. Assigning azimuthal phase to $\E$ alone without updating $\H$ through $\curl\E=-\jj\omega\mu_0\H$ violates local Maxwell compatibility \cite{stratton1941}.

\subsection{Maxwell equations: differential, integral, vacuum constitutive relations}
\label{sec:ch01.1.1}

A time-varying electric field cannot remain an isolated scalar pattern in space: Faraday's law ties its curl to $\H$, and Amp\`ere--Maxwell ties the curl of $\H$ back to $\E$ and to the sources that drive the system \cite{maxwell1873,stratton1941}. Specifying only the spatial phase of one Cartesian component---for example a forward $\exp(-\jj k z)$ factor on $E_x$---does not yet define an electromagnetic solution until the paired $\H$ satisfies both curl equations and the divergence constraints on the same domain \cite{stratton1941,jackson1999}. The local Maxwell system is therefore the starting point for every later transport, flux, and angular-momentum statement in this chapter \cite{stratton1941}.

Harmonic analysis uses the time factor $\exp(+\jj\omega t)$ in a piecewise smooth region $\Omega\subset\mathbb{R}^{3}$ with boundary $\Gamma=\partial\Omega$ \cite{stratton1941}. The phasors $\E(\mathbf{r})$ and $\H(\mathbf{r})$ then satisfy $\partial/\partial t\to +\jj\omega$. Impressed current density $\Jimp(\mathbf{r})$ and charge density $\rho(\mathbf{r})$ enter Amp\`ere--Maxwell and Gauss's electric law; vacuum permittivity $\varepsilon_0$ and permeability $\mu_0$ are SI constants linking field intensities to flux densities \cite{stratton1941,harrington2001}. Three complementary presentations of the same physics will be used: the \emph{differential} form for pointwise coupling, the \emph{integral} form for finite control surfaces and volumes, and the \emph{vacuum constitutive relations} that close the system in free space and fix propagation scales \cite{stratton1941}.

\begin{figure}[t]
\centering
\VolIPhysFig{\VolIPhysVacuumClosureMap}
\caption{Vacuum constitutive relations linking $(\E,\D)$ and $(\H,\B)$ through $\varepsilon_0$ and $\mu_0$, and the outgoing wave whose kinematics those constants determine.}
\label{fig:ch1.vacuum-closure-map}
\end{figure}
Figure~\ref{fig:ch1.vacuum-closure-map} summarizes how $\varepsilon_0$ and $\mu_0$ connect $(\E,\H)$ to $(\mathbf{D},\mathbf{B})$ and how the same constants enter wave speed, impedance, and wavenumber in later formulas \cite{maxwell1873,stratton1941}. Any update to $\E$ that is not propagated to $\H$ through Faraday's law breaks that closure and cannot represent a Maxwell field.

The differential Maxwell system on $\Omega$ reads
\begin{equation}
\label{eq:ch1.maxwell-local}
\begin{aligned}
\curl \E&=-\jj\omega\mu_0\H, &
\curl \H&=\Jimp+\jj\omega\varepsilon_0\E,\\
\divg\!\left(\varepsilon_0\E\right)&=\rho, &
\divg\!\left(\mu_0\H\right)&=0.
\end{aligned}
\end{equation}
Faraday's law states that a circulating $\E$ requires a time-harmonic $\H$ proportional to $\curl\E$. Amp\`ere--Maxwell adds the impressed current $\Jimp$ and the displacement current $\jj\omega\varepsilon_0\E$ that makes the law consistent with time-varying charge \cite{maxwell1873}. Gauss's electric law, $\divg(\varepsilon_0\E)=\rho$, specifies how much electric flux originates from charge at a point; it is not the same statement as charge continuity, which constrains how $\rho$ and $\Jimp$ vary together \cite{stratton1941}. Gauss's magnetic law, $\divg(\mu_0\H)=0$, expresses the absence of magnetic monopoles in the classical model \cite{maxwell1873}.

Charge continuity in harmonic form follows from the divergence of Amp\`ere--Maxwell. Because $\divg(\curl\H)=0$ identically,
\begin{equation}
\label{eq:ch1.charge-continuity-harmonic}
\divg\Jimp+\jj\omega\rho=0,
\end{equation}
obtained by taking $\divg$ of $\curl\H=\Jimp+\jj\omega\varepsilon_0\E$ and substituting Gauss's law \cite{maxwell1873,stratton1941}. The impressed sources are therefore not independent data: any $\rho$ paired with $\Jimp$ must satisfy Eq.~\eqref{eq:ch1.charge-continuity-harmonic} or the four rows of Eq.~\eqref{eq:ch1.maxwell-local} cannot hold simultaneously \cite{hansen1988}.

Integrating Eq.~\eqref{eq:ch1.charge-continuity-harmonic} over a control volume $V$ with boundary $\partial V$ and outward unit normal $\hat{\mathbf{n}}$ gives
\begin{equation}
\label{eq:ch1.continuity-integral}
\int_{\partial V}\Jimp\cdot\hat{\mathbf{n}}\,dS
+\jj\omega\int_V\rho\,dV=0.
\end{equation}
The surface integral is net impressed current crossing $\partial V$ in the outward direction; the volume integral is the harmonic charge content inside $V$. Under $\exp(+\jj\omega t)$, a positive imaginary contribution from $\jj\omega\rho$ corresponds to a $90^\circ$ phase relation between charge and its rate of change in the time domain; the balance in Eq.~\eqref{eq:ch1.continuity-integral} states that current leaving the volume must be accompanied by the harmonic charge variation required by Eq.~\eqref{eq:ch1.charge-continuity-harmonic} \cite{maxwell1873,stratton1941}.

Cartesian components make the same coupling visible on a single row:
\begin{equation}
\label{eq:ch1.maxwell-component-example}
\begin{aligned}
\partial_y E_z-\partial_z E_y&=-\jj\omega\mu_0 H_x,\\
\partial_z H_x-\partial_x H_z&=J_{\mathrm{imp},y}+\jj\omega\varepsilon_0 E_y,\\
\partial_x(\varepsilon_0 E_x)+\partial_y(\varepsilon_0 E_y)+\partial_z(\varepsilon_0 E_z)&=\rho.
\end{aligned}
\end{equation}
Each line is one scalar component of Eq.~\eqref{eq:ch1.maxwell-local} \cite{stratton1941}. Subscripts $x,y,z$ label fixed Cartesian axes; $J_{\mathrm{imp},y}$ is the $y$-component of $\Jimp$ \cite{harrington2001}. A patch of $E_z$ alone still demands paired $H_x$ through Faraday's row and paired $\rho$ through Gauss's row, so field maps that display only one component can misrepresent the underlying Maxwell state \cite{maxwell1873}.

The integral form of the same laws uses Stokes' theorem on surfaces $S\subset\Omega$ with boundary $\partial S$ and the divergence theorem on volumes $V\subset\Omega$:
\begin{equation}
\label{eq:ch1.maxwell-integral}
\begin{aligned}
\oint_{\partial S}\E\cdot d\boldsymbol{\ell}
&=-\jj\omega\mu_0\int_S\H\cdot\hat{\mathbf{n}}\,dS,\\
\oint_{\partial S}\H\cdot d\boldsymbol{\ell}
&=\int_S\Jimp\cdot\hat{\mathbf{n}}\,dS
+\jj\omega\varepsilon_0\int_S\E\cdot\hat{\mathbf{n}}\,dS,\\
\int_{\partial V}\varepsilon_0\E\cdot\hat{\mathbf{n}}\,dS
&=\int_V\rho\,dV,\qquad
\int_{\partial V}\mu_0\H\cdot\hat{\mathbf{n}}\,dS=0.
\end{aligned}
\end{equation}
The circulations of $\E$ and $\H$ around $\partial S$ are tied to fluxes of $\H$ and $\E$ through $S$; the normal flux of $\varepsilon_0\E$ through $\partial V$ equals the charge in $V$ \cite{stratton1941}. These statements match Eq.~\eqref{eq:ch1.maxwell-local} when $\E$, $\H$, $\Jimp$, and $\rho$ are sufficiently smooth for the classical integral theorems; singular or distributional sources require the same laws in weak or limit form \cite{stratton1941}. \phantomsection\label{prop:ch1.local-integral-equivalence}
Shrinking $V$ or $S$ to an interior point recovers the differential rows.

Vacuum constitutive relations complete the model in free space:
\begin{equation}
\label{eq:ch1.vacuum-constitutive}
\mathbf{D}=\varepsilon_0\E,\qquad
\mathbf{B}=\mu_0\H,\qquad
c=\frac{1}{\sqrt{\varepsilon_0\mu_0}}.
\end{equation}
Here $\mathbf{D}$ is electric flux density (C\,m$^{-2}$), $\mathbf{B}$ is magnetic flux density (T), and $c$ is the speed of light in vacuum (m\,s$^{-1}$) \cite{stratton1941}. Without Eq.~\eqref{eq:ch1.vacuum-constitutive}, the curl laws alone leave $(\mathbf{D},\mathbf{B})$ undetermined.

Combining Eq.~\eqref{eq:ch1.maxwell-local} with Eq.~\eqref{eq:ch1.vacuum-constitutive} and eliminating $\H$ from the curl pair yields second-order operators whose source-free solutions propagate at speed $c$. The intrinsic impedance, wavenumber, and wavelength used throughout the chapter are
\begin{equation}
\label{eq:ch1.vacuum-impedance-wavenumber}
\eta_0=\sqrt{\frac{\mu_0}{\varepsilon_0}},\qquad
k_0=\omega\sqrt{\mu_0\varepsilon_0}=\frac{\omega}{c},\qquad
\lambda_0=\frac{2\pi}{k_0}.
\end{equation}
The quantity $\eta_0$ (ohms) relates transverse $\E$ and $\H$ amplitudes in a uniform plane wave; $k_0$ (rad\,m$^{-1}$) measures phase accumulation per unit distance at angular frequency $\omega$; $\lambda_0=2\pi/k_0$ is the corresponding vacuum wavelength \cite{stratton1941}. These definitions are not separate postulates: they follow from the same $\varepsilon_0$, $\mu_0$, and $\omega$ that appear in Eq.~\eqref{eq:ch1.maxwell-local}.

To obtain the inhomogeneous wave equation for $\E$, use the vector identity
\begin{equation}
\label{eq:ch1.curl-curl-identity}
\curl(\curl\E)=\grad(\divg\E)-\nabla^2\E,
\end{equation}
valid for continuously differentiable fields \cite{stratton1941}. With $\exp(+\jj\omega t)$ phasors, $\curl\E=-\jj\omega\mu_0\H$ and $\curl\H=\Jimp+\jj\omega\varepsilon_0\E$. Applying $\curl$ to Faraday's law and substituting Amp\`ere--Maxwell,
\begin{equation}
\label{eq:ch1.wave-derivation-intermediate}
\curl(\curl\E)=-\jj\omega\mu_0\Jimp+k_0^2\E,
\end{equation}
because $k_0^2=\omega^2\mu_0\varepsilon_0$. Gauss's law gives $\divg\E=\rho/\varepsilon_0$, so Eq.~\eqref{eq:ch1.curl-curl-identity} rearranges to
\begin{equation}
\label{eq:ch1.vacuum-wave-source-driven}
\nabla^2\E+k_0^2\E=\frac{1}{\varepsilon_0}\nabla\rho+\jj\omega\mu_0\Jimp,\qquad
\nabla^2\H+k_0^2\H=-\curl\Jimp.
\end{equation}
The electric equation carries both the charge-gradient forcing $\varepsilon_0^{-1}\nabla\rho$ and the impressed-current term $+\jj\omega\mu_0\Jimp$ under the $+\jj\omega t$ convention adopted here \cite{jackson1999,stratton1941}. The magnetic companion is obtained independently: $\curl\H=\Jimp+\jj\omega\varepsilon_0\E$ gives $\curl(\curl\H)=\curl\Jimp+\jj\omega\varepsilon_0\curl\E=\curl\Jimp-k_0^2\H$, and the same curl--curl identity yields $\nabla^2\H+k_0^2\H=-\curl\Jimp$. Homogeneous regions with $\Jimp=\mathbf{0}$ and $\rho=0$ satisfy the vector Helmholtz pair
\begin{equation}
\label{eq:ch1.vacuum-helmholtz-pair}
\nabla^2\E+k_0^2\E=\mathbf{0},\qquad
\nabla^2\H+k_0^2\H=\mathbf{0},
\end{equation}
but Eq.~\eqref{eq:ch1.vacuum-helmholtz-pair} alone does not select a unique physical solution: boundary data on $\Gamma$, radiation conditions at infinity, and the impressed source specification still determine which field belongs to a given antenna or scatterer problem \cite{stratton1941}.

For a uniform plane-wave ansatz $\E=\E_0 e^{-\jj\mathbf{k}\cdot\mathbf{r}}$ with $\mathbf{k}\cdot\E_0=0$, Eq.~\eqref{eq:ch1.vacuum-helmholtz-pair} requires $|\mathbf{k}|=k_0$. Faraday's law gives $\H=(\mathbf{k}\times\E)/(\omega\mu_0)=\hat{\mathbf{k}}\times\E/\eta_0$ with $\hat{\mathbf{k}}=\mathbf{k}/|\mathbf{k}|$, and Amp\`ere--Maxwell confirms $\mathbf{k}\cdot\H=0$ \cite{stratton1941,jackson1999}. Transverse vacuum plane waves therefore satisfy
\begin{equation}
\label{eq:ch1.local-plane-wave-relations}
\H=\frac{1}{\eta_0}\,\hat{\mathbf{k}}\times\E,\qquad
\hat{\mathbf{k}}\cdot\E=0,\qquad
\hat{\mathbf{k}}\cdot\H=0,
\end{equation}
where $\eta_0=\sqrt{\mu_0/\varepsilon_0}$ \cite{stratton1941,jackson1999}. This is a \emph{local} Maxwell-compatible plane-wave sector: it does not describe, by itself, the full field of a bounded radiator, which requires outgoing selection and source integrals developed in the next subsection \cite{stratton1941}.

Take $\E=E_0 e^{-\jj k_0 z}\hat{\mathbf{x}}$ propagating in $+\hat{\mathbf{z}}$. Then $\curl\E=-\jj k_0 E_0 e^{-\jj k_0 z}\hat{\mathbf{y}}=-\jj\omega\mu_0\H$, so $\H=(E_0/\eta_0)e^{-\jj k_0 z}\hat{\mathbf{y}}$. Substituting into $\curl\H=\jj\omega\varepsilon_0\E$ (source-free) closes Amp\`ere--Maxwell with no remaining free parameters. Real-time phase surfaces satisfy $\omega t-k_0 z=\mathrm{const}$, hence $dz/dt=\omega/k_0=c>0$: the disturbance advances in $+\hat{\mathbf{z}}$, and the factor $e^{-\jj k_0 z}$ is the outgoing spatial branch for $\exp(+\jj\omega t)$ \cite{stratton1941,jackson1999}. The time-average Poynting direction is $+\hat{\mathbf{z}}$ with magnitude $|E_0|^2/(2\eta_0)$ for this polarization \cite{stratton1941}.

The same source-free physics may be written in curl-curl form:
\begin{equation}
\label{eq:ch1.vacuum-curl-second-order}
\curl\!\left(\mu_0^{-1}\curl \E\right)-\omega^2\varepsilon_0\E=\mathbf{0},\qquad
\curl\!\left(\varepsilon_0^{-1}\curl \H\right)-\omega^2\mu_0\H=\mathbf{0}.
\end{equation}
Eq.~\eqref{eq:ch1.vacuum-curl-second-order} is equivalent to Eq.~\eqref{eq:ch1.vacuum-helmholtz-pair} in vacuum and is often the preferred starting point for variational formulations with tangential boundary data \cite{stratton1941}.

Eqs.~\eqref{eq:ch1.maxwell-local}--\eqref{eq:ch1.local-plane-wave-relations} fix how $\E$ and $\H$ must balance locally for given $\Jimp$ and $\rho$, but they do not yet specify which source distributions, interface traces, and outer radiation conditions define a realizably fed structure on $\Omega$. That completion---Green-function reconstruction and the Maxwell boundary-value problem---is taken up immediately below \cite{stratton1941}.
\input{ch01-foundation-seeds-passK-bridges}

\begin{figure}[t]
\centering
\VolIPhysFig{\VolIPhysDipoleLoop}
\caption{Azimuthal filament loop of radius $b$ carrying $I_0\hat{\boldsymbol{\phi}}$: charge-continuous impressed current on compact support. Illustrative shells at $k_0 b=0.063$ (electrically small source) and $k_0 r=6.3$ ($r\approx\lambda_0$); the latter is an order-of-magnitude illustration, not by itself a proof of Fraunhofer or full far-zone asymptotics.}
\label{fig:ch1.wp-solenoidal-loop}
\end{figure}

Figure~\ref{fig:ch1.wp-solenoidal-loop} is the minimal charge-continuous source geometry: azimuthal
current on a filament loop produces dipole radiation because $\divg\Jimp=0$ on the support while
the impressed current closes the Amp\`ere row in Eq.~\eqref{eq:ch1.maxwell-local}. Replacing the
loop by a tangential $\E_t$ pattern on the same circle would satisfy a boundary display but not
the source row unless a matching $\Jimp$ is prescribed \cite{stratton1941,harrington2001}.

With solenoidal \(\Jimp\), Eq.~\eqref{eq:ch1.vacuum-wave-source-driven} yields a charge-continuous Maxwell field. Inserting a spurious \(\rho=\rho_0\delta(\mathbf{r})\) without adjusting \(\Jimp\) violates Eq.~\eqref{eq:ch1.charge-continuity-harmonic} and breaks the Amp\`ere--Gauss closure of Eqs.~\eqref{eq:ch1.maxwell-local} \cite{maxwell1873,stratton1941}.

\begin{example}[An electrically small current loop]
\label{ex:ch1.filament-loop}
A filament loop of radius \(b=3\,\mathrm{cm}\) carries phasor current \(I_0=1\,\mathrm{A}\) at \(f=100\,\mathrm{MHz}\). Then \(\lambda_0=c/f=3\,\mathrm{m}\), \(k_0=2\pi/\lambda_0\approx2.09\,\mathrm{rad\,m^{-1}}\), and \(k_0 b=0.063\ll1\): the loop is electrically small. The impressed azimuthal current is solenoidal on its support (\(\divg\Jimp=0\) in the distributional sense for a closed toroidal filament), so Eq.~\eqref{eq:ch1.vacuum-wave-source-driven} drives primarily magnetic-dipole radiation at asymptotic distances (Subsection~\ref{sec:ch01.2.4}) \cite{stratton1941,harrington2001}. At \(k_0 b=0.063\) reactive near fields dominate on shells comparable to \(b\). The value \(k_0 r=6.3\) corresponds to \(r=\lambda_0\) and is used \emph{illustratively} in Figure~\ref{fig:ch1.wp-solenoidal-loop}: it indicates order unity in \(k_0 r\) but does not, without the separate conditions \(k_0 b^2/r\ll1\) and established multipole dominance, replace a full far-zone proof. Fraunhofer-type phase accuracy for this loop would require \(k_0 b^2/r\ll1\), which holds easily at \(r=\lambda_0\) because \(b\ll\lambda_0\) \cite{stratton1941,jackson1999}.
\end{example}

Dispersive media require
frequency-dependent \(\varepsilon(\omega)\) and \(\mu(\omega)\), breaking the
algebraic completion of Eq.~\eqref{eq:ch1.vacuum-constitutive}. Relativistic or
quantum corrections lie outside the present harmonic class. Within the vacuum harmonic class,
Eqs.~\eqref{eq:ch1.maxwell-local}--\eqref{eq:ch1.local-plane-wave-relations}
form the complete local law set every later angular argument inherits
\cite{maxwell1873}.

\subsection{Sources, boundary conditions, and the Maxwell boundary-value problem}
\label{sec:ch01.1.2}

A localized Green-function representation maps impressed sources to outgoing fields in free space, but a device or scatterer is specified by the full Maxwell boundary-value problem: a spatial domain $\Omega$ with boundary $\Gamma=\partial\Omega$, piecewise constitutive laws in subregions $\{\Omega_a\}$, impressed volume and surface sources, transmission conditions on internal interfaces, outer boundary traces, and---for open exterior regions---an outgoing radiation selection \cite{stratton1941,maxwell1873}. Prescribing a tangential field on an aperture without a charge-compatible $\Jimp$ and compatible traces on $\Gamma$ need not define any realizable Maxwell state \cite{harrington2001,chu1948}.

At angular frequency $\omega>0$ under $\exp(+\jj\omega t)$, let $(\Jimp,\rho)$ denote impressed volume sources (A\,m$^{-2}$ and C\,m$^{-3}$), $\mathbf{J}_{s,ab}$ impressed electric surface current on $\Sigma_{ab}=\partial\Omega_a\cap\partial\Omega_b$ (A\,m$^{-1}$), and $\rho_{s,ab}$ free surface charge (C\,m$^{-2}$). In each $\Omega_a$, $\mathbf{D}=\varepsilon_a\E$ and $\mathbf{B}=\mu_a\H$ with piecewise constant $\varepsilon_a,\mu_a$ unless stated otherwise \cite{stratton1941}. The weighted volume test of Subsection~\ref{sec:ch01.0.3} (Eq.~\eqref{eq:ch1.energy-class}) is distinct from global $L^2$ field energy on an unbounded exterior. Eq.~\eqref{eq:ch1.global-charge-compatibility} is the integrated continuity balance obtained from the divergence theorem \emph{when} $\Jimp$ admits a classical outward normal trace on $\Gamma$; it is not a general trace-regularity theorem for distributional sources \cite{stratton1941}.

\paragraph{Interface transmission.}
On each $\Sigma_{ab}$, let $\hat{\mathbf{n}}_{ab}$ be the unit normal directed from $\Omega_a$ to $\Omega_b$. Assume fields are smooth up to the interface except for the allowed jumps and that no impressed magnetic sheet current is present. Faraday's integral in Eq.~\eqref{eq:ch1.maxwell-integral}, applied to a rectangular loop of area $\delta h$ straddling $\Sigma_{ab}$ with sides parallel and normal to the interface, gives $\oint \E\cdot d\boldsymbol{\ell}=-\jj\omega\mu_0\int_S \H\cdot\hat{\mathbf{n}}\,dS$. The right-hand side is $\mathcal{O}(\delta h)$, while the tangential contribution to the line integral is $[\hat{\mathbf{n}}_{ab}\times(\E_b-\E_a)]\cdot\Delta\boldsymbol{\ell}_t$ of order unity as the loop height $\delta\to 0$; hence $\hat{\mathbf{n}}_{ab}\times(\E_b-\E_a)=\mathbf{0}$ \cite{maxwell1873,jackson1999}. Amp\`ere--Maxwell on the same loop yields $\hat{\mathbf{n}}_{ab}\times(\H_b-\H_a)=\mathbf{J}_{s,ab}$ in the limit. Gauss's laws applied to a pillbox of height $\delta$ give $\hat{\mathbf{n}}_{ab}\cdot(\mathbf{D}_b-\mathbf{D}_a)=\rho_{s,ab}$ and, without magnetic monopoles, $\hat{\mathbf{n}}_{ab}\cdot(\mathbf{B}_b-\mathbf{B}_a)=0$ \cite{stratton1941}. Collecting,
\begin{equation}
\label{eq:ch1.interface-conditions}
\begin{aligned}
\hat{\mathbf{n}}_{ab}\times(\E_b-\E_a)&=\mathbf{0},\\
\hat{\mathbf{n}}_{ab}\times(\H_b-\H_a)&=\mathbf{J}_{s,ab},\\
\hat{\mathbf{n}}_{ab}\cdot(\mathbf{D}_b-\mathbf{D}_a)&=\rho_{s,ab},\\
\hat{\mathbf{n}}_{ab}\cdot(\mathbf{B}_b-\mathbf{B}_a)&=0.
\end{aligned}
\end{equation}
Eq.~\eqref{eq:ch1.interface-conditions} is the complete transmission law used below; omitting $\rho_{s,ab}$ or $\mathbf{J}_{s,ab}$ when a coated or fed interface carries them misstates reflected power and source closure \cite{collin2001,jackson1999}.

\paragraph{Outer boundary traces.}
Partition $\Gamma=\Gamma_{\mathrm{PEC}}\cup\Gamma_{\mathrm{PMC}}\cup\Gamma_Z$ with outward unit normal $\hat{\mathbf{n}}$ on $\Gamma$. A PEC boundary is an ideal limit of large tangential $\E$ screening: $\hat{\mathbf{n}}\times\E=\mathbf{0}$. A PMC boundary is the dual magnetic limit: $\hat{\mathbf{n}}\times\H=\mathbf{0}$. Neither is universally interchangeable with a finite conductor model; both exclude penetration physics retained in richer models \cite{stratton1941}. An impedance sheet of surface impedance $Z_s$ ($\Omega$) relates tangential traces by
\begin{equation}
\label{eq:ch1.outer-boundary-classes}
\begin{aligned}
\hat{\mathbf{n}}\times\E&=\mathbf{0}\quad &&\text{on }\Gamma_{\mathrm{PEC}},\\
\hat{\mathbf{n}}\times\H&=\mathbf{0}\quad &&\text{on }\Gamma_{\mathrm{PMC}},\\
\hat{\mathbf{n}}\times\E&=Z_s\,(\hat{\mathbf{n}}\times\H)\times\hat{\mathbf{n}}
\quad &&\text{on }\Gamma_Z,
\end{aligned}
\end{equation}
where $(\hat{\mathbf{n}}\times\H)\times\hat{\mathbf{n}}$ is the tangential magnetic field with its vector direction aligned to the electric tangential trace it couples. On $\Gamma_Z$, cycle-averaged power density normal to the sheet is
\begin{equation}
\label{eq:ch1.impedance-boundary-power}
P_Z=
\frac{1}{2}\Re\!\left\{
\hat{\mathbf{n}}\cdot\left(\E\times\H^*\right)
\right\}
=\frac{1}{2}\Re\!\left\{
Z_s\left|(\hat{\mathbf{n}}\times\H)\times\hat{\mathbf{n}}\right|^2
\right\},
\end{equation}
by Eq.~\eqref{eq:ch1.outer-boundary-classes}; when $\Re(Z_s)>0$, $P_Z\ge 0$ \cite{stratton1941}. All outer classes are collected by the boundary operator
\begin{equation}
\label{eq:ch1.boundary-data-operator}
\mathcal{B}_\Gamma[\E,\H]=\mathbf{g}_\Gamma,
\end{equation}
with prescribed data $\mathbf{g}_\Gamma$ on $\Gamma$ \cite{stratton1941}.
\phantomsection\label{eq:ch1.data-specification-tuple}

\paragraph{Outgoing selection.}
For exterior antenna and scattering problems, fields at large radius $R=|\mathbf{r}|$ must be asymptotically outgoing under $\exp(+\jj\omega t)$:
\begin{equation}
\label{eq:ch1.radiation-trace-condition}
\hat{\mathbf{n}}\times\H
+\eta_0^{-1}(\hat{\mathbf{n}}\times\E)\times\hat{\mathbf{n}}
=\mathbf{o}(R^{-1}),\qquad R\to\infty,
\end{equation}
with $\eta_0=\sqrt{\mu_0/\varepsilon_0}$ and $\hat{\mathbf{n}}=\hat{\mathbf{r}}$ \cite{stratton1941,sommerfeld1949}. Eq.~\eqref{eq:ch1.radiation-trace-condition} excludes the incoming branch proportional to $e^{+\jj k_0 R}/R$ and matches $\E_{\mathrm{rad}}\sim \E_0(\theta,\phi)\,e^{-\jj k_0 R}/R$, $\H_{\mathrm{rad}}=\eta_0^{-1}\hat{\mathbf{r}}\times\E_{\mathrm{rad}}$. It is a selection rule on the exterior solution, not a proof of existence, uniqueness, finite global energy, or full trace regularity for later torque integrals \cite{harrington2001}. Absorbing layers at finite radius approximate the same physics on a truncated domain \cite{hansen1988}.

\paragraph{Source classes and global balance.}
Bounded excitation occupies $\Omega_s\subseteq\Omega$ and must satisfy Eq.~\eqref{eq:ch1.charge-continuity-harmonic}. The regular class is
\begin{equation}
\label{eq:ch1.source-charge-continuous}
\Jimp\in L^2(\Omega_s)^3,\qquad
\operatorname{supp}(\Jimp)\subseteq\overline{\Omega_s},\qquad
\divg\Jimp+\jj\omega\rho=0
\ \text{in }\mathcal{D}'(\Omega).
\end{equation}
Square-integrable $\Jimp$ with distributional continuity does not automatically imply finite exterior field energy, classical normal traces of $\Jimp$ on $\Gamma$, or finite quadratic surface fluxes; those require separate hypotheses tied to the domain and outgoing selection \cite{stratton1941}. Closed filaments and segmented dipoles in examples are \emph{distributional} limits with $\divg\Jimp+\jj\omega\rho=0$ in $\mathcal{D}'(\Omega)$, not members of $L^2$ unless regularized by finite width \cite{stratton1941,jackson1999}. Integrating Eq.~\eqref{eq:ch1.charge-continuity-harmonic} over $\Omega$ and applying the divergence theorem yields, \emph{when $\Jimp$ admits a classical outward normal trace on $\Gamma$},
\begin{equation}
\label{eq:ch1.global-charge-compatibility}
\int_{\Omega}\!\left(\divg\Jimp+\jj\omega\rho\right)\,dV
=\int_{\Gamma}\!\Jimp\cdot\hat{\mathbf{n}}\,dS+\jj\omega\int_{\Omega}\rho\,dV=0.
\end{equation}
If $\Jimp$ is supported strictly inside $\Omega$, the surface term vanishes; otherwise impressed current crossing $\Gamma$ must be specified. A general distributional trace theorem for $\Jimp|_\Gamma$ is not established in this manuscript and remains an open extension \cite{maxwell1873,stratton1941}.

\phantomsection\label{prop:ch1.source-boundary-compatibility}
\begin{example}[Reflection at a lossless dielectric interface]
\label{ex:ch1.dielectric-interface}
Let $\Omega_1$ be vacuum ($z>0$) and $\Omega_2$ a semi-infinite dielectric with $\varepsilon_r=2.25$, $\mu_r=1$ occupying $z<0$, with interface at $z=0$ and $\hat{\mathbf{n}}_{12}=-\hat{\mathbf{z}}$ directed from $\Omega_1$ into $\Omega_2$, consistent with the definition of $\hat{\mathbf{n}}_{ab}$ in Subsection~\ref{sec:ch01.1.2}. A normally incident plane wave in $\Omega_1$ travels toward the interface from $z>0$ with $\E=E_0 e^{+\jj k_0 z}\hat{\mathbf{x}}$ under $\exp(+\jj\omega t)$. Eq.~\eqref{eq:ch1.interface-conditions} with $\mathbf{J}_{s,12}=\mathbf{0}$ gives continuous tangential $\E$ and $\H$. With intrinsic impedances $\eta_1=\eta_0$ and $\eta_2=\eta_0/\sqrt{\varepsilon_r}$, the reflection coefficient is $\Gamma=(\eta_2-\eta_1)/(\eta_2+\eta_1)=(1-\sqrt{\varepsilon_r})/(1+\sqrt{\varepsilon_r})=-0.2$, so $R=|\Gamma|^2=0.04$ and $T=1-R=0.96$ for a lossless interface; the transmitted field in $\Omega_2$ carries wavenumber $k_2=k_0\sqrt{\varepsilon_r}$ and amplitude $|E_t|=(1+\Gamma)|E_0|=0.8|E_0|$ \cite{stratton1941,jackson1999}.
\end{example}

\begin{figure}[t]
\centering
\VolIPhysFig{\csname VolIPhysWaveguideTE10\endcsname}
\caption{Rectangular waveguide cross section $a\times b$ with PEC sidewalls, TE$_{10}$ transverse $\E$ profile, and charge-compatible impressed-current feed at the aperture.}
\label{fig:ch1.wp-te10-waveguide}
\end{figure}
\begin{example}[TE$_{10}$ mode in a closed rectangular guide]
\label{ex:ch1.te10-waveguide}
A hollow PEC guide with cross section $a\times b$ supports $\E=E_0\sin(\pi x/a)\,e^{-\jj\beta z}\hat{\mathbf{y}}$ with cutoff $\omega_c=\pi c/a$ and $\beta=\sqrt{k_0^2-(\pi/a)^2}$ when $k_0>\pi/a$. Tangential $\E$ vanishes on $x=0,a$ and $y=0,b$, implementing Eq.~\eqref{eq:ch1.outer-boundary-classes} on $\Gamma_{\mathrm{PEC}}$. This is a \emph{closed} boundary-value problem: Eq.~\eqref{eq:ch1.radiation-trace-condition} at infinity is not part of the closed guide eigenproblem, though an open-ended or fed guide coupled to free space requires a different exterior completion \cite{stratton1941,collin2001}. At $f=10\,\mathrm{GHz}$ with $a=2.3\,\mathrm{cm}$, $\omega_c\approx 6.52\,\mathrm{GHz}$ and $\beta/k_0\approx 0.759$; below cutoff the same geometry is evanescent and storage-dominated \cite{harrington2001}.
\end{example}

Compatible $(\Jimp,\rho)$, constitutive data, Eq.~\eqref{eq:ch1.interface-conditions}, outer data Eq.~\eqref{eq:ch1.boundary-data-operator}, and outgoing selection Eq.~\eqref{eq:ch1.radiation-trace-condition} where applicable define the class of Maxwell fields under consideration. Existence of a solution in that class and uniqueness under additional hypotheses are treated in Subsection~\ref{sec:ch01.1.3}; flux, power, and angular-momentum integrals presuppose this closure \cite{maxwell1873,stratton1941}. Piecewise smooth interfaces and linear local constitutive laws are assumed; rough or nonlinear boundaries require enriched models before the same conservation interpretations apply \cite{stratton1941}.

\subsection{Boundary-Value Problems and Uniqueness}
\label{sec:ch01.1.3}

Power and torque integrals name a physical field only if the source, boundary, and radiation data select one Maxwell solution \cite{stratton1941,sommerfeld1949}. Let $(\E_1,\H_1)$ and $(\E_2,\H_2)$ be two solutions with identical impressed sources, interface data, and outer traces. Their difference $(\Delta\E,\Delta\H)$ is source-free ($\Delta\Jimp=\mathbf{0}$) and satisfies the homogeneous boundary operator $\mathcal{B}_\Gamma[\Delta\E,\Delta\H]=\mathbf{0}$ \cite{stratton1941}. Under $\exp(+\jj\omega t)$, the harmonic difference fields obey the same curl system with domain constitutive $\varepsilon,\mu$. The vector identity $\divg(\Delta\E\times\Delta\H^*)=\Delta\H^*\cdot\curl\Delta\E-\Delta\E\cdot\curl\Delta\H^*$ together with $\curl\Delta\E=-\jj\omega\mu\Delta\H$ and $\curl\Delta\H=\jj\omega\varepsilon\Delta\E$ (and $\curl\Delta\H^*=-\jj\omega\varepsilon^*\Delta\E^*$ for complex $\varepsilon$) gives the local relation
\begin{equation}
\label{eq:ch1.uniqueness-local-div}
\divg\mathbf{S}_\Delta=-2\jj\omega\left(u_{\Delta H}-u_{\Delta E}\right)
\end{equation}
for $\mathbf{S}_\Delta=\tfrac12\Delta\E\times\Delta\H^*$, with $u_{\Delta E}=\tfrac14\varepsilon|\Delta\E|^{2}$ and $u_{\Delta H}=\tfrac14\mu|\Delta\H|^{2}$ \cite{stratton1941}. The factor $2$ matches Eq.~\eqref{eq:ch1.uniqueness-energy-identity}: a half-sized local divergence would not reproduce the integrated identity. Integrating over $\Omega$ and applying the divergence theorem yields
\begin{equation}
\label{eq:ch1.uniqueness-energy-identity}
\int_{\Gamma}\hat{\mathbf{n}}\cdot\mathbf{S}_\Delta\,dS
=-2\jj\omega\!\int_{\Omega}\!\left(u_{\Delta H}-u_{\Delta E}\right)dV
\end{equation}
when $\Delta\E,\Delta\H$ are smooth enough for the classical divergence theorem on $\Omega$ \cite{stratton1941}. Eq.~\eqref{eq:ch1.uniqueness-energy-identity} does not by itself prove existence or global regularity; it constrains the difference field.

For real $\varepsilon$ and $\mu$ on $\Omega$, Eq.~\eqref{eq:ch1.uniqueness-energy-identity} has a purely imaginary right-hand side, so its \emph{real part} gives $\Re\{\int_{\Gamma}\hat{\mathbf{n}}\cdot\mathbf{S}_\Delta\,dS\}=0$ for any difference field: lossless interiors do not by themselves exclude nontrivial source-free modes at cavity eigenfrequencies \cite{stratton1941}. For complex scalar isotropic $\varepsilon,\mu$ in the source-free difference problem, the same vector identity with $\curl\Delta\H^*=-\jj\omega\varepsilon^*\Delta\E^*$ gives
\begin{equation}
\label{eq:ch1.uniqueness-local-div-complex}
\nabla\cdot\mathbf{S}_\Delta
=\frac{\jj\omega}{2}\left(\varepsilon^*|\Delta\E|^{2}-\mu|\Delta\H|^{2}\right),
\end{equation}
which reduces to Eq.~\eqref{eq:ch1.uniqueness-local-div} when $\varepsilon$ and $\mu$ are real \cite{stratton1941,harrington2001}. When dissipation is included through $\varepsilon=\varepsilon'-\jj\varepsilon''$ and/or $\mu=\mu'-\jj\mu''$ with $\varepsilon'',\mu''\ge0$ (passive loss under $\exp(+\jj\omega t)$),
\[
\Re\{\nabla\cdot\mathbf{S}_\Delta\}
=-\frac{\omega}{2}\left(\varepsilon''|\Delta\E|^{2}+\mu''|\Delta\H|^{2}\right)
\]
for isotropic local $\varepsilon,\mu$ in the interior (volume divergence, not a surface flux density) \cite{stratton1941,harrington2001}. Integration over $\Omega$ and the divergence theorem yield
\[
\Re\!\left\{\int_{\Gamma}\hat{\mathbf{n}}\cdot\mathbf{S}_\Delta\,dS\right\}
=-\frac{\omega}{2}\int_{\Omega}\left(\varepsilon''|\Delta\E|^{2}+\mu''|\Delta\H|^{2}\right)dV
\le 0,
\]
with equality forcing $\Delta\E=\mathbf{0}$ on subsets where $\varepsilon''>0$ and $\Delta\H=\mathbf{0}$ where $\mu''>0$, subject to the integrated balance and boundary data \cite{stratton1941,harrington2001}. Vanishing of \emph{both} difference fields throughout a connected lossy region requires both loss mechanisms (or a separate unique-continuation argument) in addition to the integral identity \cite{stratton1941,harrington2001}. On a bounded lossy domain with homogeneous PEC or PMC data on the entire $\Gamma$, the real surface flux of $\mathbf{S}_\Delta$ vanishes; when $\varepsilon''>0$ and $\mu''>0$ throughout $\Omega$, the volume integral then forces $\Delta\E=\Delta\H=\mathbf{0}$ \cite{stratton1941}. Passive impedance sheets on $\Gamma_Z$ contribute a boundary term through Eq.~\eqref{eq:ch1.outer-boundary-classes} and the normal complex-Poynting density Eq.~\eqref{eq:ch1.impedance-boundary-power}; when $\Re(Z_s)>0$, that contribution is sign-definite and enters the integrated identity rather than being assumed zero \cite{stratton1941,harrington2001}. A \emph{lossless} closed cavity at an eigenfrequency may still support a nontrivial homogeneous mode---the usual exception \cite{stratton1941,sommerfeld1949}. On an unbounded exterior radiating problem, Eq.~\eqref{eq:ch1.uniqueness-energy-identity} must be combined with the outgoing selection Eq.~\eqref{eq:ch1.radiation-trace-condition} so the difference field carries no net incoming power at infinity \cite{sommerfeld1949,harrington2001}. For an \emph{exterior} problem, let $\Omega_{\mathrm{ext}}$ be the complement of a bounded obstacle $\mathcal{O}$ with piecewise smooth boundary $\Gamma=\partial\mathcal{O}$, filled with homogeneous lossless vacuum ($\varepsilon_0,\mu_0$). Take $\hat{\mathbf{n}}_{\mathrm{ext}}$ as the unit normal \emph{outward from $\Omega_{\mathrm{ext}}$} on $\Gamma$ (hence $\hat{\mathbf{n}}_{\mathrm{ext}}=-\hat{\mathbf{n}}_{\mathrm{body}}$). Impedance data on the entire $\Gamma$ are the third line of Eq.~\eqref{eq:ch1.outer-boundary-classes} with $Z_s$ constant scalar and $\Re(Z_s)>0$. For a homogeneous difference field $(\Delta\E,\Delta\H)$ with no impressed sources in $\Omega_{\mathrm{ext}}$, Eq.~\eqref{eq:ch1.impedance-boundary-power} gives $\Re\{\hat{\mathbf{n}}_{\mathrm{ext}}\cdot\mathbf{S}_\Delta\}=\tfrac12\Re(Z_s)|\Delta\H_t|^{2}\ge0$ on $\Gamma$, i.e.\ nonnegative power leaving the exterior field domain across the sheet. Integrating the lossless-exterior form of Eq.~\eqref{eq:ch1.uniqueness-local-div} over a ball of radius $R$ truncated outside $\mathcal{O}$ and letting $R\to\infty$, the real part of Eq.~\eqref{eq:ch1.uniqueness-energy-identity} gives $\Re\{\int_{\Gamma\cup S_R}\hat{\mathbf{n}}_{\mathrm{ext}}\cdot\mathbf{S}_\Delta\,dS\}=0$ for the difference field. This does \emph{not} mean that a physical outgoing radiator has zero real power through $S_R$---a plane wave carries $\Re\{\mathbf{S}\cdot\hat{\mathbf{r}}\}>0$---but for $(\Delta\E,\Delta\H)$ any strictly positive outgoing far-sphere contribution would contradict that balance together with the passive impedance term on $\Gamma$. The outgoing selection Eq.~\eqref{eq:ch1.radiation-trace-condition} then excludes a nontrivial incoming difference branch at infinity, and uniqueness of the exterior Maxwell state follows from the cited impedance-boundary theory under the stated scalar passive class \cite{colton1981impedance,sommerfeld1949,stratton1941}. Nonlocal, anisotropic, or approximate higher-order impedance operators are \emph{not} covered by this scalar local passive class \cite{stratton1941}. Passivity, loss, boundary class, and radiation closure remain distinct hypotheses; uniqueness is stated only for the class actually proved \cite{stratton1941}.
\phantomsection\label{prop:ch1.uniqueness-passive-closure}
\phantomsection\label{eq:ch1.uniqueness-condition-index}

\subsection{Energy, linear momentum, and Maxwell stress tensor}
\label{sec:ch01.1.4}

Energy, linear momentum, and angular momentum are not separate postulates; they follow from Maxwell's curl--divergence structure on a declared control volume \cite{poynting1884,maxwell1873}. The causal chain developed here and completed in Subsection~\ref{sec:ch01.1.5} is: Eqs.~\eqref{eq:ch1.maxwell-local} $\Rightarrow$ instantaneous Poynting identity Eq.~\eqref{eq:ch1.poynting-instantaneous} $\Rightarrow$ harmonic complex balance Eq.~\eqref{eq:ch1.poynting-complex} $\Rightarrow$ Lorentz force and symmetric Maxwell stress Eqs.~\eqref{eq:ch1.momentum-balance-instantaneous} and~\eqref{eq:ch1.maxwell-stress} $\Rightarrow$ stress moment and outward angular flux Eq.~\eqref{eq:ch1.outward-am-flux}. Let $V$ be a control volume with outward unit normal $\hat{\mathbf{n}}$ on $\partial V$; let $V_m$ enclose the matter whose force or torque is evaluated; let $\Gamma_R$ be a predominantly radiative enclosing shell when radiated power or flux is read \cite{stratton1941}. Finite surface integrals require the regularity hypotheses of Subsection~\ref{sec:ch01.0.3} on the same volume---Eq.~\eqref{eq:ch1.energy-class} alone does not supply every trace needed later.

Begin with the instantaneous vacuum fields. The identity
$\divg(\E\times\H)=\H\cdot(\curl\E)-\E\cdot(\curl\H)$ together with Maxwell's equations
gives
\begin{equation}
\label{eq:ch1.poynting-instantaneous}
\divg(\E\times\H)
=-\Jimp\cdot\E
-\frac{\partial}{\partial t}\!\left(\frac{\varepsilon_0}{2}|\E|^2+\frac{\mu_0}{2}|\H|^2\right),
\end{equation}
where $\Jimp$ is impressed current density \cite{stratton1941}. Define instantaneous
Poynting vector $\mathbf{s}=\E\times\H$ and energy density
$u=\tfrac12(\varepsilon_0|\E|^2+\mu_0|\H|^2)$; then
$\divg\mathbf{s}+\partial u/\partial t+\Jimp\cdot\E=0$. The term $\Jimp\cdot\E$ is
power density at which the field performs work on impressed current per unit volume; $\divg\mathbf{s}$ is net energy
outflow from an infinitesimal volume. Integrating over volume $V$ and applying the
divergence theorem yields the power balance through $\partial V$ \cite{poynting1884}.

Instantaneous transport uses real fields: $\mathbf{s}=\E\times\H$, $\mathbf{g}_{\mathrm{em}}=\varepsilon_0\,\E\times\B$, and $\Tinst$ in Eqs.~\eqref{eq:ch1.maxwell-stress-components}--\eqref{eq:ch1.momentum-balance-instantaneous} \cite{maxwell1873,poynting1884}. Harmonic phasors use $\mathbf{S}_\mathrm{c}=\tfrac12\E\times\H^*$, $\langle\mathbf{S}\rangle=\Re\mathbf{S}_\mathrm{c}$, and cycle-averaged stress $\Tavg$ of Eq.~\eqref{eq:ch1.maxwell-stress} with $\Mavg=\mathbf{r}\times\Tavg$ \cite{stratton1941}. The two descriptions are linked by $\exp(+\jj\omega t)$ but must not be interchanged when reading reactive storage or radiative flux.

For harmonic fields with $\exp(+\jj\omega t)$, time averages reduce bilinear products to
phasor forms. Define the complex Poynting vector
$\mathbf{S}_\mathrm{c}=\tfrac12\E\times\H^*$, time-averaged real power flux
$\langle\mathbf{S}\rangle=\Re(\mathbf{S}_\mathrm{c})$, and harmonic storage densities
$u_E=\tfrac14\varepsilon_0\E\cdot\E^*$, $u_H=\tfrac14\mu_0\H\cdot\H^*$. Applying
$\divg(\E\times\H^*)=\H^*\!\cdot(\curl\E)-\E\!\cdot(\curl\H^*)$ to the phasor Maxwell
equations gives Eq.~\eqref{eq:ch1.poynting-complex} below \cite{stratton1941,harrington2001}.
With $\curl\E=-\jj\omega\mu_0\H$ and
$\curl\H=\Jimp+\jj\omega\varepsilon_0\E$, and $\curl\H^*=\Jimp^*-\jj\omega\varepsilon_0\E^*$,
\begin{equation}
\label{eq:ch1.poynting-intermediate}
\divg(\E\times\H^*)
=-\E\cdot\Jimp^*
+\jj\omega\bigl(\varepsilon_0\E\cdot\E^*-\mu_0\H^*\!\cdot\H\bigr).
\end{equation}
Under $\exp(+\jj\omega t)$, the complex Poynting vector is
\(\mathbf{S}_\mathrm{c}=\tfrac{1}{2}\E\times\H^*\), where \((\cdot)^*\) denotes complex
conjugation; time-averaged real power flux is
\(\langle\mathbf{S}\rangle=\Re(\mathbf{S}_\mathrm{c})\). The harmonic stored-energy
densities are
\(u_E=\tfrac{1}{4}\varepsilon_0\E\cdot\E^*\) and
\(u_H=\tfrac{1}{4}\mu_0\H\cdot\H^*\), so
$\varepsilon_0\E\cdot\E^*-\mu_0\H^*\!\cdot\H=4(u_E-u_H)$ \cite{stratton1941}.
Dividing Eq.~\eqref{eq:ch1.poynting-intermediate} by two then gives
\begin{equation}
\label{eq:ch1.poynting-average}\divg\mathbf{S}_\mathrm{c}
=-\frac{1}{2}\E\cdot\Jimp^*
-2\jj\omega\left(u_H-u_E\right).
\end{equation}
\addtocounter{equation}{-1}\refstepcounter{equation}\label{eq:ch1.dual-invariant-poynting}%
\addtocounter{equation}{-1}\refstepcounter{equation}\label{eq:ch1.poynting-complex}%
Physically, Eq.~\eqref{eq:ch1.poynting-complex} states that net complex power leaving a differential volume equals electromagnetic work on impressed current plus harmonic exchange between electric and magnetic storage \cite{poynting1884}. The real part $\divg\Re(\mathbf{S}_\mathrm{c})$ balances impressed work; the imaginary part $\divg\Im(\mathbf{S}_\mathrm{c})$ tracks reactive exchange in the complex-power balance and is not, without further definition, a unique local stored-energy density \cite{stratton1941}. Momentum in material media involves additional stress--energy forms (Abraham, Minkowski); vacuum-side traction and angular flux use the Maxwell stress of Eq.~\eqref{eq:ch1.maxwell-stress} \cite{abraham1910,minkowski1908,stratton1941}.

\begin{equation}
\label{eq:ch1.poynting-real-imag-split}
\begin{aligned}
\divg\Re(\mathbf{S}_\mathrm{c})&=-\frac12\Re\!\left(\E\cdot\Jimp^*\right),\\
\divg\Im(\mathbf{S}_\mathrm{c})&=-\frac12\Im\!\left(\E\cdot\Jimp^*\right)-2\omega(u_H-u_E).
\end{aligned}
\end{equation}
\begin{equation}
\label{eq:ch1.dual-invariant-energy}\int_{\partial V}\mathbf{S}_\mathrm{c}\cdot\hat{\mathbf{n}}\,dS
=-\frac12\int_V\E\cdot\Jimp^*\,dV
-2\jj\omega\int_V(u_H-u_E)\,dV.
\end{equation}
\addtocounter{equation}{-1}\refstepcounter{equation}\label{eq:ch1.power-balance-integral}%
Here $V\subset\Omega$ is a control volume with outward normal $\hat{\mathbf{n}}$ on $\partial V$ \cite{stratton1941}. The surface integral is net complex power through $\partial V$; the volume terms are impressed work and reactive storage imbalance inside $V$ \cite{poynting1884}. Integrating the local balance over any finite region therefore relates power crossing the boundary to source work and stored reactive energy within the volume---the finite-volume form of Poynting's theorem \cite{stratton1941}.

Mechanical force on matter follows from the Lorentz law,
\begin{equation}
\label{eq:ch1.lorentz-force}
\mathbf{f}_{\mathrm{matter}}=\rho_{\mathrm{phys}}\E+\mathbf{J}_{\mathrm{phys}}\times\B,
\end{equation}
where $\rho_{\mathrm{phys}}$ and $\mathbf{J}_{\mathrm{phys}}$ are the charge and current densities of the carriers that experience the force (N\,m$^{-3}$) \cite{maxwell1873,stratton1941}.
In Eqs.~\eqref{eq:ch1.maxwell-local}, $\rho$ is the charge density in Gauss's law and $\Jimp$ is the impressed current in Amp\`ere's law---the mathematical source pair $(\rho,\Jimp)$ that closes Maxwell's equations \cite{maxwell1873,stratton1941}. Route~A identifies $\rho_{\mathrm{phys}}=\rho$ and $\mathbf{J}_{\mathrm{phys}}=\Jimp$ on the support where impressed sources represent the actual carriers, so $\mathbf{f}_{\mathrm{matter}}$ equals the source-exchange density $\mathbf{f}_{\mathrm{src}}$ below. Route~B treats $(\rho,\Jimp)$ as impressed or equivalent excitation; then $\mathbf{f}_{\mathrm{matter}}$ applies on material supports, net force and torque on a body follow from $\oint_{\partial V_b}\Tinst\cdot\hat{\mathbf{n}}\,dS$ on a closed surface in source-free vacuum enclosing that body, and the Poynting term $-\tfrac12\E\cdot\Jimp^*$ is impressed-source exchange, not automatic mechanical work on an arbitrary load \cite{poynting1884,stratton1941,harrington2001}.
The symmetric vacuum stress tensor Eqs.~\eqref{eq:ch1.maxwell-stress-components}--\eqref{eq:ch1.maxwell-stress} and the local identity below are derived for regions where $(\E,\B)$ obey vacuum constitutive relations $\varepsilon_0,\mu_0$ and the displayed curl--divergence structure; they are not asserted pointwise inside arbitrary macroscopic media without an applicable material stress model \cite{stratton1941}. Using Maxwell's equations in such a vacuum/source region to eliminate the same $(\rho,\Jimp)$ as in Eq.~\eqref{eq:ch1.maxwell-local}, in favour of $(\E,\B)$,
one forms the electromagnetic momentum density
\begin{equation}
\label{eq:ch1.momentum-density-instantaneous}
\mathbf{g}_{\mathrm{em}}=\varepsilon_0\,\E\times\B
=\mathbf{s}/c^2,
\end{equation}
where $\mathbf{s}=\E\times\H$ is the instantaneous Poynting vector \cite{maxwell1873,poynting1884}.
Radiation pressure on an absorber is momentum delivered through $\Tinst$ or, in harmonic form, through $\Tavg$ \cite{maxwell1873,stratton1941}. The instantaneous symmetric Maxwell stress in Cartesian components is
\begin{equation}
\label{eq:ch1.maxwell-stress-components}
T^{(i)}_{ij}=\varepsilon_0\!\left(E_iE_j-\tfrac12\delta_{ij}E^2\right)
+\frac{1}{\mu_0}\!\left(B_iB_j-\tfrac12\delta_{ij}B^2\right),
\end{equation}
with $E^2=\mathbf{E}\cdot\mathbf{E}$ and $B^2=\mathbf{B}\cdot\mathbf{B}$ \cite{maxwell1873,stratton1941}.
Index contraction with Gauss's and Amp\`ere's laws gives the field momentum balance
\begin{equation}
\label{eq:ch1.momentum-balance-instantaneous}
\frac{\partial g_{i}}{\partial t}+f_{\mathrm{src},i}=\partial_j T^{(i)}_{ij},
\end{equation}
equivalently $\partial_t\mathbf{g}_{\mathrm{em}}+\mathbf{f}_{\mathrm{src}}=\divg\Tinst$, where
\begin{equation}
\label{eq:ch1.maxwell-source-force}
\mathbf{f}_{\mathrm{src}}=\rho\E+\Jimp\times\B
\end{equation}
uses the same $\rho$ and $\Jimp$ as Eqs.~\eqref{eq:ch1.maxwell-local} \cite{maxwell1873,jackson1999}. With real fields $\E(\mathbf{r},t)=\Re\{\E(\mathbf{r})e^{+\jj\omega t}\}$ and similarly for $\H$, $\B$, $\Jimp$, and $\rho$, cycle averaging gives
$\langle\mathbf{f}_{\mathrm{src}}\rangle=\tfrac12\Re\{\rho\E^*+\Jimp\times\B^*\}$,
distinct from the instantaneous product $\mathbf{f}_{\mathrm{src}}=\rho\E+\Jimp\times\B$ \cite{stratton1941}. Eq.~\eqref{eq:ch1.momentum-balance-instantaneous} is field--source bookkeeping; $\mathbf{f}_{\mathrm{matter}}$ equals $\mathbf{f}_{\mathrm{src}}$ only under Route~A on the relevant support.
Throughout the stress, Poynting, and torque balances of
Subsections~\ref{sec:ch01.1.4}--\ref{sec:ch01.1.5}, $V$ denotes a control volume
enclosing the matter whose mechanical force or torque is being evaluated, and
$\hat{\mathbf{n}}$ is the unit normal outward from $V$ \cite{stratton1941}.
Integrating over that $V$ with outward normal $\hat{\mathbf{n}}$,
\begin{equation}
\label{eq:ch1.momentum-integral-balance}
\mathbf{F}_{\mathrm{matter}}
=\int_V \mathbf{f}_{\mathrm{matter}}\,dV
=\int_{\partial V}\Tinst\cdot\hat{\mathbf{n}}\,dS
-\frac{d}{dt}\int_V \mathbf{g}_{\mathrm{em}}\,dV
\end{equation}
when $V$ lies in a vacuum/source region where Eq.~\eqref{eq:ch1.momentum-balance-instantaneous} holds and Route~A applies throughout $V$, so $\mathbf{f}_{\mathrm{matter}}=\mathbf{f}_{\mathrm{src}}$ there \cite{stratton1941}. Under Route~B, $\int_V \mathbf{f}_{\mathrm{src}}\,dV$ is source-exchange bookkeeping on the chosen volume, while mechanical force on a body is $\mathbf{F}=\oint_{\partial V_b}\Tinst\cdot\hat{\mathbf{n}}\,dS$ (minus storage) for a closed $\partial V_b$ in source-free vacuum enclosing that body, with traction evaluated on the vacuum field side \cite{stratton1941,harrington2001}. Traction must not be identified with momentum density $\mathbf{g}_{\mathrm{em}}$ or with Poynting flux $\mathbf{S}_\mathrm{c}$.

\begin{figure}[t]
\centering
\VolIPhysFig{\VolIPhysStressTraction}
\caption{Differential surface element with outward normal $\hat{\mathbf{n}}$: cycle-averaged traction $\mathbf{t}=\Tavg\cdot\hat{\mathbf{n}}$ on a material-enclosing control volume $V_m$ \cite{maxwell1873,stratton1941}.}
\label{fig:ch1.stress-traction-element}
\end{figure}

Figure~\ref{fig:ch1.stress-traction-element} shows cycle-averaged traction $\mathbf{t}=\Tavg\cdot\hat{\mathbf{n}}$ on a material-enclosing volume $V_m$.

For harmonic fields with $\exp(+\jj\omega t)$, time averaging yields the cycle-averaged dyadic
\begin{equation}
\label{eq:ch1.maxwell-stress}
\Tavg
=\frac{1}{2}\Re\!\left[
\varepsilon_0\,\E\E^*
+\mu_0\,\H\H^*
-\frac{1}{2}\mathbf{I}
\left(\varepsilon_0\,\E\cdot\E^*+\mu_0\,\H\cdot\H^*\right)
\right],
\end{equation}
where \(\mathbf{I}\) is the identity dyad and juxtaposition such as \(\E\E^*\)
is dyadic product. Operand by operand, the first two dyads are directional field
intensity contributions, while the isotropic subtraction term
\(\tfrac{1}{2}\mathbf{I}(\cdots)\) removes trace pressure so that directional
tractions are represented correctly. The real-part operator keeps physically
observable time-averaged transfer \cite{stratton1941}.
Together, $\mathbf{S}_\mathrm{c}$ and $\langle\mathbf{S}\rangle$ give energetic flux densities and $\Tavg$ gives cycle-averaged mechanical traction; Subsection~\ref{sec:ch01.1.5} uses $\Mavg=\mathbf{r}\times\Tavg$ and stress-moment surface density $\Mavg\cdot\hat{\mathbf{n}}$ \cite{poynting1884}.

On a material-enclosing volume $V_m$, traction $\mathbf{t}=\Tavg\cdot\hat{\mathbf{n}}$ is cycle-averaged force per unit area on matter, and $\mathbf{F}=\int_{\partial V_m}\Tavg\cdot\hat{\mathbf{n}}\,dS$ \cite{stratton1941}.
For a plane wave in $+\hat{\mathbf{z}}$ incident on a perfectly absorbing sheet at
$z=0$ with vacuum control volume $z<0$, the outward unit normal on the interface is
$\hat{\mathbf{n}}=+\hat{\mathbf{z}}$. Radiation pressure on the absorber is
$p_{\mathrm{abs}}=|\langle\mathbf{S}\rangle|/c$; a perfect reflector at normal incidence receives
$p_{\mathrm{refl}}=2|\langle\mathbf{S}\rangle|/c$ \cite{stratton1941,maxwell1873}.

\begin{example}[Radiation pressure on an absorbing surface]
\label{ex:ch1.radiation-pressure}
With $\langle\mathbf{S}\rangle=(|E_0|^2/2\eta_0)\hat{\mathbf{z}}$ from Eq.~\eqref{eq:ch1.plane-wave-transverse}, $p_{\mathrm{abs}}=|E_0|^2/(2\eta_0 c)$. If $|E_0|=100\,\mathrm{V/m}$, then $\langle S_z\rangle=13.2\,\mathrm{W/m^2}$ and $p_{\mathrm{abs}}\approx4.4\times10^{-8}\,\mathrm{N/m^2}$; a perfect reflector doubles the pressure \cite{stratton1941}. The same result follows from the momentum balance: in vacuum $\langle\mathbf{g}_{\mathrm{em}}\rangle=\Re(\mathbf{S}_\mathrm{c})/c^2$ (Eq.~\eqref{eq:ch1.rotation-generator-split}), so the cycle-averaged momentum flux through the interface is $\langle\mathbf{g}_{\mathrm{em}}\rangle\cdot\hat{\mathbf{n}}$ times $c$ when the wave is absorbed; specular reversal reverses $\langle\mathbf{g}_{\mathrm{em}}\rangle$ and doubles the magnitude of the momentum delivered to the reflector \cite{stratton1941}.
\end{example}

\begin{figure}[t]
\centering
\VolIPhysFig{\VolIPhysPlaneWaveRadiationPressure}
\caption{Uniform plane wave in $+\hat{\mathbf{z}}$ on field control volume $V_f$ ($z<0$): outward interface normal $\hat{\mathbf{n}}=+\hat{\mathbf{z}}$ at $z=0$; perfect absorber receives radiation pressure $|\langle\mathbf{S}\rangle|/c$, perfect reflector at normal incidence $2|\langle\mathbf{S}\rangle|/c$; transverse $\E$, $\H$, and $\langle\mathbf{S}\rangle$ remain right-handed \cite{stratton1941,maxwell1873}.}
\label{fig:ch1.plane-wave-handedness}
\end{figure}
Integrating the momentum flux over a control volume yields the harmonic steady-state balance
\begin{equation}
\label{eq:ch1.linear-momentum-balance}
\langle\mathbf{f}_{\mathrm{src}}\rangle=\divg\Tavg,
\end{equation}
where \(\langle\mathbf{f}_{\mathrm{src}}\rangle=\tfrac12\Re\{\rho\E^*+\Jimp\times\B^*\}\) is the cycle-averaged source-exchange density (distinct from the instantaneous \(\mathbf{f}_{\mathrm{src}}\) of Eq.~\eqref{eq:ch1.maxwell-source-force}) and \(\partial_t\langle\mathbf{g}_{\mathrm{em}}\rangle=\mathbf{0}\) on the harmonic branch \cite{stratton1941}. Eq.~\eqref{eq:ch1.linear-momentum-balance} is the steady phasor limit of Eq.~\eqref{eq:ch1.momentum-balance-instantaneous}; mechanical force on carriers remains \(\mathbf{f}_{\mathrm{matter}}\) of Eq.~\eqref{eq:ch1.lorentz-force}, equal to \(\mathbf{f}_{\mathrm{src}}\) only under Route~A.

\phantomsection\label{eq:ch1.energy-momentum-directional-ratio}
The plane wave provides the simplest standard example in which energy flow and momentum flow become visibly parallel \cite{poynting1884,stratton1941}. For an outgoing wave with $\mathbf{k}=k\hat{\mathbf{z}}$ and transverse electric field $\E=E_0 e^{-\jj k z}\hat{\mathbf{x}}$,
\begin{equation}
\label{eq:ch1.plane-wave-transverse}
\H=\frac{1}{\eta_0}\,\hat{\mathbf{k}}\times\E,
\qquad
\langle\mathbf{S}\rangle=\frac{|\E|^{2}}{2\eta_0}\,\hat{\mathbf{k}},
\end{equation}
where $\eta_0=\sqrt{\mu_0/\varepsilon_0}$ is the vacuum intrinsic impedance \cite{stratton1941}. The time-averaged Poynting vector and normal Maxwell traction then point along $\hat{\mathbf{k}}$, with radiation pressure $|\mathbf{t}|=|\langle\mathbf{S}\rangle|/c$ on the same Maxwell branch. In structured near fields energetic and mechanical transport need not be co-directed---a common situation in reactive storage zones \cite{stratton1941}.

\phantomsection\label{prop:ch1.energy-momentum-consistency}
For a Maxwell solution branch fixed by impressed sources, interface data, boundary completion, and vacuum constitutive law on the integration region, Eqs.~\eqref{eq:ch1.poynting-complex} and~\eqref{eq:ch1.maxwell-stress} are quadratic functionals of the same \((\E,\H)\). Poynting and stress surface integrals on \(\partial V\) are defined only when the corresponding integrands are integrable on that surface---typically when \((\E,\H)\) has sufficient regularity on \(V\) (Subsection~\ref{sec:ch01.0.3}); Eq.~\eqref{eq:ch1.energy-class} on \(V\) is helpful but not sufficient by itself, and finite \(P_{\mathrm{rad}}\) on a shell \(S_R\) alone does not replace those hypotheses \cite{stratton1941}. Both balances then describe transport generated by one field state, which is the common basis for the angular-flux discussion later in the chapter \cite{stratton1941}.

Eqs.~\eqref{eq:ch1.poynting-complex} and
\eqref{eq:ch1.maxwell-stress} presuppose harmonic averaging in passive media with
local \(\varepsilon,\mu\). Dispersive, anisotropic, or strongly time-varying
constitutive laws modify both storage terms and traction operators; active boundaries
can inject net gain that invalidates the passive-balance route used in
Subsection~\ref{sec:ch01.1.3}
\cite{stratton1941}.

\input{ch01-classical-tables}

Figure~\ref{fig:ch1.wp-reactive-vs-radiative} illustrates how the reactive--radiative distinction enters these balances.
\begin{figure}[!htbp]
\centering
\VolIPhysFig{\VolIPhysReactiveVsRadiative}
\caption{Concentric evaluation shells at $kr=0.5$, $kr=3$, and $kr=10$ for a small dipole: reactive storage dominates on the inner shell, while outward $\langle\mathbf{S}\rangle$ and steady traction emerge on the radiative shell \cite{stratton1941,chu1948}.}
\label{fig:ch1.wp-reactive-vs-radiative}
\end{figure}

Figure~\ref{fig:ch1.wp-reactive-vs-radiative} names three concentric evaluation shells at $kr=0.5$, $kr=3$, and $kr=10$ for an electrically small dipole source ($k\ell\ll1$ for source scale $\ell$). At $kr=0.5$ reactive
storage dominates: Eq.~\eqref{eq:ch1.poynting-real-imag-split} gives
\(\int_{4\pi}\Re(\mathbf{S}_\mathrm{c}\cdot\hat{\mathbf{r}})\,dS=P_{\mathrm{rad}}\ll
\int_{4\pi}|\Im(\mathbf{S}_\mathrm{c}\cdot\hat{\mathbf{r}})|\,dS=|P_{\mathrm{reac}}|\), and traction on a
small patch is oscillatory rather than proportional to \(P_{\mathrm{rad}}/c\). At
$kr=3$ the transition regime begins; at
$kr=10$ outward $\Re(\mathbf{S}_\mathrm{c})$ dominates and the shell supports radiation-pressure and torque integrals developed in
section~\ref{sec:ch01.6.4} \cite{stratton1941,chu1948}.
The conservation law holds on any admissible closed surface; a predominantly radiative shell is preferred when the purpose is to interpret the integral as outward radiated transport rather than local reactive exchange \cite{chu1948,stratton1941}.

Using the uniform plane-wave example of Subsection~\ref{sec:ch01.1.1}, take $\mathbf{S}_\mathrm{c}=(|E_0|^2/2\eta_0)\hat{\mathbf{z}}$
and evaluate power through a disk of radius $a$ normal to $\hat{\mathbf{z}}$:
\begin{equation}
\label{eq:ch1.worked-aperture-power}
P(a)=\int_0^{2\pi}\!\!\int_0^a \Re(\mathbf{S}_\mathrm{c}\cdot\hat{\mathbf{z}})\,\rho\,d\rho\,d\phi
=\frac{\pi a^2 |E_0|^2}{2\eta_0}.
\end{equation}
For a uniform plane wave, $P(a)$ in Eq.~\eqref{eq:ch1.worked-aperture-power} is the active power intercepted by a disk of radius $a$ whose normal aligns with $\langle\mathbf{S}\rangle$ \cite{poynting1884}.

Within the harmonic vacuum class treated here, Eqs.~\eqref{eq:ch1.poynting-complex} and~\eqref{eq:ch1.maxwell-stress} supply the energetic and mechanical flux quantities that Subsection~\ref{sec:ch01.1.5} rotates into angular form through $\Mavg=\mathbf{r}\times\Tavg$ \cite{stratton1941}.

\subsection{Angular momentum conservation and torque balance}
\label{sec:ch01.1.5}

Instantaneous angular momentum follows the same Maxwell pair as linear momentum. From Eq.~\eqref{eq:ch1.momentum-balance-instantaneous}, multiply by the Levi-Civita symbol and position vector to obtain, in index notation,
\begin{equation}
\label{eq:ch1.angular-momentum-index-start}
\varepsilon_{ijk}\,r_j\!\left(\partial_t g_k+f_{\mathrm{src},k}-\partial_l T_{kl}\right)
=0.
\end{equation}
The product rule expands the stress term as
\begin{equation}
\label{eq:ch1.angular-momentum-product-rule}
\varepsilon_{ijk}\,r_j\,\partial_l T_{kl}
=\partial_l\!\left(\varepsilon_{ijk}\,r_j T_{kl}\right)
-\varepsilon_{ijk}\,T_{kj},
\end{equation}
and symmetry of the Maxwell stress tensor, $T_{kj}=T_{jk}$, makes the contraction
$\varepsilon_{ijk}T_{kj}=0$ \cite{maxwell1873,stratton1941}. Define the instantaneous
stress-moment density $\Minst=\mathbf{r}\times\Tinst$ and angular-momentum density
$\mathbf{j}_{\mathrm{em}}=\mathbf{r}\times\mathbf{g}_{\mathrm{em}}$. Eqs.~\eqref{eq:ch1.angular-momentum-index-start}--\eqref{eq:ch1.angular-momentum-product-rule} then yield
\begin{equation}
\label{eq:ch1.angular-momentum-balance}
\frac{\partial\mathbf{j}_{\mathrm{em}}}{\partial t}+\boldsymbol{\tau}_{\mathrm{src}}
=\divg\Minst,
\end{equation}
with source-exchange torque density $\boldsymbol{\tau}_{\mathrm{src}}=\mathbf{r}\times\mathbf{f}_{\mathrm{src}}$ (N\,m\,m$^{-3}$) obtained from the same $\rho$ and $\Jimp$ as Eq.~\eqref{eq:ch1.momentum-balance-instantaneous} \cite{maxwell1873,stratton1941}. Mechanical torque on carriers is $\boldsymbol{\tau}_{\mathrm{matter}}=\mathbf{r}\times\mathbf{f}_{\mathrm{matter}}$ of Eq.~\eqref{eq:ch1.lorentz-force}; $\boldsymbol{\tau}_{\mathrm{src}}=\boldsymbol{\tau}_{\mathrm{matter}}$ only under Route~A on the support where $\mathbf{f}_{\mathrm{src}}=\mathbf{f}_{\mathrm{matter}}$. The resultant mechanical torque $\mathbf{T}_{\mathrm{matter}}=\int_{V_m}\boldsymbol{\tau}_{\mathrm{matter}}\,dV$ (N\,m) follows from Eq.~\eqref{eq:ch1.angular-balance-integral} under Route~A throughout $V_m$, or from the closed stress integral on $\partial V_m$ in source-free vacuum under Route~B (Eq.~\eqref{eq:ch1.torque-stress-surface}) \cite{stratton1941,beth1936}. In harmonic steady state with vanishing storage rate, $\mathbf{T}_{\mathrm{matter}}=\boldsymbol{\mathcal{T}}_J=-\boldsymbol{\Phi}_{J,\mathrm{out}}$ when the integration domain and route assumptions support that identification \cite{beth1936}. Eq.~\eqref{eq:ch1.angular-momentum-balance} is the instantaneous rotational counterpart of Eq.~\eqref{eq:ch1.momentum-balance-instantaneous}: stored field angular momentum changes when source-exchange torque density acts or when stress-moment flux diverges.

For time-harmonic fields under $\exp(+\jj\omega t)$, the same conservation laws may be cycle-averaged. Averaging Eq.~\eqref{eq:ch1.momentum-balance-instantaneous} gives $\partial_t\langle\mathbf{g}_{\mathrm{em}}\rangle+\langle\mathbf{f}_{\mathrm{src}}\rangle=\divg\Tavg$, and averaging Eq.~\eqref{eq:ch1.angular-momentum-balance} gives
\begin{equation}
\label{eq:ch1.angular-momentum-balance-averaged}
\frac{\partial\langle\mathbf{j}_{\mathrm{em}}\rangle}{\partial t}+\langle\boldsymbol{\tau}_{\mathrm{src}}\rangle
=\divg\Mavg,
\end{equation}
where $\Mavg=\mathbf{r}\times\Tavg$, $\langle\mathbf{j}_{\mathrm{em}}\rangle=\mathbf{r}\times\langle\mathbf{g}_{\mathrm{em}}\rangle$, and $\langle\boldsymbol{\tau}_{\mathrm{src}}\rangle=\mathbf{r}\times\langle\mathbf{f}_{\mathrm{src}}\rangle$ \cite{stratton1941}. Under Route~A, $\langle\boldsymbol{\tau}_{\mathrm{matter}}\rangle=\langle\boldsymbol{\tau}_{\mathrm{src}}\rangle$ on the relevant support; under Route~B, $\langle\boldsymbol{\tau}_{\mathrm{src}}\rangle$ remains field--source bookkeeping and mechanical torque on a body is read from the stress surface integral (Eq.~\eqref{eq:ch1.torque-stress-surface}). Angle brackets mark cycle-averaged harmonic quantities; $\Minst$ and $\Mavg$ are not interchangeable in the same displayed balance.

In linear isotropic vacuum the cycle-averaged momentum density is
\begin{equation}
\label{eq:ch1.rotation-generator-split}
\langle\mathbf{g}_{\mathrm{em}}\rangle=\varepsilon_0\mu_0\,\langle\mathbf{S}\rangle
=\varepsilon_0\mu_0\,\Re(\mathbf{S}_\mathrm{c}),
\end{equation}
and the corresponding angular-momentum density is
\begin{equation}
\label{eq:ch1.angular-momentum-density}
\langle\mathbf{j}_{\mathrm{em}}\rangle=\mathbf{r}\times\langle\mathbf{g}_{\mathrm{em}}\rangle.
\end{equation}
\addtocounter{equation}{-1}\refstepcounter{equation}\label{eq:ch1.angular-density-split}%

Integrating Eq.~\eqref{eq:ch1.angular-momentum-balance-averaged} over a control volume $V_m$ with outward normal $\hat{\mathbf{n}}$ gives a valid volume identity only where that local balance holds pointwise (vacuum $\varepsilon_0,\mu_0$ region with the same $(\rho,\Jimp)$ as Maxwell's equations). When $V_m$ encloses matter, the surface balance Eqs.~\eqref{eq:ch1.torque-stress-surface} and~\eqref{eq:ch1.outward-am-flux} on $\partial V_m$ in source-free vacuum supply the mechanical resultant under Route~B; the volume integral of $\langle\boldsymbol{\tau}_{\mathrm{src}}\rangle$ is then source-exchange bookkeeping, not automatic load torque \cite{stratton1941}. With that scope,
\begin{equation}
\label{eq:ch1.angular-balance-integral}
\frac{d}{dt}\int_{V_m} \langle\mathbf{j}_{\mathrm{em}}\rangle\,dV
+\int_{V_m} \langle\boldsymbol{\tau}_{\mathrm{src}}\rangle\,dV
=\oint_{\partial V_m}\Mavg\cdot\hat{\mathbf{n}}\,dS.
\end{equation}
The volume term $\int_{V_m}\langle\boldsymbol{\tau}_{\mathrm{src}}\rangle\,dV$ is the integrated source-exchange torque associated with $(\rho,\Jimp)$; it equals $\int_{V_m}\langle\boldsymbol{\tau}_{\mathrm{matter}}\rangle\,dV$ only under Route~A throughout $V_m$ \cite{stratton1941}. The surface term is the stress moment $\boldsymbol{\mathcal{T}}_J=\oint\Mavg\cdot\hat{\mathbf{n}}\,dS$ \cite{stratton1941}. At harmonic steady state with $\partial_t\langle\mathbf{j}_{\mathrm{em}}\rangle=\mathbf{0}$, Eq.~\eqref{eq:ch1.torque-stress-surface} and~\eqref{eq:ch1.outward-am-flux} give the \emph{resultant} mechanical torque $\mathbf{T}_{\mathrm{matter}}=\boldsymbol{\mathcal{T}}_J-d\mathbf{J}_{\mathrm{em,store}}/dt=-\boldsymbol{\Phi}_{J,\mathrm{out}}-d\mathbf{J}_{\mathrm{em,store}}/dt$ when the enclosing surface lies in source-free vacuum and the route assumptions are satisfied (N\,m), not the local density $\boldsymbol{\tau}_{\mathrm{src}}$ or $\boldsymbol{\tau}_{\mathrm{matter}}$ \cite{beth1936}. Beth \cite{beth1936} measured load torque from net flux through the plate's control surface (incident minus transmitted), not from equating an impressed-source volume integral to load mechanics without that surface balance \cite{stratton1941}.

Figure~\ref{fig:ch1.momentum-angular-geometry} shows the cycle-averaged chain at a field point and on a surface element: $\langle\mathbf{g}_{\mathrm{em}}\rangle$, $\langle\mathbf{j}_{\mathrm{em}}\rangle=\mathbf{r}\times\langle\mathbf{g}_{\mathrm{em}}\rangle$, traction $\Tavg\cdot\hat{\mathbf{n}}$, and stress moment $\mathbf{r}\times(\Tavg\cdot\hat{\mathbf{n}})$ \cite{maxwell1873,beth1936}.

\begin{figure}[t]
\centering
\VolIPhysFigWide{\VolIPhysMomentumAngularGeometry}
\caption{Cycle-averaged angular-momentum geometry: at observation point $\mathbf{r}$, $\langle\mathbf{g}_{\mathrm{em}}\rangle$ and $\langle\mathbf{j}_{\mathrm{em}}\rangle=\mathbf{r}\times\langle\mathbf{g}_{\mathrm{em}}\rangle$; on enclosing surface element with outward $\hat{\mathbf{n}}$, traction $\Tavg\cdot\hat{\mathbf{n}}$ and stress moment $\mathbf{r}\times(\Tavg\cdot\hat{\mathbf{n}})$ that enters torque integrals \cite{maxwell1873,beth1936,stratton1941}.}
\label{fig:ch1.momentum-angular-geometry}
\end{figure}

The cycle-averaged stress-moment density $\Mavg=\mathbf{r}\times\Tavg$ satisfies
$\Mavg\cdot\hat{\mathbf{n}}=\mathbf{r}\times(\Tavg\cdot\hat{\mathbf{n}})$ on $\partial V$ (stress-moment surface density).
At harmonic steady state the cycle-averaged \emph{resultant} mechanical torque on matter enclosed in $V_m$ is
\begin{equation}
\label{eq:ch1.torque-stress-surface}
\mathbf{T}_{\mathrm{matter}}
\equiv
\int_{V_m}\langle\boldsymbol{\tau}_{\mathrm{matter}}\rangle\,dV
=\oint_{\partial V_m}\Mavg\cdot\hat{\mathbf{n}}\,dS
-\frac{d}{dt}\int_{V_m} \langle\mathbf{j}_{\mathrm{em}}\rangle\,dV
=\oint_{\partial V_m}\mathbf{r}\times(\Tavg\cdot\hat{\mathbf{n}})\,dS
-\frac{d}{dt}\int_{V_m} \langle\mathbf{j}_{\mathrm{em}}\rangle\,dV,
\end{equation}
with outward $\hat{\mathbf{n}}$ on $\partial V_m$ \cite{beth1936,stratton1941}. The volume identity $\int_{V_m}\langle\boldsymbol{\tau}_{\mathrm{matter}}\rangle\,dV=\int_{V_m}\langle\boldsymbol{\tau}_{\mathrm{src}}\rangle\,dV$ holds under Route~A throughout $V_m$; under Route~B, $\mathbf{T}_{\mathrm{matter}}$ is obtained from the stress surface integral when $\partial V_m$ lies in source-free vacuum. The stress moment on the same closed surface is
\begin{equation}
\label{eq:ch1.stress-moment-integral}
\boldsymbol{\mathcal{T}}_J
\equiv
\oint_{\partial V_m}\mathbf{r}\times(\Tavg\cdot\hat{\mathbf{n}})\,dS
=\oint_{\partial V_m}\Mavg\cdot\hat{\mathbf{n}}\,dS.
\end{equation}
Eq.~\eqref{eq:ch1.torque-stress-surface} gives $\mathbf{T}_{\mathrm{matter}}=\boldsymbol{\mathcal{T}}_J-d\mathbf{J}_{\mathrm{em,store}}/dt$ with stored field angular momentum $\mathbf{J}_{\mathrm{em,store}}=\int_{V_m}\langle\mathbf{j}_{\mathrm{em}}\rangle\,dV$. Outward electromagnetic angular-momentum flux through $\partial V_m$ is
\begin{equation}
\label{eq:ch1.outward-am-flux}
\boldsymbol{\Phi}_{J,\mathrm{out}}
\equiv
\dot{\mathbf{J}}_{\mathrm{em,out}}
=-\boldsymbol{\mathcal{T}}_J
=-\oint_{\partial V_m}\mathbf{r}\times(\Tavg\cdot\hat{\mathbf{n}})\,dS,
\end{equation}
so $\mathbf{T}_{\mathrm{matter}}=-\boldsymbol{\Phi}_{J,\mathrm{out}}-d\mathbf{J}_{\mathrm{em,store}}/dt$. The stress moment $\boldsymbol{\mathcal{T}}_J$ is not the outward flux; outward transport is $\boldsymbol{\Phi}_{J,\mathrm{out}}=-\boldsymbol{\mathcal{T}}_J$ \cite{stratton1941,beth1936}. In harmonic steady state with $d\mathbf{J}_{\mathrm{em,store}}/dt=\mathbf{0}$, $\mathbf{T}_{\mathrm{matter}}=\boldsymbol{\mathcal{T}}_J=-\boldsymbol{\Phi}_{J,\mathrm{out}}$; source and load resultants $\boldsymbol{\tau}_{\mathrm{source}}=-\boldsymbol{\Phi}_{J,\mathrm{out}}$ and $\boldsymbol{\tau}_{\mathrm{load}}=+\boldsymbol{\Phi}_{J,\mathrm{out}}$ apply only when the declared control surface encloses the support and the load intercepts the net flux \cite{beth1936}.

Belinfante's 1940 field-theoretic analysis \cite{belinfante1940} showed that canonical
energy--momentum and symmetric (kinetic) tensors differ by the divergence of a spin
current; correspondingly, total angular momentum
$\int_V \mathbf{j}_{\mathrm{em}}\,dV$ and closed-surface flux
$\oint \Mavg\cdot\hat{\mathbf{n}}\,dS$ are gauge invariant, whereas local spin--orbit splits are convention-dependent \cite{belinfante1940}. Primary quantities are $\mathbf{j}_{\mathrm{em}}=\mathbf{r}\times\mathbf{g}_{\mathrm{em}}$ and stress-moment flux through $\Mavg\cdot\hat{\mathbf{n}}$. A separate
``residual spin'' obtained by subtracting these identical definitions is not used.
Detailed spin--orbit and canonical decompositions are representation-dependent and are
developed where the modal basis is fixed (Chapter~\ref{ch:03}) \cite{belinfante1940,tai1994}.
Beth's experiment on a wave plate under circular polarization
\cite{beth1936} showed that the spin-associated part carries mechanical consequences; later orbital-indexed beam experiments
\cite{allen1992} extended the vocabulary to spatial phase gradients; the measurable
balance remains Eq.~\eqref{eq:ch1.angular-balance-integral} \cite{belinfante1940}.

For torque about axis $\hat{\mathbf{u}}$, project Eq.~\eqref{eq:ch1.angular-balance-integral}:
\begin{equation}
\label{eq:ch1.torque-observable}
\tau_{\hat{\mathbf{u}}}(V_m)
=\hat{\mathbf{u}}\cdot
\left(
\oint_{\partial V_m}\Mavg\cdot\hat{\mathbf{n}}\,dS
-\frac{d}{dt}\int_{V_m} \langle\mathbf{j}_{\mathrm{em}}\rangle\,dV
\right),
\end{equation}
which is the axis-projected form of Eq.~\eqref{eq:ch1.angular-balance-integral}.
Torque is therefore a projected surface integral of the stress moment, not a scalar read from a phase portrait \cite{beth1936,stratton1941}.

To compare directional concentration of stress moment on a sphere of radius $r$, note that $\avgStressMoment\cdot\hat{\mathbf{r}}$ states how the local stress moment projects onto each solid-angle sector. A main-lobe concentration in $\|\avgStressMoment\cdot\hat{\mathbf{r}}\|$ does not by itself establish net torque about an axis until Eq.~\eqref{eq:ch1.torque-observable} is evaluated with the declared origin and surface \cite{stratton1941,beth1936}.
\phantomsection\label{eq:ch1.angular-flux-directional-index}

Torque about axis $\hat{\mathbf{u}}$ is fixed once origin $\mathbf{r}_0$, axis $\hat{\mathbf{u}}$, and control surface $\partial V$ entering $\Mavg=\mathbf{r}\times\Tavg$ are declared \cite{stratton1941}. Changing any of these geometric choices rotates flux projections without altering the underlying Maxwell pair; published torque values therefore state $(\mathbf{r}_0,\hat{\mathbf{u}},\partial V)$ alongside the field solution \cite{beth1936}.

Consider two Maxwell field states \(\mathcal{F}_1,\mathcal{F}_2\) at common radius
\(r\), with equal total radiated power but different directional stress-moment
densities \(\mathbf{M}_1\cdot\hat{\mathbf{r}}\neq
\mathbf{M}_2\cdot\hat{\mathbf{r}}\) on a subset of $4\pi$ of nonzero measure. Equal $P_{\mathrm{rad}}$ constrains only the energetic integral; it does not fix the stress-moment distribution direction by direction. Projected torque through Eq.~\eqref{eq:ch1.torque-observable} can therefore differ between the two states even when their radiation patterns share the same total power \cite{stratton1941}.

\phantomsection\label{prop:ch1.torque-observable-criterion}
\begin{figure}[t]
\centering
\VolIPhysFig{\VolIPhysTorqueBalanceSurface}
\caption{Spherical control shell $S_R$ with axis $\hat{\mathbf{z}}$: axisymmetric radiation with azimuthal order $m=0$ may carry zero $\dot J_{z,\mathrm{out}}$ from the stress-moment integral even when $P_{\mathrm{rad}}$ is finite; circular polarization fixes local handedness but does not by itself guarantee nonzero $\dot J_{z,\mathrm{out}}$ on $S_R$ under Eq.~\eqref{eq:ch1.outward-am-flux}. Azimuthal modes with $m\neq 0$ can break symmetry and yield nonzero net flux when the Maxwell-valid field and surface integral support it \cite{beth1936,stratton1941}.}
\label{fig:ch1.wp-torque-balance-surface}
\end{figure}

Figure~\ref{fig:ch1.wp-torque-balance-surface} contrasts symmetric radiation with azimuthally phased
radiation on the same spherical shell: when the stress flux integrand possesses definite azimuthal
symmetry about $\hat{\mathbf{z}}$, the solid-angle integral of
$[\mathbf{r}\times(\Tavg\cdot\hat{\mathbf{n}})]\cdot\hat{\mathbf{z}}$ may vanish even though
local phase structure is visible on a near-field display---the symmetry argument Beth exploited for
circular polarization \cite{beth1936,stratton1941}.

For a matter-enclosing control volume $V_m$ whose stress-evaluation surface lies in source-free vacuum,
Eq.~\eqref{eq:ch1.angular-balance-integral} reduces at harmonic steady state to
\(T_{z,\mathrm{matter}}=\oint_{\partial V_m}[\mathbf{r}\times(\Tavg\cdot\hat{\mathbf{n}})]\cdot\hat{\mathbf{z}}\,dS\)
and \(\dot J_{z,\mathrm{em,out}}=-T_{z,\mathrm{matter}}\) by Eq.~\eqref{eq:ch1.outward-am-flux} (here $T_{z,\mathrm{matter}}$ is the resultant torque in N\,m, not the local density $\boldsymbol{\tau}_{\mathrm{src}}$ of Eq.~\eqref{eq:ch1.angular-momentum-balance}).
On sphere \(S_R\), use \(\mathbf{r}=R\hat{\mathbf{r}}\) and
\(dS=R^2\sin\theta\,d\theta\,d\phi\). For axisymmetric radiation with azimuthal index $m=0$, the far-zone pattern carries no azimuthal phase circulation about $\hat{\mathbf{z}}$, so $\dot J_{z,\mathrm{out}}$ from orbital winding vanishes while $P_{\mathrm{rad}}$ remains finite. For a uniform circularly polarized plane wave, the symmetric Maxwell stress tensor on a centered transverse disk gives $\dot J_{z,\mathrm{em,out}}=0$ while power crossing the disk is nonzero (Subsection~\ref{sec:ch01.6.1}); $(m+\sigma)P/\omega$ and $\sigma P/\omega$ symbols elsewhere are modal-class evaluators unless derived from a Maxwell-valid field on the declared surface \cite{beth1936,stratton1941}.
Torque is an integral consequence of the Maxwell stress tensor through Eq.~\eqref{eq:ch1.torque-observable}, not a local phase label
\cite{stratton1941}.

Eqs.~\eqref{eq:ch1.angular-momentum-balance}
and \eqref{eq:ch1.angular-balance-integral} assume a fixed origin, a symmetric
stress-tensor completion, and local linear constitutive behavior in the present
regime; strongly nonlocal or dissipative media may require additional closure before
torque partition is interpreted without ambiguity \cite{stratton1941}.

\section{Angular Structure of Electromagnetic Fields}
\label{sec:ch01.2}

This section applies the flux conventions of Subsection~\ref{sec:ch01.1.5} to angular organization of $(\E,\H)$: phase and polarization describe local structure; signed transport follows from $\oint \Mavg\cdot\hat{\mathbf{n}}\,dS$ under the separable conditions stated where used \cite{allen1992,stratton1941,beth1936}.
\subsection{Rotating source versus rotating pattern}
\label{sec:ch01.2.0}

A spiral isophase map on outgoing sphere $S_R$ does not by itself fix radiated angular momentum, nor does it distinguish a mechanically rotating feed from a stationary phased array \cite{beth1936}. In one configuration the source---currents, polarization basis, or feed geometry---rotates in space about \(\hat{\mathbf{z}}\); in the other the feed stays fixed while an azimuthally progressive law \(\exp(\jj m\phi)\) is imposed by phasing or modal synthesis \cite{harrington2001}. Both may show the same spiral contours on $S_R$; Maxwell stress integrals and torque on matter decide whether rotational action is transported \cite{beth1936,stratton1941}. Equal counter-rotating superpositions with matched power cancel net signed flux while retaining strong phase structure (Subsection~\ref{sec:ch01.6.2}) \cite{allen1992}.

Figure~\ref{fig:ch1.signature-rotation-duality} compares physical source rotation with a stationary feed carrying azimuthally progressive phase; both can yield the same spiral isophase pattern on \(S_R\), yet only the flux--torque pair of Eqs.~\eqref{eq:ch1.directional-am-flux} and~\eqref{eq:ch1.torque-observable} distinguishes radiated action from appearance.

\begin{figure}[t]
\centering
\VolIPhysFig{\VolIPhysRotationDuality}
\caption{Physical source rotation (a) versus stationary phased feed (b): both can yield identical spiral isophase contours on far-zone shell $S_R$, yet signed flux and torque observables distinguish radiated angular action from appearance alone.}
\label{fig:ch1.signature-rotation-duality}
\end{figure}

Panel~(a) applies active rotation to a compact electric dipole \(\mathbf{p}\) as in Subsection~\ref{sec:ch01.0.2}. A mechanically rotating source is generally time-dependent and can produce sidebands or other harmonic content depending on how the motion is modeled; the single-frequency phasor at fixed \(\omega\) describes only a frozen orientation, a separated harmonic component, or another explicitly stated single-$\omega$ approximation \cite{stratton1941}.
Panel~(b) keeps the feed stationary and sets aperture phasing
\(\E_a(\phi)=\E_0\exp(\jj m\phi)\hat{\boldsymbol{\rho}}\) on a ring of radius
\(a\ll\lambda\). In the paraxial limit the same azimuthal index \(m\) appears in
both cases, yet the source angular momentum balance differs: panel~(a) requires a
motorized or mechanically constrained feed; panel~(b) can be realized with a passive
corrugated phased feed network whose elements are fixed in space \cite{harrington2001}.

On sphere \(S_R\) in the radiative zone, define the directional observables of
Eqs.~\eqref{eq:ch1.radiation-intensity-Ur}--\eqref{eq:ch1.directional-am-flux}.
For configuration~(b) with single azimuthal index \(m\) and axisymmetric magnitude,
\begin{equation}
\label{eq:ch1.signature-phase-ramp}
\E(R,\theta,\phi)\propto \E_\theta(\theta)\,e^{\jj m\phi},
\end{equation}
a schematic azimuthal phase label on $S_R$ for configuration~(b). Signed transport $\dot J_{z,\mathrm{em,out}}$ is fixed only by Eq.~\eqref{eq:ch1.ex61-flux-from-stress} (or Eq.~\eqref{eq:ch1.signature-flux-torque-pair}) for a Maxwell-valid $(\E,\H)$ on the declared surface; Eq.~\eqref{eq:ch1.signature-phase-ramp} alone does not imply $\dot J_{z,\mathrm{em,out}}=(m/\omega)P_{\mathrm{rad}}$ \cite{stratton1941,allen1992}. For configuration~(a) at harmonic frequency
\(\omega\) with slow rotation \(\Omega_s\ll\omega\), the radiated field is a
sideband superposition; when \(\Omega_s\) is commensurate with the phasor period the
observed map on \(S_R\) can match Eq.~\eqref{eq:ch1.signature-phase-ramp} at a
snapshot, but the time-averaged torque on a Beth-type absorber includes the
mechanical power supplied to rotate the feed. Beth's experiment used a suspended half-wave plate; reversal of the circular polarization changes the angular momentum carried by the field, and the corresponding mechanical reaction appears as the measured torque on the plate \cite{beth1936}. A quarter-wave plate that converts circular polarization to linear polarization is a separate polarization-conversion example: it measures \(\tau_{z,\mathrm{load}}\) through Eq.~\eqref{eq:ch1.torque-observable} without invoking Beth's original arrangement. The rotating-source case transfers torque from feed mechanics and from field flux, whereas the phase-pattern case transfers torque through field flux alone when the feed is passive \cite{stratton1941}.

Closed-surface integrals, not the phase portrait alone, fix transport:
\begin{equation}
\label{eq:ch1.signature-flux-torque-pair}
\Phi_{J,z,\mathrm{out}}(R)
=-\hat{\mathbf{z}}\cdot\int_{4\pi}\mathbf{j}_J(\theta,\phi;R)\,d\Omega
=-\hat{\mathbf{z}}\cdot\boldsymbol{\Phi}_{J,\mathrm{out}}(R),
\qquad
\dot J_{z,\mathrm{em,out}}(R)=\Phi_{J,z,\mathrm{out}}(R),
\end{equation}
\begin{equation}
\label{eq:ch1.signature-source-torque-z}
\tau_{z,\mathrm{source}}(R)
=\oint_{S_R}\big[\mathbf{r}\times(\Tavg\cdot\hat{\mathbf{n}})\big]\cdot\hat{\mathbf{z}}\,dS
=-\Phi_{J,z,\mathrm{out}}(R)
\end{equation}
at harmonic steady state when storage in the enclosing volume averages to zero \cite{beth1936,stratton1941}.
Configuration~(b) with \(|m|\ge 1\) and \(|P_{\mathrm{reac}}(R)|/P_{\mathrm{rad}}\ll 1\) on \(S_R\) describes a predominantly radiative shell, but net signed \(\Phi_{J,z,\mathrm{out}}(R)\) and \(\tau_{z,\mathrm{source}}(R)\) follow only from evaluating Eqs.~\eqref{eq:ch1.signature-flux-torque-pair} and~\eqref{eq:ch1.signature-source-torque-z} for the Maxwell-valid \((\E,\H)\) on that surface; azimuthal index and a small reactive-power ratio do not by themselves guarantee nonzero transport \cite{stratton1941,allen1992}.
Configuration~(a) matched to the same snapshot pattern can share \(U_r(\theta,\phi;R)\)
yet carry additional mechanical boundary terms on the feed support that do not appear
in a pure phase-synthesis model; omitting those terms violates
Eq.~\eqref{eq:ch1.angular-balance-integral}. Conversely, a balanced pair of
counter-rotating phase ramps \(\exp(\pm\jj m\phi)\) with equal power superpose to a standing azimuthal pattern \(\sqrt{2}\cos(m\phi)\); under the fixed linear-polarization and modal assumptions of section~\ref{sec:ch01.6.2}, the signed stress integral in Eq.~\eqref{eq:ch1.directional-am-total} yields \(\boldsymbol{\Phi}_{J,\mathrm{out}}(R)=\mathbf{0}\) for that template---an orbital-bookkeeping counterexample, not a universal cancellation of all spin-plus-orbital contributions for arbitrary polarization \cite{allen1992,stratton1941}.

Azimuthal index $m$ organizes phase on a chart; Eqs.~\eqref{eq:ch1.signature-flux-torque-pair} measure transport. Panels~(a) and~(b) cannot be separated by phase on $S_R$ alone, but $U_r$, the stress integrand, and $\tau_{z,\mathrm{source}}$ on the same surface can \cite{stratton1941,beth1936}.
\subsection{Directional dependence and angular momentum}
\label{sec:ch01.2.1}


Radiation intensity $U_r$ and the Maxwell-stress angular-momentum integrand tie Poynting transport to torque on declared surfaces \cite{poynting1884,belinfante1940,beth1936}. Modal labels are not substitutes for those integrals \cite{stratton1941}.

Radiated power and radiated angular momentum are resolved on the same sphere but through different field functionals \cite{stratton1941,beth1936}. An axisymmetric dipole may carry finite $P_{\mathrm{rad}}$ yet vanishing net $z$-flux when opposing sectors cancel \cite{stratton1941}. Let \(\hat{\mathbf{r}}=\hat{\mathbf{r}}(\theta,\phi)\) be radial unit vector on a
radius-\(r\) sphere and \(d\Omega=\sin\theta\,d\theta\,d\phi\) the solid-angle
measure. With \(\mathbf{S}_\mathrm{c}\) from Eq.~\eqref{eq:ch1.intro-poynting}, the radiation intensity is
\begin{equation}
\label{eq:ch1.radiation-intensity-Ur}
U_r(r,\theta,\phi)=
r^2\,\Re\!\left\{\mathbf{S}_\mathrm{c}(r,\theta,\phi)\cdot\hat{\mathbf{r}}\right\},
\end{equation}
in watts per steradian \cite{poynting1884,harrington2001}. The far-zone limit is $U(\theta,\phi)=\lim_{r\to\infty}U_r$, total radiated power is
\begin{equation}
\label{eq:ch1.directional-power-integral}
P_{\mathrm{rad}}(r)=\int_{4\pi}U_r(r,\theta,\phi)\,d\Omega,
\end{equation}
and directivity is
\begin{equation}
\label{eq:ch1.directivity-definition}
D(\theta,\phi)=\frac{4\pi\,U(\theta,\phi)}{P_{\mathrm{rad}}}.
\end{equation}
Equal integrated power does not imply equal directional angular-momentum structure \cite{stratton1941}.

For rotational transport, define from
Eq.~\eqref{eq:ch1.angular-momentum-balance}
\begin{equation}
\label{eq:ch1.directional-am-flux}
\mathbf{j}_J(\theta,\phi;r)=
r^2\,\big[\mathbf{r}\times(\Tavg\cdot\hat{\mathbf{r}})\big].
\end{equation}
Here $\mathbf{j}_J$ is the directional stress-moment density (N\,m/sr): solid-angle resolution of $\Mavg=\mathbf{r}\times\Tavg$ on a sphere of radius $r$ \cite{beth1936,stratton1941}. Radiation intensity $U_r$ from Eq.~\eqref{eq:ch1.radiation-intensity-Ur} describes energetic transport; $\mathbf{j}_J$ describes the corresponding rotational transport density.

The shell integrals are
\begin{equation}
\label{eq:ch1.directional-am-total}
\boldsymbol{\mathcal{T}}_J(r)=
\int_{4\pi}\mathbf{j}_J(\theta,\phi;r)\,d\Omega,
\qquad
\boldsymbol{\Phi}_{J,\mathrm{out}}(r)=-\boldsymbol{\mathcal{T}}_J(r).
\end{equation}
This vector integral cannot replace direction-resolved structure because cancellation among
opposing solid-angle sectors can leave a small total despite large directional flux
magnitudes.

A direction with outward \(U_r>0\) need not carry corresponding signed angular-momentum transport; the stress integrand and \(U_r\) must be evaluated separately on the same shell \cite{stratton1941,belinfante1940}.

\phantomsection\label{prop:ch1.directional-nonequivalence}
Let two radiative states \(\mathcal{F}_1,\mathcal{F}_2\) at radius \(r\) satisfy
\(P_{\mathrm{rad}}^{(1)}(r)=P_{\mathrm{rad}}^{(2)}(r)\). If there exists a set
\(\mathcal{U}\subset4\pi\) of nonzero measure where
\(\mathbf{j}_J^{(1)}\neq\mathbf{j}_J^{(2)}\) while both fields carry the same \(P_{\mathrm{rad}}\) on \(\mathcal{U}\), then the states are rotationally non-equivalent for
direction-selective interactions.

Equal \(P_{\mathrm{rad}}\) does not fix direction-resolved angular-momentum transport.
Weighting $\mathbf{j}_J$ differently over solid angle changes the projected torque about a given axis even when total radiated power is unchanged \cite{stratton1941}.

Integrated totals are terminal reductions of richer direction-resolved fields. Different solid-angle topology
can produce nearly identical totals while yielding different torque outcomes.

Origin declaration remains essential because angular transport is moment-based.
Changing origin re-weights directional moments even when global balances remain
unchanged. Directional flux and torque statements therefore require the same origin convention established in section~\ref{sec:ch01.1.5}.

Directional observables represent physical transport
only when outward flux is well-defined on the chosen shell and Eq.~\eqref{eq:ch1.am-net-transport-condition} holds for the field under study.
In strongly reactive near zones, strongly nonlocal media, or rapidly modulated
boundaries, extra decomposition terms may be necessary before \(U_r\) and
\(\mathbf{j}_J\) are treated as direct rotational-transfer observables
\cite{stratton1941}.

For fully polarized harmonic transverse fields, the two independent electric
components \(E_\theta\) and \(E_\phi\) on a chosen surface fix the oscillation
ellipse. Left- and right-circular polarization correspond to \(\E\) and \(\H\)
tracing the same rotation sense about the propagation direction
\cite{stratton1941}. Beth's quartz-plate experiment \cite{beth1936} is the
mechanical consequence when such a field transfers angular momentum to matter
through Maxwell stress at the interface.
\phantomsection\label{eq:ch1.polarization-ellipse-angles}
\phantomsection\label{eq:ch1.jones-rotation}
\phantomsection\label{eq:ch1.coherency-matrix}
\phantomsection\label{eq:ch1.polarization-degree}

Directional dependence is a question of solid-angle structure: which sheet on $S^{2}$ carries flux, and does that sheet carry nonzero signed angular-momentum flux before any Beth or OAM label is attached \cite{belinfante1940}? Circular polarization fixes the handedness of the transverse \(\E\)--\(\H\) triad locally; outward angular flux still requires a signed stress integral on an enclosing surface \cite{stratton1941,beth1936}. Helical phase topology is separated from radiated action below \cite{allen1992}.

\subsection{Helical phase structure and transport}
\label{sec:ch01.2.2}

Helical phase on a transverse slice fixes local winding of $\arg E$; net transport is the signed solid-angle integral of angular-flux density on a control surface enclosing the source (Figure~\ref{fig:ch1.phase-vs-transport-sketch}) \cite{belinfante1940,allen1992}. A $\lambda/4$ plate can change $\oint \Mavg\cdot\hat{\mathbf{n}}\,dS$ on $\Gamma_R$ while a near-field phase portrait appears unchanged \cite{allen1992,beth1936}. Counterexamples with matched modal powers appear in Subsection~\ref{sec:ch01.6.2}.

On the left of Figure~\ref{fig:ch1.phase-vs-transport-sketch}, helical isophases show how $\arg E$ winds about an axis. On the right, control surface $\Gamma_R$ encloses the source; outward arrows mark opposing angular-flux contributions whose solid-angle integral $\boldsymbol{\Phi}_{J,\mathrm{out}}=-\int_{4\pi}\mathbf{j}_J\,d\Omega$ may vanish when counter-rotating sectors cancel---even though the transverse phase portrait looks identical to a transporting beam \cite{allen1992,beth1936}.

\begin{figure}[t]
\centering
\ChOneFigBlock{%
\begin{tikzpicture}[ch1geo]
 \begin{scope}[shift={(-2.1,0)}]
 \draw[ch1axis] (-0.95,0) -- (0.95,0) node[right, font=\scriptsize] {$x$};
 \draw[ch1axis] (0,-0.95) -- (0,0.95) node[above, font=\scriptsize] {$y$};
 \draw[ch1shellO, dotted] (0,0) circle (0.12);
 \foreach \r in {0.28,0.48,0.68} {
   \draw[gray!55!black, line width=0.55pt] (0,0) circle (\r);
 }
 \draw[-{Latex[length=1.5mm,width=1mm]}, gray!45!black, line width=0.5pt]
   (0.75,0) arc [start angle=0, end angle=35, radius=0.75];
 \draw[-{Latex[length=1.5mm,width=1mm]}, gray!50!black, line width=0.5pt]
   ({0.75*cos(60)},{0.75*sin(60)}) arc [start angle=60, end angle=95, radius=0.75];
 \draw[-{Latex[length=1.5mm,width=1mm]}, gray!50!black, line width=0.5pt]
   ({0.75*cos(120)},{0.75*sin(120)}) arc [start angle=120, end angle=155, radius=0.75];
 \draw[-{Latex[length=1.5mm,width=1mm]}, gray!50!black, line width=0.5pt]
   ({0.75*cos(180)},{0.75*sin(180)}) arc [start angle=180, end angle=215, radius=0.75];
 \draw[-{Latex[length=1.5mm,width=1mm]}, gray!50!black, line width=0.5pt]
   ({0.75*cos(240)},{0.75*sin(240)}) arc [start angle=240, end angle=275, radius=0.75];
 \draw[-{Latex[length=1.5mm,width=1mm]}, gray!50!black, line width=0.5pt]
   ({0.75*cos(300)},{0.75*sin(300)}) arc [start angle=300, end angle=335, radius=0.75];
 \node[ch1oncanvas, font=\scriptsize, align=center] at (0,-1.15)
 {winding of $\arg E$};
 \end{scope}
 \begin{scope}[shift={(2.1,0)}]
 \draw[ch1shell] (0,0) circle (0.72);
 \fill (0,0) circle (1.5pt);
 \draw[-{Latex[length=2mm,width=1.3mm]}, gray!45!black, line width=0.65pt]
   (0.72,0) -- (1.05,0);
 \draw[-{Latex[length=2mm,width=1.3mm]}, gray!65!black, line width=0.65pt]
   (0,0.72) -- (0,1.05);
 \node[ch1oncanvas, font=\scriptsize] at (1.15,0) {$+\mathbf{j}_J$};
 \node[ch1oncanvas, font=\scriptsize] at (0,1.18) {$-\mathbf{j}_J$};
 \node[ch1oncanvas, font=\scriptsize] at (0.95,-0.55) {$\Gamma_R$};
 \node[ch1oncanvas, font=\scriptsize, align=center] at (0,-1.15)
 {$\boldsymbol{\Phi}_{J,\mathrm{out}}$};
 \end{scope}
\end{tikzpicture}%
}
\caption{Phase topology on a transverse plane (left) and angular-flux transport on a control surface $\Gamma_R$ (right). Helical isophases show the winding of $\arg E$; angular-momentum transport is established only when direction-resolved $\mathbf{j}_J$ integrates to nonzero $\boldsymbol{\Phi}_{J,\mathrm{out}}$ without cancellation between counter-rotating sectors \cite{allen1992,beth1936}.}
\label{fig:ch1.phase-vs-transport-sketch}
\end{figure}

Allen \emph{et al.} \cite{allen1992} identified paraxial modes; the following is a waist-plane paraxial modal factor with the outgoing axial phase retained for the present discussion, not the full propagated Laguerre--Gaussian solution with $w(z)$, wavefront curvature, and Gouy phase \cite{allen1992}.
\begin{equation}
\label{eq:ch1.phase-winding-number}E_\phi(\rho,\phi,z)
=
E_0
\left(\frac{\rho}{w_0}\right)^{|m|}
L_p^{|m|}\!\left(\frac{2\rho^2}{w_0^2}\right)
e^{-\rho^2/w_0^2}\,
e^{\jj(m\phi-k_0 z)},
\end{equation}
\addtocounter{equation}{-1}\refstepcounter{equation}\label{eq:ch1.lg-vortex-field}%
where \(L_p^{|m|}\) is an associated Laguerre polynomial, \(w_0\) is beam waist,
and integer \(m\) is azimuthal index. For \(p=0\) and \(|m|=1\), the intensity vanishes
on the axis (\(\rho=0\)), while the phase advances by \(2\pi m\) in one circuit about the axis.
Phase topology is the winding structure of the transverse phase; transport is the outward real-power and angular-flux integral on \(\Gamma\), and topology can persist while transport cancels on closed surfaces. The Poynting vector therefore has an appreciable azimuthal component on the intensity ring even when the beam is predominantly collimated. Net structured angular-momentum transport requires a nonzero signed solid-angle integral: superposing the same outgoing modal factor of Eq.~\eqref{eq:ch1.lg-vortex-field} with the opposite winding $m\to -m$ produces the same standing pattern as
Eq.~\eqref{eq:ch1.ex62-counter-rotating-superposition} in
Subsection~\ref{sec:ch01.6.2}. Topology and polarization describe local field structure; transport is fixed only after the signed solid-angle integral of $\mathbf{j}_J$ is formed on a
control surface \cite{allen1992}.

In the paraxial limit of Eq.~\eqref{eq:ch1.lg-vortex-field} with $p=0$ and $|m|=1$,
the leading-order relation $\H^{(0)}\approx(\hat{\mathbf{z}}\times\E^{(0)})/\eta_0$ establishes the axial power flow.
The first paraxial corrections to the Maxwell field generate the transverse, including azimuthal, Poynting component associated with the $\partial_\phi$ phase gradient.
The complex Poynting vector $\mathbf{S}_\mathrm{c}=\tfrac12\E\times\H^*$ then has an axial part
$S_z\propto |E|^2$ that carries time-averaged power along $+\hat{\mathbf{z}}$.
The azimuthal component can appear in $\Re(\mathbf{S}_\mathrm{c})$ as real energy-flow density associated with azimuthal phase winding; locally it may also contribute to $\Im(\mathbf{S}_\mathrm{c})$ where electric and magnetic storage are out of phase \cite{poynting1884}. The axial component integrates to outward power
through a transverse plane. Phase winding therefore produces strong local Poynting
circulation without by itself guaranteeing a nonzero
$\oint \Mavg\cdot\hat{\mathbf{n}}\,dS$ on an enclosing sphere when opposing
solid-angle sectors cancel \cite{allen1992,poynting1884}.

Poynting circulation can be locally strong while the outward real-power integral on a surrounding shell remains small or unchanged: the complex Poynting vector separates reactive exchange through $\Im\{\mathbf{S}_\mathrm{c}\cdot\hat{\mathbf{n}}\}$ from time-averaged radiation through $\Re\{\mathbf{S}_\mathrm{c}\cdot\hat{\mathbf{n}}\}$ \cite{poynting1884}. Poynting's original flow-line pictures \cite{poynting1884} emphasized this
distinction long before Stokes vectors were standard: circulation in $\mathbf{P}=\langle\mathbf{S}\rangle=\Re\{\mathbf{S}_\mathrm{c}\}$ marks directed energy flow and interference texture, while $\oint_{\partial V}\Mavg\cdot\hat{\mathbf{n}}\,dS$ on an enclosing surface expresses outward angular-momentum transport when interpreted on a predominantly radiative shell \cite{poynting1884,beth1936}. Subsection~\ref{sec:ch01.2.3} develops the signed angular-momentum flux that follows from $\Mavg=\mathbf{r}\times\Tavg$ \cite{stratton1941}.

Directional dependence tracks radiation geometry on $\Gamma_R$: the far-zone pattern may exhibit strong angular variation while source and field angular-momentum balances remain distinct objects tied to impressed currents and flux integrals \cite{belinfante1940}.
\subsection{Physical relations for angular-momentum transport}
\label{sec:ch01.2.3}

Azimuthal phase structure and radiation intensity $U_r$ are distinct: net outward angular-momentum transport follows from the Maxwell-stress balance and
\begin{equation}
\label{eq:ch1.am-net-transport-condition}
\boldsymbol{\Phi}_{J,\mathrm{out}}(r)=-\oint_{S_r}\mathbf{r}\times(\Tavg\cdot\hat{\mathbf{n}})\,dS
=-\int_{4\pi}\mathbf{j}_J(\theta,\phi;r)\,d\Omega,
\end{equation}
and the axis-projected scalar transport is
\begin{equation}
\label{eq:ch1.am-net-transport-z}
\Phi_{J,z,\mathrm{out}}(r)=\hat{\mathbf{z}}\cdot\boldsymbol{\Phi}_{J,\mathrm{out}}(r)
=-\int_{4\pi}\hat{\mathbf{z}}\cdot\mathbf{j}_J(\theta,\phi;r)\,d\Omega
=\dot J_{z,\mathrm{em,out}}(r).
\end{equation}
with $\mathbf{j}_J=r^2[\mathbf{r}\times(\Tavg\cdot\hat{\mathbf{r}})]$ from Eq.~\eqref{eq:ch1.directional-am-flux} \cite{stratton1941,beth1936}. Equal counter-rotating components with $P_+=P_-$ give $\dot J_{z,\mathrm{orb,out}}=(m/\omega)(P_+-P_-)=0$ for the declared orbital modal bookkeeping scale while $|E_t|\propto|\cos(m\phi)|$ may remain large (Subsection~\ref{sec:ch01.6.2}); total Maxwell transport $\dot J_{z,\mathrm{em,out}}$ from Eq.~\eqref{eq:ch1.am-net-transport-z} must still be evaluated on the stated field and surface \cite{allen1992,stratton1941}. Outward power in a sector requires $\int_{\Delta\Omega}U_r\,d\Omega>0$; outward angular momentum requires $\boldsymbol{\Phi}_{J,\mathrm{out}}\neq\mathbf{0}$---neither implies the other \cite{poynting1884}.
\phantomsection\label{eq:ch1.flux-qualified-decision}%
\phantomsection\label{eq:ch1.transport-charge-continuous-index}%
\phantomsection\label{tab:ch1.triad-numeric-dipole}

\subsection{Reactive and radiative field regimes}
\label{sec:ch01.2.4}

The field of a finite electric dipole exhibits the classical $1/r^3$, $1/r^2$, and $1/r$ radial hierarchy in its exact exterior solution \cite{stratton1941}. Higher multipoles possess their own higher-order near-field radial structure, while every radiating spherical multipole carries an outgoing $1/r$ branch at infinity \cite{stratton1941,hansen1988}. Near the source the dipole $1/r^3$ and $1/r^2$ terms may be comparable with, or larger than, the $1/r$ radiation field. The near zone can therefore contain substantial electric--magnetic energy exchange and complicated phase structure while the outward real power through an enclosing sphere already equals the total radiated power of the source.

\begin{figure}[t]
\centering
\VolIPhysFig{\VolIPhysRadialFieldZones}
\caption{Radial organization of a localized radiator: $r^{-3}$ electrostatic storage, $r^{-2}$ induction, and $r^{-1}$ radiation dominate at increasing distance. The boundaries are qualitative; all three terms coexist at every radius, but their relative weight shifts with $kr$ \cite{stratton1941,poynting1884}.}
\label{fig:ch1.radial-field-zones}
\end{figure}

Figure~\ref{fig:ch1.radial-field-zones} matches the multipole counting developed in Subsection~\ref{sec:ch01.6.4}: moving outward does not create new radiated power, but changes which radial term carries the local Poynting and stress flux.

At increasing distance the $1/r$ terms dominate. The electric and magnetic fields then approach the transverse relation characteristic of an outgoing wave, and the real Poynting flux becomes the principal carrier of energy away from the source \cite{poynting1884}. Figure~\ref{fig:ch1.wp-reactive-vs-radiative} illustrates this shell dependence: reactive circulation in $\Im\{\mathbf{S}_\mathrm{c}\cdot\hat{\mathbf{n}}\}$ can be large on a small-$kr$ surface even when the net outward real power is unchanged, whereas on a large-$kr$ surface the same source supports a field whose local phase relation is predominantly radiative.

The distinction is evaluated on a spherical observation surface $S_r$ of radius $r$ with outward unit normal $\hat{\mathbf{n}}$. Define
\begin{equation}
\label{eq:ch1.ex64-power-functionals}
P_{\mathrm{rad}}=\int_{S_r}\Re\!\left\{\mathbf{S}_\mathrm{c}\cdot\hat{\mathbf{n}}\right\}dS,\qquad
P_{\mathrm{reac}}(r)=\int_{S_r}\Im\!\left\{\mathbf{S}_\mathrm{c}\cdot\hat{\mathbf{n}}\right\}dS,
\end{equation}
where \(\mathbf{S}_\mathrm{c}=\tfrac{1}{2}\E\times\H^*\) is the complex Poynting vector
\cite{poynting1884}. In a lossless source-free region between enclosing surfaces, $P_{\mathrm{rad}}$ is independent of $r$ \cite{stratton1941}: it is the total time-averaged power radiated by the source. The integral $P_{\mathrm{reac}}(r)$ measures reactive power associated with local oscillatory exchange between electric and magnetic storage on the chosen shell; $|P_{\mathrm{reac}}(r)|$ generally varies with $r$ even when $P_{\mathrm{rad}}$ does not.

The real and imaginary parts of the complex Poynting vector distinguish active power transport from reactive harmonic exchange because they weight different phase relationships between \(\E\) and \(\H\). On a reactive shell, \(\Im\{\mathbf{S}_\mathrm{c}\cdot\hat{\mathbf{n}}\}\) can integrate to a large \(|P_{\mathrm{reac}}(r)|\) while \(P_{\mathrm{rad}}\) remains fixed; on a radiative shell the ratio \(|P_{\mathrm{reac}}(r)|/P_{\mathrm{rad}}\) is typically small even though the same $P_{\mathrm{rad}}$ crosses every enclosing sphere \cite{poynting1884,stratton1941}.

Near the source, with \(kr\ll 1\), the \(r^{-3}\) and \(r^{-2}\) multipole terms dominate and \(|P_{\mathrm{reac}}(r)|/P_{\mathrm{rad}}\) is typically much greater than unity: reactive circulation dominates the normal flux pattern on that shell even though \(P_{\mathrm{rad}}\) is unchanged \cite{stratton1941}. The scale \(kr\sim 1\) marks a transition region in which no single radial power law fixes the shell integrals: the ratio \(|P_{\mathrm{reac}}(r)|/P_{\mathrm{rad}}\) must be evaluated for the actual field on \(S_r\), not inferred from \(kr\) alone \cite{stratton1941}. When \(kr\gg 1\) and the \(r^{-1}\) radiation field dominates, \(|P_{\mathrm{reac}}(r)|/P_{\mathrm{rad}}\ll 1\) typically holds and the field on \(S_r\) is predominantly radiative \cite{poynting1884,stratton1941}. Any angular-momentum interpretation that ignores shell radius therefore risks
attributing storage-dominated structure to propagative torque flux
\cite{stratton1941}. Reactive shells carry standing-wave structure whose phase rotates
locally while time-averaged Poynting flux through the same surface may vanish
\cite{stratton1941}. The same small electric dipole at $kr=0.5$ versus $kr=5$ illustrates why
shell placement matters: the inner shell is storage-dominated with \(|P_{\mathrm{reac}}(r)|/P_{\mathrm{rad}}\gg 1\); the outer shell carries \(|P_{\mathrm{reac}}(r)|/P_{\mathrm{rad}}\ll 1\) and exhibits a predominantly radiative field \cite{harrington2001}. Subsection~\ref{sec:ch01.6.4} develops the two-shell comparison on the same dipole using these integrals.

\section{Rotational covariance and reference frames}
\label{sec:ch01.3}

Flux integrals in section~\ref{sec:ch01.2} are evaluated in a fixed reference frame with fixed origin. If the source is rotated while laboratory axes remain fixed, the field is physically rotated only when the source, medium, and boundaries share the corresponding rotational symmetry \cite{stratton1941,hansen1988}. Relabeling coordinate axes while the impressed current stays fixed changes component charts but not the source. Vectors, the Maxwell stress tensor, and traction then transform as physical objects under rotation; gauge relabeling of potentials is a separate non-uniqueness treated in section~\ref{sec:ch01.4} \cite{belinfante1940,maxwell1873}.

\phantomsection\label{sec:ch01.3.1}

\subsubsection*{Active and passive rotations; transformation of vectors, tensors, and dyads}

\begin{figure}[t]
\centering
\ChOneFigBlock{%
\begin{tikzpicture}[ch1geo]
 \node[ch1oncanvas, font=\scriptsize] at (-1.55,1.45) {(a) source rotation};
 \draw[-{Latex[length=1.8mm,width=1.1mm]}, gray!60!black, line width=0.55pt]
   (-0.15,-1.1) -- (-0.15,1.1);
 \draw[-{Latex[length=1.8mm,width=1.1mm]}, gray!60!black, line width=0.55pt]
   (-1.1,-0.15) -- (1.1,-0.15);
 \node[ch1oncanvas, font=\scriptsize] at (0.05,1.05) {$\hat{\mathbf{z}}$};
 \node[ch1oncanvas, font=\scriptsize] at (1.0,-0.05) {$\hat{\mathbf{x}}$};
 \draw[-{Latex[length=2mm,width=1.3mm]}, gray!40!black, line width=0.65pt]
   (0,0) -- (0.55,0.75);
 \draw[gray!40!black, line width=0.55pt] (0,0) circle (3pt);
 \node[ch1oncanvas, font=\scriptsize] at (0.35,0.45) {$\Jimp$};
 \draw[gray!50!black, dashed, line width=0.5pt] (0,0) circle (0.85);
 \draw[-{Latex[length=1.5mm,width=1mm]}, gray!60!black, line width=0.5pt]
   (0.72,0.35) -- (1.0,0.55);
 \node[ch1oncanvas, font=\scriptsize] at (0.95,0.7) {$\E$};
 \node[ch1oncanvas, font=\scriptsize] at (0,-1.45) {fixed reference axes};

 \node[ch1oncanvas, font=\scriptsize] at (2.45,1.45) {(b) chart rotation};
 \draw[-{Latex[length=1.8mm,width=1.1mm]}, gray!60!black, line width=0.55pt]
   (2.35,-1.1) -- (2.35,1.1);
 \draw[-{Latex[length=1.8mm,width=1.1mm]}, gray!55!black, line width=0.55pt]
   (1.5,0.35) -- (3.3,1.15);
 \node[ch1oncanvas, font=\scriptsize] at (3.05,1.25) {$\hat{\mathbf{x}}'$};
 \draw[-{Latex[length=1.8mm,width=1.1mm]}, gray!55!black, line width=0.55pt]
   (2.05,-0.35) -- (2.85,-0.95);
 \node[ch1oncanvas, font=\scriptsize] at (2.55,-1.05) {$\hat{\mathbf{y}}'$};
 \draw[-{Latex[length=2mm,width=1.3mm]}, gray!40!black, line width=0.65pt]
   (2.35,0) -- (2.9,0);
 \draw[gray!40!black, line width=0.55pt] (2.35,0) circle (3pt);
 \node[ch1oncanvas, font=\scriptsize] at (2.75,0.2) {$\Jimp$};
 \draw[gray!50!black, dashed, line width=0.5pt] (2.35,0) circle (0.85);
 \draw[-{Latex[length=1.5mm,width=1mm]}, gray!60!black, line width=0.5pt]
   (3.07,0) -- (3.35,0.2);
 \node[ch1oncanvas, font=\scriptsize] at (3.45,0.25) {$\E$};
 \node[ch1oncanvas, font=\scriptsize] at (2.35,-1.45) {same source, relabeled axes};
\end{tikzpicture}%
}
\caption{Physical source rotation in a fixed reference frame (a) versus passive relabeling of coordinate axes while the impressed current remains fixed (b). Maxwell radiation changes in (a) when source, medium, and boundary do not share rotational symmetry; passive relabeling in (b) reorders field components without moving the source \cite{hansen1988,stratton1941}.}
\label{fig:ch1.source-vs-chart-rotation}
\end{figure}

Figure~\ref{fig:ch1.source-vs-chart-rotation} separates transformations that are often conflated. In (a) the impressed current rotates while reference axes stay fixed; in (b) the source stays fixed while the component chart is relabeled. A third situation---watching azimuthal phase roll past a fixed observation direction without moving either source or axes---exhibits angular \emph{appearance} in one frame; it is not automatically the same as (a) or (b) \cite{hansen1988}.

For panel (a), when source, medium, and boundary geometry share the same rotational symmetry about the selected axis, rotating the impressed current and then solving Maxwell's equations gives the same physical field as solving first and rotating the result \cite{stratton1941,hansen1988}. An anisotropic or spatially varying medium breaks that equivalence only when its constitutive and boundary data are not invariant under that rotation \cite{stratton1941}.

Passive relabeling in (b) changes only the component chart. Traction $\mathbf{t}=\Tavg\cdot\hat{\mathbf{n}}$ and its moment about a declared origin are built from $\Tavg$ (Eq.~\eqref{eq:ch1.maxwell-stress}). Vectors rotate with one application of $R$; $\Tavg$ rotates as $\Tavg'=R\Tavg R^\top$ \cite{beth1936,stratton1941}.

For proper rotation \(R\in\mathrm{SO}(3)\), with \(R^\top R=I\), \(\det R=1\), let
\(\mathbf{v}=v_j\hat{\mathbf{e}}_j\) be a vector in orthonormal basis
\(\hat{\mathbf{e}}_j\). With the active convention used here, a physical rotation sends $\mathbf{v}'_{\mathrm{active}}=R\mathbf{v}$. For a passive rotation of the basis by the same $R$, the components transform as $\mathbf{v}'_{\mathrm{passive}}=R^\top\mathbf{v}$. The two descriptions represent the same geometric vector when the active and passive rotations are taken as inverse operations. Eq.~\eqref{eq:ch1.vector-rotation-components} states the active component law
\begin{equation}
\label{eq:ch1.infinitesimal-active-rotation}v_i' = R_{ij}v_j.
\end{equation}
\addtocounter{equation}{-1}\refstepcounter{equation}\label{eq:ch1.vector-rotation-components}%
\addtocounter{equation}{-1}\refstepcounter{equation}\label{eq:ch1.active-rotation-action}%
For a finite rotation through angle \(\varphi\) about unit axis \(\hat{\mathbf{u}}\), the active Rodrigues formula reads
\begin{equation}
\label{eq:ch1.rodrigues-rotation}
R(\varphi,\hat{\mathbf{u}})\,\mathbf{v}
=
\mathbf{v}\cos\varphi
+(\hat{\mathbf{u}}\times\mathbf{v})\sin\varphi
+\hat{\mathbf{u}}(\hat{\mathbf{u}}\cdot\mathbf{v})(1-\cos\varphi),
\end{equation}
equivalent on vector fields to the generator exponential developed below and used in Chapter~\ref{ch:03} \cite{jackson1999,stratton1941}.

Rotating a field is not the same operation as rotating its component chart. For a proper rotation $R\in\mathrm{SO}(3)$, the active action on a scalar field $F$ and on a vector field $\mathbf{F}$ is
\begin{equation}
\label{eq:ch1.field-rotation-operator}
(T_R F)(\mathbf{r})=F(R^{-1}\mathbf{r}),
\qquad
(T_R\mathbf{F})(\mathbf{r})=R\,\mathbf{F}(R^{-1}\mathbf{r}),
\end{equation}
so the spatial argument and the vector direction transform together \cite{jackson1999}. Evaluating $\E$ at $\mathbf{r}$ after an active rotation is not equivalent to multiplying $\E(\mathbf{r})$ by $R$ at fixed $\mathbf{r}$ unless the configuration is merely relabeled \cite{stratton1941}. Maxwell's equations in source-free vacuum are covariant under $T_R$ because $\curl$ and $\divg$ commute with the spatial part of the rotation action when $\varepsilon$ and $\mu$ are scalars \cite{maxwell1873}; impressed sources break source--field equivalence under rotation, and an anisotropic medium does so only when its constitutive and boundary data fail to share the selected rotation symmetry, even though the transformation law itself remains valid \cite{stratton1941}.

After an active source rotation, the same $|E|$ on $S_R$ is expressed in rotated $\hat{\mathbf{x}},\hat{\mathbf{y}}$ components \cite{hansen1988}. Rotating the \emph{source} (Subsection~\ref{sec:ch01.0.2}) changes $\Jimp$ and alters radiation when the source, medium, and boundary geometry do not share the required rotational symmetry; passively rotating the \emph{component chart} while holding the source fixed relabels components without moving the excitation \cite{stratton1941}.

The Maxwell stress tensor links $\mathbf{T}$, $\hat{\mathbf{n}}$, and $\mathbf{t}$: each column of $\mathbf{T}$ specifies the traction that would act on a surface element whose normal points along a coordinate direction. Because $\mathbf{T}$ is a geometric object that maps normals to force densities, it must rotate with \emph{two} applications of $R$---one for each index:
\begin{equation}
\label{eq:ch1.passive-rotation-action}\Tavg' = R\,\Tavg\,R^\top,
\qquad
T'_{ij}=R_{ip}R_{jq}T_{pq}.
\end{equation}
\addtocounter{equation}{-1}\refstepcounter{equation}\label{eq:ch1.tensor-rotation-rule}%
\addtocounter{equation}{-1}\refstepcounter{equation}\label{eq:ch1.vector-observable-covariance}%
The two factors \(R_{ip}\), \(R_{jq}\) rotate each tensor slot. Eq.~\eqref{eq:ch1.tensor-rotation-rule} ensures that stress tensor
Eq.~\eqref{eq:ch1.maxwell-stress} rotates as a single geometric object acting on surface normals, not as nine independent scalars.
Mishandling this law---treating $T_{ij}$ as entries that rotate independently---is a frequent source of spurious torque when stress-integral results are compared with mechanical reaction on a polarizable body \cite{beth1936}.

When the reference frame rotates, the surface element rotates with it: traction on the rotated patch is still $\mathbf{t}=\Tavg\cdot\hat{\mathbf{n}}$, and
\begin{equation}
\label{eq:ch1.traction-rotation-covariance}
\mathbf{t}'=\Tavg'\cdot\hat{\mathbf{n}}' = R\mathbf{t}.
\end{equation}
Eq.~\eqref{eq:ch1.traction-rotation-covariance} states the mechanical consequence: the force per unit area on the rotated surface element is the same physical force density, expressed in the rotated reference basis \cite{beth1936}. Vacuum Maxwell stress is symmetric, \(\Tavg=\Tavg^\top\), under \(R\in\mathrm{SO}(3)\) \cite{maxwell1873}.

For an axially symmetric field about \(\hat{\mathbf{z}}\), azimuthal dependence separates as \(\psi(\phi)=e^{\jj m\phi}\), and the orbital angular index satisfies
\begin{equation}
\label{eq:ch1.azimuthal-eigenvalue}
-\jj\,\frac{\partial\psi}{\partial\phi}=m\,\psi.
\end{equation}
Eq.~\eqref{eq:ch1.azimuthal-eigenvalue} is the scalar harmonic index of azimuthal phase; it is not by itself the full vector-field angular-momentum operator on \((\E,\H)\) \cite{stratton1941}. Comparisons of torque on a polarizable body require a fixed convention: whether $R$ acted on the source, on the component chart, or on both, together with a fixed reference origin \cite{beth1936,stratton1941}.

\phantomsection\label{eq:ch1.active-passive-duality}
\phantomsection\label{eq:ch1.tensor-invariant-trace}
\phantomsection\label{eq:ch1.stress-eigen-covariance}
\phantomsection\label{eq:ch1.rotation-group-action}
\phantomsection\label{eq:ch1.dyadic-covariance}
\phantomsection\label{eq:ch1.rank1-dyad-rotation}
\phantomsection\label{sec:ch01.3.2b}
\phantomsection\label{eq:ch1.total-angular-generator}
\phantomsection\label{eq:ch1.orbital-angular-generator}
\phantomsection\label{eq:ch1.spin-angular-generator}
\phantomsection\label{eq:ch1.rotation-from-generators}

\phantomsection\label{sec:ch01.3.3}
\subsubsection*{Origin dependence}

Torque is always torque \emph{about a point}. Translating the reference origin from $\mathbf{r}_0$ to $\mathbf{r}_0'=\mathbf{r}_0+\Delta\mathbf{r}$ leaves the resultant force unchanged but adds the couple $\boldsymbol{\tau}_{\mathbf{r}_0'}=\boldsymbol{\tau}_{\mathbf{r}_0}-\Delta\mathbf{r}\times\mathbf{F}$, with $\mathbf{F}=\oint_{S_m}\Tavg\cdot\hat{\mathbf{n}}\,dS$, even when $(\E,\H)$ and $\mathbf{T}$ on $S_m$ are unchanged \cite{beth1936,stratton1941}. Figure~\ref{fig:ch1.origin-shift-torque} shows that force is a property of the traction distribution, whereas torque depends on the chosen origin. Angular-momentum flux $\Mavg=\mathbf{r}\times\Tavg$ carries the same origin dependence \cite{stratton1941}.

\begin{figure}[t]
\centering
\ChOneFigBlock{%
\begin{tikzpicture}[ch1geo]
 \coordinate (O) at (0,0);
 \coordinate (Op) at (0.55,0.35);
 \coordinate (P) at (1.35,0.55);
 \fill (O) circle (2pt);
 \fill (Op) circle (2pt);
 \fill (P) circle (2pt);
 \node[ch1oncanvas, font=\scriptsize, anchor=east] at (O) {$O$};
 \node[ch1oncanvas, font=\scriptsize, anchor=west] at (Op) {$O'$};
 \node[ch1oncanvas, font=\scriptsize, anchor=west] at (P) {$P$};
 \draw[gray!55!black, line width=0.45pt] (-0.35,-0.55) -- (1.65,-0.55);
 \node[ch1oncanvas, font=\scriptsize] at (0.65,-0.75) {material surface $S_m$};
 \draw[-{Latex[length=2mm,width=1.3mm]}, gray!50!black, line width=0.6pt]
   (P) -- (1.75,0.55);
 \node[ch1oncanvas, font=\scriptsize] at (1.95,0.55) {$\mathbf{t}$};
 \draw[-{Latex[length=1.8mm,width=1.1mm]}, gray!60!black, line width=0.55pt]
   (O) -- (P);
 \node[ch1oncanvas, font=\scriptsize] at (0.75,0.15) {$\mathbf{r}$};
 \draw[-{Latex[length=1.8mm,width=1.1mm]}, gray!55!black, line width=0.55pt]
   (Op) -- (P);
 \node[ch1oncanvas, font=\scriptsize] at (1.05,0.05) {$\mathbf{r}'$};
 \draw[-{Latex[length=1.6mm,width=1mm]}, gray!40!black, line width=0.55pt, dashed]
   (O) -- (Op);
 \node[ch1oncanvas, font=\scriptsize] at (0.15,0.25) {$\Delta\mathbf{r}$};
 \draw[-{Latex[length=1.8mm,width=1.1mm]}, gray!45!black, line width=0.55pt]
   (O) -- (-0.15,0.55);
 \node[ch1oncanvas, font=\scriptsize] at (-0.35,0.65) {$\boldsymbol{\tau}_O$};
 \draw[-{Latex[length=1.8mm,width=1.1mm]}, gray!45!black, line width=0.55pt]
   (Op) -- (0.95,0.95);
 \node[ch1oncanvas, font=\scriptsize] at (1.15,1.05) {$\boldsymbol{\tau}_{O'}$};
\end{tikzpicture}%
}
\caption{Force density $\mathbf{t}$ at point $P$ on $S_m$ contributes to resultant force independent of origin, but torque $\boldsymbol{\tau}_O$ about $O$ differs from $\boldsymbol{\tau}_{O'}$ about $O'=O+\Delta\mathbf{r}$ because the moment arm changes \cite{beth1936}.}
\label{fig:ch1.origin-shift-torque}
\end{figure}

\phantomsection\label{sec:ch01.3.4}

\FloatBarrier
\section{Gauge structure, helicity, and spin--orbit limits}
\label{sec:ch01.4}

\phantomsection\label{sec:ch01.4.2}
\phantomsection\label{prin:ch1.angular-em-structure}

\subsubsection*{Gauge freedom and angular structure}

Rotation and origin declaration fixed which mechanical reports built from $(\E,\H)$ are objective. Potentials introduce a further question that rotation alone does not settle: can $(\mathbf{A},\phi)$ be changed without changing the physical field? On a compact range two field maps may disagree on $\mathbf{A}(\mathbf{r})$ at the same sample point yet reproduce identical tangential $\E$ and $\H$ on the scan surface \cite{stratton1941,hansen1988}. That non-uniqueness is not a failure of Maxwell theory; it is a property of the potential representation. Wattmeters, torsion balances, and stress integrals respond to $(\E,\H)$, $\mathbf{S}_\mathrm{c}$, and $\Tavg$; they do not read $\mathbf{A}$ directly \cite{maxwell1873,beth1936}.

\begin{figure}[t]
\centering
\ChOneFigBlock{%
\begin{tikzpicture}[ch1geo]
 \node[ch1oncanvas, font=\scriptsize] at (-1.35,1.35) {(a) $\mathbf{A}_1$};
 \draw[-{Latex[length=1.8mm,width=1.1mm]}, gray!60!black, line width=0.55pt]
   (-0.15,-0.95) -- (-0.15,0.95);
 \draw[-{Latex[length=1.8mm,width=1.1mm]}, gray!60!black, line width=0.55pt]
   (-0.95,-0.15) -- (0.95,-0.15);
 \draw[-{Latex[length=2mm,width=1.3mm]}, gray!60!black, line width=0.65pt]
   (-0.55,0.35) arc (145:35:0.65);
 \draw[-{Latex[length=2mm,width=1.3mm]}, gray!40!black, line width=0.65pt]
   (0.15,0.15) -- (0.75,0.75);
 \node[ch1oncanvas, font=\scriptsize] at (0.55,0.85) {$\E$};
 \node[ch1oncanvas, font=\scriptsize] at (-0.55,0.55) {$\mathbf{A}_1$};

 \node[ch1oncanvas, font=\scriptsize] at (1.65,1.35) {(b) $\mathbf{A}_2=\mathbf{A}_1+\nabla\chi$};
 \draw[-{Latex[length=1.8mm,width=1.1mm]}, gray!60!black, line width=0.55pt]
   (1.35,-0.95) -- (1.35,0.95);
 \draw[-{Latex[length=1.8mm,width=1.1mm]}, gray!60!black, line width=0.55pt]
   (0.55,-0.15) -- (2.45,-0.15);
 \draw[-{Latex[length=2mm,width=1.3mm]}, gray!60!black, line width=0.65pt, dashed]
   (1.05,0.55) .. controls (1.55,0.15) .. (2.05,0.45);
 \draw[-{Latex[length=2mm,width=1.3mm]}, gray!40!black, line width=0.65pt]
   (1.65,0.15) -- (2.25,0.75);
 \node[ch1oncanvas, font=\scriptsize] at (2.05,0.85) {$\E$};
 \node[ch1oncanvas, font=\scriptsize] at (1.15,0.55) {$\mathbf{A}_2$};

 \node[ch1oncanvas, font=\scriptsize, align=center] at (0.55,-1.25)
 {same $\E$, $\H$; different $\mathbf{A}$, $\phi$};
\end{tikzpicture}%
}
\caption{Two potential pairs can describe the same physical field: $\mathbf{A}_2=\mathbf{A}_1+\nabla\chi$ relabels the vector potential while $\E=-\nabla\phi-\jj\omega\mathbf{A}$ and $\H=\curl\mathbf{A}/\mu$ remain unchanged \cite{maxwell1873}.}
\label{fig:ch1.gauge-same-field}
\end{figure}

The curl of $\mathbf{A}$ fixes $\mathbf{B}=\mu\H$; adding $\nabla\chi$ to $\mathbf{A}$ and adjusting $\phi$ leaves $(\E,\H)$ unchanged (Figure~\ref{fig:ch1.gauge-same-field}). Quantities built from $(\E,\H)$ alone---including closed-surface flux and stress integrals---are gauge invariant; pointwise maps that invoke $\phi$ or $\mathbf{A}$ are not \cite{stratton1941,belinfante1940}.
\subsection{Gauge-invariant observables, helicity, and dual symmetry}
\label{sec:ch01.4.3}

Let the phasor field state be \(\mathcal{F}=(\E,\H)\) in a linear isotropic region
with scalar permittivity \(\varepsilon\), scalar permeability \(\mu\), angular
frequency \(\omega\), and intrinsic impedance
\(\eta=\sqrt{\mu/\varepsilon}\). Let \((\mathbf{A},\phi)\) denote a potential pair
related to the physical fields by
\begin{equation}
\label{eq:ch1.potential-field-relations}
\E=-\nabla\phi-\jj\omega\mathbf{A},\qquad \mathbf{B}=\curl\mathbf{A}=\mu\H,
\end{equation}
The vector potential $\mathbf{A}$ is not a field observable; it is a convenience for solving Maxwell's equations whose curl fixes $\mathbf{B}$ and whose combination with $\phi$ fixes $\E$. Because $\curl\nabla\chi=\mathbf{0}$, adding $\nabla\chi$ to $\mathbf{A}$ changes $\mathbf{A}$ but not $\mathbf{B}$; adjusting $\phi$ by $-\jj\omega\chi$ restores $\E$:
\begin{equation}
\label{eq:ch1.gauge-transform}
\mathbf{A}'=\mathbf{A}+\nabla\chi,\qquad \phi'=\phi-\jj\omega\chi,\qquad \E'=\E,\ \H'=\H.
\end{equation}
Substituting into $\E=-\nabla\phi-\jj\omega\mathbf{A}$ gives
$\E'=-\nabla(\phi-\jj\omega\chi)-\jj\omega(\mathbf{A}+\nabla\chi)
=-\nabla\phi-\jj\omega\mathbf{A}=\E$,
with $\mathbf{B}'=\curl\mathbf{A}'=\mathbf{B}$ unchanged. Eq.~\eqref{eq:ch1.gauge-transform} is the relabeling shown in Figure~\ref{fig:ch1.gauge-same-field}. What remains observable are $\E$, $\B$, $\mathbf{S}_\mathrm{c}$, $\Tavg$, and closed-surface integrals built from $(\E,\H)$ without explicit potential dependence \cite{maxwell1873,belinfante1940}. \phantomsection\label{eq:ch1.gauge-observable-invariance}\phantomsection\label{eq:ch1.observable-admissibility-triad}\phantomsection\label{eq:ch1.observable-charge-continuous-triad}
Potential-tagged local maps require an explicit gauge declaration \cite{belinfante1940}.

Coulomb and Lorenz gauges yield the same \((\E,\H)\) yet different \(\phi\) and \(\nabla\phi\) at interior points, so a potential-based local map need not agree pointwise between them even when traction and angular-flux integrals built from $\Tavg$ and $\Mavg=\mathbf{r}\times\Tavg$ remain stable on declared surfaces \cite{belinfante1940,maxwell1873}. For illustration, define the following gauge-dependent potential map. It is introduced only to demonstrate the non-uniqueness of potential-based local constructions; it is not an electromagnetic angular-momentum density.

\begin{equation}
\label{eq:ch1.representation-linked-jloc}
\mathbf{q}_{\mathrm{loc}}^{(\mathrm{pot})}
=
\frac{\varepsilon_0}{\jj\omega}\,\mathbf{r}\times\nabla\phi,
\end{equation}
where \(\phi\) is the scalar potential in the chosen gauge of Eq.~\eqref{eq:ch1.potential-field-relations}. Under \(\phi'=\phi-\jj\omega\chi\), \(\nabla\phi'=\nabla\phi-\jj\omega\nabla\chi\), so \(\mathbf{q}_{\mathrm{loc}}^{(\mathrm{pot})}\) shifts by \(-\varepsilon_0\,\mathbf{r}\times\nabla\chi\). Eq.~\eqref{eq:ch1.representation-linked-jloc} is therefore a gauge-linked orbital map associated with \(\phi\), not the gauge-invariant electromagnetic angular-momentum density \(\mathbf{j}_{\mathrm{em}}\) and not a Belinfante-class local spin density \cite{belinfante1940,stratton1941}.

Within a control volume $V_g$ of radius $R=2\lambda$ about a dipole, Coulomb and Lorenz gauges yield the same \((\E,\H)\) on \(\partial V_g\) yet different \(\nabla\phi\) and therefore different Eq.~\eqref{eq:ch1.representation-linked-jloc} at interior points; \(\oint_{\partial V_g}\Mavg\cdot\hat{\mathbf{n}}\,dS\) is unchanged \cite{belinfante1940,stratton1941}. Helicity measures the handedness with which $\E$ and $\H$ rotate about the propagation direction. For a monochromatic circularly polarized plane wave, $\E$ and $\H$ trace the same sense of rotation while remaining orthogonal to $\mathbf{k}$; right-handed and left-handed states differ by that sense, not by a change in radiated power \cite{stratton1941}. Figure~\ref{fig:ch1.helicity-cp-wave} shows the geometry before the density is written.

\begin{figure}[t]
\centering
\ChOneFigBlock{%
\begin{tikzpicture}[ch1geo]
 \draw[-{Latex[length=2mm,width=1.3mm]}, gray!60!black, line width=0.65pt]
   (-1.1,0) -- (1.1,0);
 \node[ch1oncanvas, font=\scriptsize] at (1.25,0) {$\mathbf{k}$};
 \draw[-{Latex[length=1.8mm,width=1.1mm]}, gray!35!black, line width=0.6pt]
   (0,0.75) -- (0,1.05);
 \draw[-{Latex[length=1.8mm,width=1.1mm]}, gray!55!black, line width=0.6pt]
   (0.75,0) -- (1.05,0);
 \draw[gray!35!black, line width=0.55pt] (0,0) circle (0.75);
 \draw[gray!55!black, line width=0.55pt] (0,0) circle (0.55);
 \node[ch1oncanvas, font=\scriptsize] at (-0.95,0.95) {$\E$};
 \node[ch1oncanvas, font=\scriptsize] at (0.95,0.55) {$\H$};
 \node[ch1oncanvas, font=\scriptsize, align=center] at (0,-1.05)
 {same rotation sense $\Rightarrow$ definite helicity};
\end{tikzpicture}%
}
\caption{Circularly polarized plane wave: $\E$ and $\H$ rotate in the same handedness about propagation axis $\mathbf{k}$. Helicity is the field-intrinsic sign of that handedness \cite{stratton1941,belinfante1940}.}
\label{fig:ch1.helicity-cp-wave}
\end{figure}

The monochromatic vacuum electromagnetic helicity density used in this chapter is
\begin{equation}
\label{eq:ch1.helicity-density}
h(\mathbf{r})=
\frac{1}{2\omega c}\,
\Im\!\left\{\E\cdot\H^*\right\},
\end{equation}
with $\exp(+\jj\omega t)$ phasors \cite{stratton1941,belinfante1940,tai1994}.
Dimensionally, $[\E\cdot\H^*]=\mathrm{W\,m^{-2}}$ and
$[h]=\mathrm{J\,s\,m^{-3}}$---angular momentum per unit volume---because $\omega$ carries $\mathrm{rad\,s^{-1}}$ and $c$ carries $\mathrm{m\,s^{-1}}$.
Eq.~\eqref{eq:ch1.helicity-density} must not be confused with optical chirality (Lipkin
zilch) or with Stokes parameter $S_3$: those are distinct pseudoscalars with different
normalizations and conservation properties \cite{stratton1941,tai1994}. Helicity
$h$ measures handedness of the transverse $\E$--$\H$ triad; local spin--orbit partitions are fixed only under an explicit decomposition convention
(section~\ref{sec:ch01.4.5}) \cite{belinfante1940,stratton1941}.

For an outgoing circularly polarized plane wave with \emph{real} amplitude $E_0>0$,
\(\E=(E_0/\sqrt{2})(\hat{\mathbf{x}}+\jj\sigma\hat{\mathbf{y}})\,e^{-\jj kz}\) with
\(\sigma=\pm1\), \(\mathbf{k}=k\hat{\mathbf{z}}\), and right-handed coordinates
\((\hat{\mathbf{x}},\hat{\mathbf{y}},\hat{\mathbf{z}})\), Faraday's law gives
\(\H=(1/\eta_0)\,\hat{\mathbf{z}}\times\E
=(E_0/\eta_0\sqrt{2})(\hat{\mathbf{y}}-\jj\sigma\hat{\mathbf{x}})\,e^{-\jj kz}\).
Then \(\E\cdot\H^*=\jj\sigma E_0^2/\eta_0\) and
\begin{equation}
\label{eq:ch1.helicity-plane-wave}
h=\frac{\sigma E_0^2}{2\omega c\eta_0}
=\frac{\sigma\varepsilon_0 E_0^2}{2\omega},
\end{equation}
with opposite signs for opposite circular polarizations and $h=0$ for linear polarization
\cite{stratton1941,belinfante1940}. Handedness labels ($\sigma=+1$ vs.\ $\sigma=-1$) are fixed
only after the spatial triad and time convention are stated; they are not interchangeable
with Stokes $S_3$ without an explicit sign map \cite{stratton1941}.

The dual rotation on field pairs is
\begin{equation}
\label{eq:ch1.duality-rotation}
\E'=\cos\xi\,\E+\sin\xi\,\eta\H,\qquad
\H'=-\sin\xi\,\E/\eta+\cos\xi\,\H,\qquad \xi\in\mathbb{R}.
\end{equation}
Eq.~\eqref{eq:ch1.duality-rotation} is not a gauge transformation: it
mixes \(\E\) and \(\H\) while leaving the source-free Maxwell system invariant in
isotropic media. Infinitesimal dual variation \(\xi=\delta\xi\) gives
\(\delta\E=\delta\xi\,\eta\H\) and \(\delta\H=-\delta\xi\,\E/\eta\) to first order.
Dual symmetry is a property of the field equations in homogeneous regions. Under a continuous
duality rotation Eq.~\eqref{eq:ch1.duality-rotation}, a helicity eigenfield retains its
helicity sign and acquires only a phase within the same transverse sector; a rotation by
\(\xi=\pi/2\) exchanges electric and magnetic roles but does not map a \(+1\) helicity
eigenstate onto a \(-1\) eigenstate \cite{stratton1941}. Helicity claims require a fixed
polarization, outgoing selection, and a field pair in which \(\E\) and \(\H\) remain
phase-coherent. Duality is an exact symmetry of the source-free Maxwell pair in isotropic, lossless regions where the impedance \(\eta=\sqrt{\mu/\varepsilon}\) is constant (including proportional profiles \(\varepsilon(\mathbf r)\propto\mu(\mathbf r)\) with fixed \(\eta\)); abrupt changes in \(\eta\) at interfaces, conduction loss, impressed sources, or broken \(\E\)--\(\H\) phase coherence break the global rotation \cite{stratton1941,fernandezcorbaton2013}.

Because Eq.~\eqref{eq:ch1.helicity-density} is built from $(\E,\H)$, $h$ is gauge invariant; potential-based canonical helicity definitions may shift under $\mathbf{A}'=\mathbf{A}+\nabla\chi$ and are not interchangeable with Eq.~\eqref{eq:ch1.helicity-density} without stating the formalism \cite{belinfante1940,tai1994,jackson1999}. For a circularly polarized plane wave, $\langle\mathbf{S}\rangle$ has the same magnitude for $\sigma=\pm1$ while $h$ changes sign with handedness; a linearly polarized wave gives $h=0$ \cite{stratton1941}. Chapter~\ref{ch:01} uses helicity as a field-based handedness scalar; a separate helicity continuity law is not developed here because energy, momentum, and angular momentum already carry the transport structure needed for radiation \cite{stratton1941}.

Helicity and dual symmetry close the gauge discussion of section~\ref{sec:ch01.4}. Maxwell's equations determine $(\E,\H)$ once sources, constitutive data, boundary conditions, and outgoing selection are specified \cite{maxwell1873,stratton1941}. On a declared surface, cycle-averaged power flux follows from $\langle\mathbf{S}\rangle=\Re\mathbf{S}_\mathrm{c}$, mechanical load from $\Tavg\cdot\hat{\mathbf{n}}$, and angular action from $\Mavg=\mathbf{r}\times\Tavg$ \cite{poynting1884,beth1936}.

\phantomsection\label{sec:ch01.4.4}
\subsection{Limits of local spin-orbit decompositions}
\label{sec:ch01.4.5}

Total electromagnetic angular momentum is a conserved quantity tied to $(\E,\H)$ through Eq.~\eqref{eq:ch1.angular-momentum-balance}. Helicity, developed above, is a gauge-invariant sign of handedness in the field pair. Local ``spin'' and ``orbit'' labels partition $\mathbf{j}_{\mathrm{em}}$ into pointwise contributions useful for interpretation but not uniquely fixed by Maxwell's equations alone \cite{belinfante1940,stratton1941}. The total density $\mathbf{j}_{\mathrm{em}}$ is fixed by $(\E,\H)$; convention-dependent $\mathbf{s}_{\mathrm{loc}}$ and $\mathbf{l}_{\mathrm{loc}}$ may redistribute under gauge choice without changing closed-surface flux \cite{belinfante1940,stratton1941}. Measurable transport remains $\oint \Mavg\cdot\hat{\mathbf{n}}\,dS$ and the balance of Eq.~\eqref{eq:ch1.angular-momentum-balance-averaged}; paraxial index $m$ is developed in Chapter~\ref{ch:02} \cite{allen1992}.

\begin{figure}[!htbp]
\centering
\adjustbox{max width=0.78\linewidth, max height=0.42\textheight, center}{%
\begin{tikzpicture}[ch1geo]
 \draw[-{Latex[length=2mm,width=1.3mm]}, gray!60!black, line width=0.65pt]
   (-1.0,0) -- (1.0,0);
 \node[ch1oncanvas, font=\scriptsize] at (1.15,0) {$\mathbf{k}$};
 \draw[gray!55!black, line width=0.55pt] (0,0) circle (0.55);
 \node[ch1oncanvas, font=\scriptsize] at (0,0.75) {helicity};
 \draw[gray!40!black, dashed, line width=0.5pt]
   (-0.45,-0.35) arc (-140:40:0.55);
 \node[ch1oncanvas, font=\scriptsize] at (-0.55,-0.55) {phase winding};
 \draw[-{Latex[length=2mm,width=1.3mm]}, gray!30!black, line width=0.65pt]
   (0,0) -- (0,0.95);
 \node[ch1oncanvas, font=\scriptsize] at (0.25,1.05) {$\mathbf{j}_{\mathrm{em}}$};
 \draw[-{Latex[length=1.6mm,width=1mm]}, gray!45!black, line width=0.55pt]
   (0.35,0.15) -- (0.75,0.35);
 \node[ch1oncanvas, font=\scriptsize] at (0.95,0.45) {$\mathbf{s}_{\mathrm{loc}}$};
 \draw[-{Latex[length=1.6mm,width=1mm]}, gray!65!black, line width=0.55pt]
   (-0.35,0.15) -- (-0.75,0.35);
 \node[ch1oncanvas, font=\scriptsize] at (-1.05,0.45) {$\mathbf{l}_{\mathrm{loc}}$};
 \node[ch1oncanvas, font=\scriptsize, align=center] at (0,-1.05)
 {chosen split of one total};
\end{tikzpicture}%
}
\caption{Total angular-momentum density $\mathbf{j}_{\mathrm{em}}$ may be partitioned into spin-like $\mathbf{s}_{\mathrm{loc}}$ and orbital-like $\mathbf{l}_{\mathrm{loc}}$ contributions for a structured beam; helicity (polarization handedness) and azimuthal phase winding are distinct physical features, and the local split is a convention, not two separate substances \cite{belinfante1940,allen1992}.}
\label{fig:ch1.spin-orbit-decomposition}
\end{figure}

\FloatBarrier
\section{Surface flux, torque, and physical observation}
\label{sec:ch01.5}

On a declared surface $\Gamma$, cycle-averaged power follows from $\langle\mathbf{S}\rangle=\Re\mathbf{S}_\mathrm{c}$, and mechanical force and torque from $\Tavg$ and $\Mavg$ \cite{poynting1884,maxwell1873,beth1936}. Harmonic coefficients of $\E_t$ organize the same field for synthesis in Chapter~\ref{ch:02}; they are not substitutes for flux integrals \cite{stratton1941,hansen1988}.

\begin{figure}[t]
\centering
\ChOneFigBlock{%
\begin{tikzpicture}[ch1geo]
 \draw[gray!50!black, dashed, line width=0.45pt] (0,0) circle (0.95);
 \node[ch1oncanvas, font=\scriptsize] at (0,1.15) {$\Gamma$};
 \fill[gray!40] (0,0) circle (0.12);
 \node[ch1oncanvas, font=\scriptsize] at (0,-0.28) {source};
 \draw[-{Latex[length=1.8mm,width=1.1mm]}, gray!45!black, line width=0.55pt]
   (0.55,0.35) -- (0.95,0.55);
 \draw[-{Latex[length=1.8mm,width=1.1mm]}, gray!45!black, line width=0.55pt]
   (-0.35,0.65) -- (-0.55,0.95);
 \node[ch1oncanvas, font=\scriptsize] at (1.05,0.65) {$\langle\mathbf{S}\rangle$};
 \draw[gray!55!black, line width=0.5pt] (-0.95,0) -- (-0.55,0);
 \draw[-{Latex[length=1.6mm,width=1mm]}, gray!55!black, line width=0.55pt]
   (-0.75,0) -- (-0.75,0.45);
 \draw[-{Latex[length=1.6mm,width=1mm]}, gray!55!black, line width=0.55pt]
   (-0.75,0) -- (-0.95,0.18);
 \node[ch1oncanvas, font=\scriptsize, align=center] at (-1.05,0.55)
 {$\mathbf{t}$};
 \draw[-{Latex[length=1.4mm,width=0.9mm]}, gray!40!black, line width=0.5pt]
   (-0.75,0) -- (-0.95,-0.35);
 \node[ch1oncanvas, font=\scriptsize] at (-1.05,-0.45) {$\mathbf{r}_0$};
 \draw[-{Latex[length=1.6mm,width=1mm]}, gray!65!black, line width=0.5pt]
   (0.65,-0.55) arc [start angle=-40, end angle=20, radius=0.75];
 \node[ch1oncanvas, font=\scriptsize] at (0.95,-0.35) {$\E_t$};
 \node[ch1oncanvas, font=\scriptsize, align=center] at (0,-1.15)
 {flux, traction, and tangential field on one shell};
\end{tikzpicture}%
}
\caption{Maxwell field content on a closed surface $\Gamma$: outward cycle-averaged power flux $\langle\mathbf{S}\rangle$ enters power balance; stress traction $\mathbf{t}=\Tavg\cdot\hat{\mathbf{n}}$ on a material patch enters force and torque about origin $\mathbf{r}_0$; tangential $\E_t$ is part of the same field pair \cite{poynting1884,beth1936,stratton1941}.}
\label{fig:ch1.measurement-surface-readout}
\end{figure}

Field--matter coupling is where Maxwell predictions acquire mechanical consequence: traction $\mathbf{t}=\Tavg\cdot\hat{\mathbf{n}}$ acts across conductor, dielectric, and polarizable interfaces, and the resulting force and torque express how absorption or polarization conversion redistributes angular action \cite{beth1936}. Beth's quartz wave-plate experiment under circularly polarized light is the canonical instance: the plate face is the surface $S_m$ for traction integration, and the observed twist follows from electromagnetic stress at that interface---not from an orbital mode index assigned in a harmonic chart \cite{beth1936,allen1992}.

\phantomsection\label{sec:ch01.5.1}

Cycle-averaged Maxwell stress $\Tavg$ is the local flux of linear momentum \cite{maxwell1873,stratton1941}. At a boundary between vacuum and matter, the traction
\begin{equation}
\label{eq:ch1.stress-traction-definition}
\mathbf{t}=\Tavg\cdot\hat{\mathbf{n}}
\end{equation}
is the force per unit area exerted by the field on the material surface element with outward normal $\hat{\mathbf{n}}$. Let $\partial V_b$ be the \emph{closed} surface bounding the material body (all faces included). The net electromagnetic force on the body is
\begin{equation}
\label{eq:ch1.stress-force-integral}
\mathbf{F}=\oint_{\partial V_b}\Tavg\cdot\hat{\mathbf{n}}\,dS.
\end{equation}
A finite material patch $S_m\subset\partial V_b$ contributes only part of this integral unless the remaining boundary segments and any edge terms are accounted for \cite{stratton1941}. Angular momentum is the moment of linear momentum about a chosen origin $\mathbf{r}_0$. The electromagnetic torque resultant is
\begin{equation}
\label{eq:ch1.stress-torque-integral}
\mathbf{T}_{\mathbf{r}_0}=\oint_{\partial V_b}(\mathbf{r}-\mathbf{r}_0)\times(\Tavg\cdot\hat{\mathbf{n}})\,dS.
\end{equation}
Here $\mathbf{r}$ is measured from the same origin as in $\Mavg=\mathbf{r}\times\Tavg$; torque is the moment of the same momentum flux carried by $\Tavg$ \cite{stratton1941,beth1936}. Because angular momentum is defined with respect to an origin, $\boldsymbol{\tau}_{\mathbf{r}_0}$ changes when $\mathbf{r}_0$ is translated even if $(\E,\H)$ and $\Tavg$ on $S_m$ are unchanged (Eq.~\eqref{eq:ch1.stress-origin-shift-torque}).

Once the electromagnetic load $\mathbf{F}$ and $\boldsymbol{\tau}_{\mathbf{r}_0}$ are fixed by the Maxwell field on $S_m$, a polarizable body converts that load into measurable mechanical exchange through its constitutive and elastic response \cite{beth1936,stratton1941}.

\subsection{Power flux, stress traction, and torque}
\label{sec:ch01.5.2}

How does cycle-averaged power leave a radiative shell, and how does mechanical torque arise at a material boundary? At each point on a surface, $\mathbf{S}_\mathrm{c}=\tfrac{1}{2}\E\times\H^*$ gives the local complex harmonic energy transport \cite{poynting1884,stratton1941}. Its outward normal component $\Re\{\mathbf{S}_\mathrm{c}\cdot\hat{\mathbf{n}}\}=\langle\mathbf{S}\rangle\cdot\hat{\mathbf{n}}$ is the cycle-averaged power crossing the surface element $dS$; integrating over a closed surface enclosing the source yields net radiated power through Eq.~\eqref{eq:ch1.power-balance-integral}. Near-field Poynting loops can be large while the outward integral on a radiative shell remains small---the reactive-versus-radiative distinction of section~\ref{sec:ch01.2.4}. In a lossless source-free exterior, the net signed real power $\int_{S_r}\Re\{\mathbf{S}_\mathrm{c}\cdot\hat{\mathbf{n}}\}\,dS$ is conserved on any enclosing sphere and equals $P_{\mathrm{rad}}$; a predominantly radiative shell is preferred when interpreting the \emph{local} field as far-zone radiation, not because the integral is undefined on other admissible surfaces \cite{poynting1884,stratton1941}.

Mechanical exchange uses the same Maxwell stress tensor. Traction $\mathbf{t}=\Tavg\cdot\hat{\mathbf{n}}$ is force per unit area on a boundary; integrating $\mathbf{t}$ yields net force, and integrating $\mathbf{r}\times\mathbf{t}$ about a fixed reference origin $\mathbf{r}_0$ yields torque \cite{maxwell1873,beth1936}. Force depends on the traction distribution alone; torque depends additionally on where moments are taken---the origin dependence established in section~\ref{sec:ch01.3}. Beth \cite{beth1936} observed mechanical rotation in quartz wave plates under circularly polarized light: the plate experienced torque because the field transferred angular action through stress at the material interface \cite{allen1992}. Stress-integral torque and Poynting flux on $\Gamma$ follow from $(\E,\H)$; harmonic labels assigned to $\E_t$ require an explicit angular representation and are not interchangeable with either integral \cite{stratton1941}. Subsection~\ref{sec:ch01.6.1} returns to Beth's pattern in fully worked form.

A polarizable body at $S_m$ converts the Maxwell load into observable mechanical motion through its constitutive response \cite{beth1936,stratton1941}. The electromagnetic prediction remains the stress integral; the observed twist on a torsion wire is the mechanical signature of that load at the boundary.

Origin declaration controls comparability of torque predictions. If origin changes
from \(\mathbf{r}_0\) to \(\mathbf{r}_0'=\mathbf{r}_0+\Delta\mathbf{r}\), the
deterministic torque integral satisfies
\begin{equation}
\label{eq:ch1.stress-origin-shift-torque}
\mathbf{T}_{\mathrm{det}}'
=
\mathbf{T}_{\mathrm{det}}-\Delta\mathbf{r}\times\mathbf{F}_{\mathrm{det}},
\end{equation}
with
\(\mathbf{F}_{\mathrm{det}}=\int_{S_m}\Tavg\cdot\hat{\mathbf{n}}\,dS\) and
\(\mathbf{T}_{\mathrm{det}}=\int_{S_m}(\mathbf{r}-\mathbf{r}_0)\times(\Tavg\cdot\hat{\mathbf{n}})\,dS\).
Eq.~\eqref{eq:ch1.stress-origin-shift-torque} is the surface form of
origin dependence required before torque comparisons are meaningful \cite{stratton1941}.

\begin{example}[Why torque depends on the chosen origin]
\label{ex:ch1.origin-shift}
With $\mathbf{r}'=\mathbf{r}-\Delta\mathbf{r}$, the stress torque transforms as $\boldsymbol{\tau}'=\boldsymbol{\tau}-\Delta\mathbf{r}\times\mathbf{F}$ (Eq.~\eqref{eq:ch1.stress-origin-shift-torque}). On a thin plate with $\mathbf{r}_0$ at the centre, shifting the origin by $\Delta\mathbf{r}=2\,\mathrm{mm}\,\hat{\mathbf{x}}$ while holding net force $\mathbf{F}_{\mathrm{det}}$ fixed changes the $z$-torque by $\Delta\tau_z=-F_y\,\Delta x$ when $\hat{\mathbf{u}}_\tau=\hat{\mathbf{z}}$ \cite{stratton1941}. Comparisons of stress torque with mechanical measurements must declare the same reference origin.
\end{example}

Surface selection is part of the field problem, not an afterthought \cite{stratton1941}. Two closed surfaces yield the same force and torque resultants only when they can be deformed into one another without crossing intervening sources, material bodies, or field singularities, the stress is evaluated on the same Maxwell-valid $(\E,\H)$ branch in the intervening source-free region, and storage terms are treated consistently \cite{maxwell1873,stratton1941}. A partial patch $S_m$ is not interchangeable with a full closed $\partial V_b$ unless the omitted segments and any edge contributions are separately accounted for \cite{stratton1941}.

\phantomsection\label{prop:ch1.stress-torque-admissibility}
Force and torque obtained from Maxwell stress integrals are meaningful in Chapter~\ref{ch:01} only when the integration surface, reference origin, and torque axis are declared with the field solution \cite{stratton1941,beth1936}.

Maxwell consistency fixes which field traces may carry flux on $\mathcal{F}_{\mathrm{rad}}$ \cite{stratton1941}.
\subsection{Field quantities, instrument readings, and inferred structure}
\label{sec:ch01.5.3}

% HARRINGTON-OPEN:sec:ch01.5.3
Maxwell theory determines $\E$ and $\H$ from sources, constitutive relations, boundary data, and the outgoing radiation condition \cite{maxwell1873,stratton1941}. A receiving structure samples the field through a physical coupling on a declared surface; an angular coefficient is obtained only after a basis and representation have been chosen \cite{poynting1884,beth1936,hansen1988}. Flux, traction, and stress-moment integrals are Maxwell field functionals; instrument readings and inferred modal labels are separate layers that must not be interchanged with those integrals without stating the connecting assumptions \cite{stratton1941}.

\subsection{Apertures, angular charts, and later radiation mappings}
\label{sec:ch01.5.4}

On a large sphere enclosing sources, tangential $\E_t(\theta,\phi)$ and $\H_t(\theta,\phi)$ specify the outgoing field angularly \cite{stratton1941}. A physical aperture selects directions and couples to $\E$ and $\H$ through its geometry and constitutive response \cite{hansen1988}. Angular spectra and harmonic charts organize the same tangential data for propagation and synthesis; they are not substitutes for conservation integrals on declared surfaces \cite{harrington2001,stratton1941}. Chapter~\ref{ch:02} connects this angular organization to radiation operators and dyadic Green functions on the outgoing branch fixed here \cite{harrington2001,tai1994}.

\phantomsection\label{prop:ch1.observable-inference-criterion}
Throughout Chapter~\ref{ch:01}, a reported number is either a Maxwell field functional on a declared surface---flux, traction, torque---or a quantity inferred from partial field information through a stated representation; the status must be clear \cite{stratton1941}.

\section{Examples and Counterexamples}
\label{sec:ch01.6}

\subsection{Angular momentum conservation in radiation}
\label{sec:ch01.6.1}

Circular polarization gives the electromagnetic field a definite handedness about the direction of propagation \cite{stratton1941,beth1936}. The measurable transport quantities of this chapter are Maxwell stress-moment integrals on declared surfaces, not phase portraits alone. Begin from the exact definition
\begin{equation}
\label{eq:ch1.ex61-flux-from-stress}
\dot J_{z,\mathrm{em,out}}
=-\oint_{S_R}\big[\mathbf{r}\times(\Tavg\cdot\hat{\mathbf{r}})\big]\cdot\hat{\mathbf{z}}\,dS,
\end{equation}
which is Eq.~\eqref{eq:ch1.outward-am-flux} with outward normal $\hat{\mathbf{n}}=\hat{\mathbf{r}}$ on a sphere $S_R$ in the source-free exterior \cite{stratton1941,beth1936}.

A single outgoing Hansen multipole on a source-free sphere $S_R$ carries radiated power through the Poynting integral and axial angular momentum through the Maxwell-stress moment. Those two surface integrals are distinct densities on the sphere; their ratio, after integration, is derived below for one electric-type or one magnetic-type multipole. A radiation-zone identity of Jackson's, of a different character, is recorded afterward. Fix $\exp(+\jj\omega t)$, outgoing $h_\ell^{(2)}(kr)=j_\ell(kr)-\jj y_\ell(kr)$, and the radiation condition $\H=\eta_0^{-1}\hat{\mathbf{r}}\times\E$ in the far zone \cite{stratton1941}. On $S_R$ in the source-free exterior define
\begin{equation}
\label{eq:ch1.ex61-power-def}
P_{\mathrm{out}}:=\oint_{S_R}\Re\{\mathbf{S}_\mathrm{c}\cdot\hat{\mathbf{r}}\}\,dS,\qquad
\mathbf{S}_\mathrm{c}=\tfrac12\E\times\H^*,
\end{equation}
with $\dot J_{z,\mathrm{em,out}}$ given by Eq.~\eqref{eq:ch1.ex61-flux-from-stress}. A single outgoing Hansen multipole of indices $(\ell,m)$ is specified by one electric-type amplitude $m_{\ell m}$ or one magnetic-type amplitude $n_{\ell m}$, with exterior fields
\begin{equation}
\label{eq:ch1.ex61-hansen-outgoing}
\E=-k^2 Z_0\bigl(m_{\ell m}\mathbf M_{\ell m}^+ + n_{\ell m}\mathbf N_{\ell m}^+\bigr),\qquad
\H=-k^2\bigl(m_{\ell m}\mathbf N_{\ell m}^+ - n_{\ell m}\mathbf M_{\ell m}^+\bigr),
\end{equation}
where $\mathbf M_{\ell m}^+=h_\ell^{(2)}(kr)\vec Y_{\ell m}^{\ell m}$ and $\mathbf N_{\ell m}^+=(1/k)\nabla\times\mathbf M_{\ell m}^+$. Write $\sigma=\sqrt{\ell(\ell+1)}$ (never $\ell(\ell+1)$ itself) and define the angular operator by $\vec L=\jj\mathbf r\times\nabla$, so that on the unit sphere
\[
L_\theta Y_\ell^m=\frac{-\jj\,\partial_\phi Y_\ell^m}{\sin\theta},\qquad
L_\phi Y_\ell^m=\jj\,\partial_\theta Y_\ell^m.
\]
The vector harmonic in \eqref{eq:ch1.ex61-hansen-outgoing} is the unit-norm field
\begin{equation}
\label{eq:ch1.ex61-vsh}
\vec Y_{\ell m}^{\ell m}=\frac{\vec L Y_\ell^m}{\sigma}
=\hat{\boldsymbol{\theta}}\,\frac{m Y_\ell^m}{\sigma\sin\theta}
+\hat{\boldsymbol{\phi}}\,\frac{\jj\,\partial_\theta Y_\ell^m}{\sigma}.
\end{equation}
For $\int_{4\pi}|Y_\ell^m|^2\,d\Omega=1$ the angular-gradient identity is $\int_{4\pi}\bigl(|\partial_\theta Y_\ell^m|^2+|\partial_\phi Y_\ell^m|^2/\sin^2\theta\bigr)\,d\Omega=\ell(\ell+1)$, hence $\int_{4\pi}|\vec Y_{\ell m}^{\ell m}|^2\,d\Omega=1$ \cite{jackson1999}. Jackson's $\mathbf X_{\ell m}$ \cite[Eqs.~(9.119)--(9.120)]{jackson1999} uses the same factor $1/\sqrt{\ell(\ell+1)}$ with his orbital operator $-\mathrm{j}\mathbf r\times\nabla$; the two operators differ by a global sign, which does not change the norm. Replacing $\sigma$ by $\ell(\ell+1)$ in \eqref{eq:ch1.ex61-vsh} would yield $\int|\vec Y|^2=1/[\ell(\ell+1)]$ and would shift the absolute power in Proposition~\ref{prop:ch1.single-hansen-power}. The manuscript $\mathbf N_{\ell m}^+=(1/k)\nabla\times\mathbf M_{\ell m}^+$ with $M_r=0$ has, on $S_R$ with $z=kR$, $h=h_\ell^{(2)}(z)$, and $u(z)=zh_\ell^{(2)}(z)$,
\begin{equation}
\label{eq:ch1.ex61-hansen-surface}
N_r=-\frac{\sigma}{z}\,h\,Y_\ell^m,\qquad
N_\theta=\frac{u'}{z}\,\jj(-Y_\phi),\qquad
N_\phi=\frac{u'}{z}\,\jj Y_\theta,
\end{equation}
where $Y_\theta,Y_\phi$ are the $\hat{\boldsymbol{\theta}},\hat{\boldsymbol{\phi}}$ components of $\vec Y_{\ell m}^{\ell m}$ in \eqref{eq:ch1.ex61-vsh} and $u'=\mathrm{d}u/\mathrm{d}z$. These components follow from the spherical curl of $h\vec Y$ with $M_r=0$ and are checked against the source-free Maxwell pair \eqref{eq:ch1.ex61-hansen-outgoing} in the proof below; they are not copied from an unverified Hansen appendix normalization \cite{stratton1941,jackson1999}. Active rotation acts as $(R_\alpha\mathbf F)(\mathbf r)=R_\alpha\mathbf F(R_\alpha^{-1}\mathbf r)$ with generator $\mathcal{J}_z=-\jj\partial_\phi+\mathcal{S}_z$ and $\mathcal{U}(\alpha)=e^{-\jj\alpha\mathcal{J}_z}$; on a pure $\vec Y_{\ell m}^{\ell m}$ envelope the orbital part contributes azimuthal index $m$, while the transverse vector structure fixes the spin contribution \cite{stratton1941}.

\begin{proposition}[Radiated power of one outgoing Hansen multipole]
\label{prop:ch1.single-hansen-power}
Let $(\E,\H)$ be the source-free outgoing vacuum pair \eqref{eq:ch1.ex61-hansen-outgoing} for fixed $(\ell,m)$ on a sphere $S_R$ in the source-free exterior, with arbitrary complex amplitudes $m_{\ell m}$, $n_{\ell m}$ (not both zero). Assume the manuscript curl pairing $\curl\E=-\jj\omega\mu_0\H$ and the orthogonality $\int_{4\pi}\vec Y_{\ell m}^{\ell m}\cdot\vec Y_{\ell'm'}^{\ell'm'*}\,d\Omega=\delta_{\ell\ell'}\delta_{mm'}$. Then
\begin{equation}
\label{eq:ch1.ex61-single-mode-power}
P_{\mathrm{out}}=\frac{k^2 Z_0}{2}\bigl(|m_{\ell m}|^2+|n_{\ell m}|^2\bigr).
\end{equation}
\end{proposition}
\begin{proof}
Let $c$ denote the single nonzero coefficient $m_{\ell m}$ or $n_{\ell m}$, and on $S_R$ write $z=kR$, $h=h_\ell^{(2)}(z)$, and $u'=u'(z)$ with $u=zh$. Insert the electric-type pair $\E=-k^2 Z_0 c\,\mathbf M_{\ell m}^+$, $\H=-k^2 c\,\mathbf N_{\ell m}^+$ from \eqref{eq:ch1.ex61-hansen-outgoing} into \eqref{eq:ch1.ex61-power-def}. Because $E_r=0$ and $H_r$ is the only radial magnetic component,
\[
\Re\{\mathbf S_\mathrm{c}\cdot\hat{\mathbf r}\}
=\frac{k^4|Z_0 c|^2}{2}\,\Re\bigl[(\mathbf M_{\ell m}^+\times\mathbf N_{\ell m}^{+*})_r\bigr],
\]
with $\mathbf M_{\ell m}^+=h\,\vec Y_{\ell m}^{\ell m}$ and $\mathbf N_{\ell m}^+$ given by \eqref{eq:ch1.ex61-hansen-surface}. Using $M_\theta=hY_\theta$, $M_\phi=hY_\phi$ and \eqref{eq:ch1.ex61-hansen-surface},
\[
(\mathbf M\times\mathbf N^*)_r
=-\jj\,\frac{h\,[u']^*}{z}\,\bigl(|Y_\theta|^2+|Y_\phi|^2\bigr),
\qquad
(\mathbf N\times\mathbf M^*)_r
=-\jj\,\frac{[u']\,h^*}{z}\,\bigl(|Y_\theta|^2+|Y_\phi|^2\bigr),
\]
with $[u']=u'(z)$ and $h=h_\ell^{(2)}(z)$. These are \emph{not} negatives as complex numbers; however $v=[u']/z$ gives $h v^*+v h^*=2\Re(h v^*)$, so
\[
\Re\{(\mathbf N\times\mathbf M^*)_r\}=-\Re\{(\mathbf M\times\mathbf N^*)_r\}
=\frac{1}{z}\,\Im(h[u']^*)\,|\vec Y_{\ell m}^{\ell m}|^2
\]
after $|Y_\theta|^2+|Y_\phi|^2=|\vec Y_{\ell m}^{\ell m}|^2$. Hence
\[
\Re\{\mathbf S_\mathrm{c}\cdot\hat{\mathbf r}\}
=\frac{k^4 Z_0 |c|^2}{2}\,\frac{1}{z}\,\Im(h[u']^*)\,|\vec Y_{\ell m}^{\ell m}|^2.
\]
Integrating over $dS=R^2\sin\theta\,d\theta\,d\phi$ and using $\int_{4\pi}|\vec Y_{\ell m}^{\ell m}|^2\,d\Omega=1$ gives
\[
P_{\mathrm{out}}=\frac{k^4 Z_0 |c|^2 R^2}{2z}\,\Im(h[u']^*).
\]
With \eqref{eq:ch1.ex61-hankel-wronskian}, $\Im(h[u']^*)=1/z$, so $P_{\mathrm{out}}=(k^2 Z_0/2)|c|^2$.
If only $n_{\ell m}\neq0$, interchange $\mathbf M_{\ell m}^+$ and $\mathbf N_{\ell m}^+$ with the signs in \eqref{eq:ch1.ex61-hansen-outgoing}; the same algebra with $H_r=0$ and $E_r$ from \eqref{eq:ch1.ex61-hansen-surface} gives \eqref{eq:ch1.ex61-single-mode-power} for a magnetic-type amplitude alone.

\medskip\noindent\emph{Same-index electric--magnetic superposition.}
Write $m=m_{\ell m}$, $n=n_{\ell m}$ and insert the general pair \eqref{eq:ch1.ex61-hansen-outgoing} into \eqref{eq:ch1.ex61-power-def}. With $M_r=0$,
\[
\E=-k^2 Z_0(m\mathbf M+n\mathbf N),\qquad
\H=-k^2(m\mathbf N-n\mathbf M),
\]
so
\[
(\E\times\H^*)_r
=k^4 Z_0\Bigl(|m|^2(\mathbf M\times\mathbf N^*)_r
-|n|^2(\mathbf N\times\mathbf M^*)_r
-mn^*(\mathbf M\times\mathbf M^*)_r
+n m^*(\mathbf N\times\mathbf N^*)_r\Bigr).
\]
Using \eqref{eq:ch1.ex61-hansen-surface} and \eqref{eq:ch1.ex61-vsh}, the tangential components satisfy $M_\theta=hY_\theta$, $M_\phi=hY_\phi$, and $Y_\theta=m Y_\ell^m/(\sigma\sin\theta)$, $Y_\phi=\jj\,\partial_\theta Y_\ell^m/\sigma$. Then
\[
(\mathbf M\times\mathbf M^*)_r
=h^2\bigl(Y_\theta Y_\phi^*-Y_\phi Y_\theta^*\bigr)
=-\,\jj\,\frac{h^2 m}{\sigma^2\sin\theta}\,\partial_\theta|Y_\ell^m|^2,
\]
which is \emph{purely imaginary} for real $\theta$ because $|Y_\ell^m|^2$ is real; hence $\Re\{(\mathbf M\times\mathbf M^*)_r\}=0$ pointwise, and the same argument with $u'/z$ in place of $h$ gives $\Re\{(\mathbf N\times\mathbf N^*)_r\}=0$. With $v=[u']/z$, $\Re\{(\mathbf N\times\mathbf M^*)_r\}=-\Re\{(\mathbf M\times\mathbf N^*)_r\}$ because $\Re(-\jj h v^* S)+\Re(-\jj v h^* S)=\Im(h v^*+v h^*)=0$. The interference term in $\Re\{(\E\times\H^*)_r\}$ therefore reduces to
\[
\frac{k^4 Z_0}{2}\,\Re\Bigl\{(|m|^2+|n|^2)(\mathbf M\times\mathbf N^*)_r\Bigr\},
\]
and \eqref{eq:ch1.ex61-single-mode-power} follows from the single-amplitude integral above with $|c|^2$ replaced by $|m|^2+|n|^2$.
\end{proof}

For real $z>0$ and outgoing $h=h_\ell^{(2)}(z)$, $u=zh$, the spherical-Hankel Wronskian gives
\begin{equation}
\label{eq:ch1.ex61-hankel-wronskian}
\Im\bigl(h\,[u']^*\bigr)=\frac{1}{z}.
\end{equation}
This follows from $j_\ell y'_\ell-j'_\ell y_\ell=1/z^2$ and $h=j_\ell-\jj y_\ell$, since $u'=h+zh'$ and a short complex-algebra reduction yields $\Im(h[u']^*)=z\,\Im(hh'^*)=1/z$. For the incoming branch $h_\ell^{(1)}=j_\ell+\jj y_\ell$, the same Wronskian reverses sign and $\Im(h[u']^*)=-1/z$, so $P_{\mathrm{out}}$ and $\dot J_{z,\mathrm{em,out}}$ change sign together while their ratio is unchanged. The same identity makes each flux in Propositions~\ref{prop:ch1.single-hansen-power} and~\ref{prop:ch1.single-hansen-transport} independent of $R$.

On the sphere $S_R$, the axial stress moment reduces geometrically to the single traction component $T_{r\phi}$. With outward normal $\hat{\mathbf{r}}$ and $\hat{\boldsymbol{\theta}}\cdot\hat{\mathbf{z}}=-\sin\theta$,
\[
\bigl[\mathbf{r}\times(\Tavg\cdot\hat{\mathbf{r}})\bigr]_z=R\sin\theta\,T_{r\phi},
\]
so that Eq.~\eqref{eq:ch1.ex61-flux-from-stress} is
\[
\dot J_{z,\mathrm{em,out}}=-R^3\int_0^{\pi}\int_0^{2\pi}T_{r\phi}\,\sin^2\theta\,d\phi\,d\theta.
\]
The cycle-averaged tensor \eqref{eq:ch1.maxwell-stress} gives $T_{r\phi}=\tfrac12\mathrm{Re}(\varepsilon_0 E_r E_\phi^*+\mu_0 H_r H_\phi^*)$. For a pure electric Hansen pair, $M_r=0$ and therefore $E_r=0$, leaving only the magnetic term. Write $z=kR$, $\sigma=\sqrt{\ell(\ell+1)}$ as in \eqref{eq:ch1.ex61-vsh}, $u(z)=z h_\ell^{(2)}(z)$, with the prime denoting $d/dz$. The radial partner is $N_r=-(\sigma/z)h Y_\ell^m$; the factor is $\sigma$, not $\ell(\ell+1)$. The $\theta$ densities of Poynting and of $T_{r\phi}$ are not pointwise proportional. Distinct-$(\ell,m)$ superpositions and full-vector $\mathcal{J}_z$ eigenfields are treated separately below.

\begin{proposition}[Axial stress flux of one outgoing Hansen multipole]
\label{prop:ch1.single-hansen-transport}
Let $(\E,\H)$ be as in Proposition~\ref{prop:ch1.single-hansen-power}, with arbitrary complex $m_{\ell m}$, $n_{\ell m}$ at fixed $(\ell,m)$, $\ell\ge 1$, and $m\in\{-\ell,\ldots,\ell\}$. Then, at every radius $R$ in the source-free exterior,
\begin{equation}
\label{eq:ch1.ex61-single-mode-transport}
\dot J_{z,\mathrm{em,out}}=\frac{m}{\omega}\,P_{\mathrm{out}}.
\end{equation}
\end{proposition}
\begin{proof}
Let $c$ be the single nonzero amplitude and write $h=h_\ell^{(2)}(z)$. For the electric-type pair, $\E=-k^2 Z_0 c\,\mathbf M_{\ell m}^+$ and $\H=-k^2 c\,\mathbf N_{\ell m}^+$, so $E_r=0$ and $T_{r\phi}=\tfrac12\mu_0\mathrm{Re}(H_r H_\phi^*)$. The components on $S_R$ are
\[
H_r=k^2 c\,\frac{\sigma}{z}\,h\,Y_\ell^m,\qquad
H_\phi=-k^2 c\,\frac{u'}{z}\,\jj Y_\theta.
\]
Conjugation on real $z$ gives $H_\phi^*=k^2 c^*(u'/z)\,\jj Y_\theta^*$, hence $H_r H_\phi^*=k^4|c|^2(m/z^2)\,\jj h[u']^*\,|Y_\ell^m|^2/\sin\theta$ after $Y_\theta=m Y_\ell^m/(\sigma\sin\theta)$ from \eqref{eq:ch1.ex61-vsh}. Replacing $\sigma$ by $\Lambda=\ell(\ell+1)$ in \eqref{eq:ch1.ex61-vsh} and \eqref{eq:ch1.ex61-hansen-surface} while holding the physical amplitude $c$ fixed rescales $|\vec Y|^2$ and breaks absolute $P_{\mathrm{out}}$; replacing $\sigma$ by $\Lambda$ in \emph{both} $N_r$ and $Y_\theta$ can cancel in the product $H_r H_\phi^*$ at fixed $|Y_\ell^m|$ without restoring unit vector-harmonic normalization---that cancellation is not a valid normalization choice and does not rescue the power coefficient (see the $\Lambda$ negative control in the audit engines). Then $\mathrm{Re}(\jj h[u']^*)=-\mathrm{Im}(h[u']^*)$. Inserting $dS=R^2\sin\theta\,d\theta\,d\phi$ into $\dot J_{z,\mathrm{em,out}}=-R^3\iint T_{r\phi}\sin^2\theta\,d\phi\,d\theta$, the two minus signs cancel and $\int_{4\pi}|Y_\ell^m|^2\,d\Omega=1$ produces
\[
\dot J_{z,\mathrm{em,out}}=\tfrac12\mu_0 k^2 |c|^2 m R\,\Im(h[u']^*).
\]
The radial Poynting density is $\Re(S_r)=\tfrac12 k^4 Z_0|c|^2 z^{-1}\mathrm{Im}(h[u']^*)\,|\vec Y_{\ell m}^{\ell m}|^2$. After $\int|\vec Y_{\ell m}^{\ell m}|^2\,d\Omega=1$,
\[
P_{\mathrm{out}}=\tfrac12 k^3 Z_0 |c|^2 R\,\Im(h[u']^*).
\]
The radial bilinear is common, so the ratio $\dot J_{z,\mathrm{em,out}}/P_{\mathrm{out}}=\mu_0 m/(k Z_0)$ is independent of both $R$ and the sign of $\mathrm{Im}(h[u']^*)$. Vacuum SI law $k Z_0=\mu_0\omega$ yields \eqref{eq:ch1.ex61-single-mode-transport}. Equation~\eqref{eq:ch1.ex61-hankel-wronskian} then gives $P_{\mathrm{out}}=k^2 Z_0|c|^2/2$ as in Proposition~\ref{prop:ch1.single-hansen-power}. For $m=0$, $Y_\theta=0$ and both sides of \eqref{eq:ch1.ex61-single-mode-transport} vanish. Negative $m$ enters only through $Y_\theta$; $P_{\mathrm{out}}$ is independent of $\mathrm{sgn}(m)$.

For the magnetic-type pair, $n_{\ell m}=c$ and $m_{\ell m}=0$ in \eqref{eq:ch1.ex61-hansen-outgoing}, so $\E=-k^2 Z_0 c\,\mathbf N_{\ell m}^+$ and $\H=k^2 c\,\mathbf M_{\ell m}^+$. Then $H_r=0$ and $T_{r\phi}=\tfrac12\varepsilon_0\mathrm{Re}(E_r E_\phi^*)$, with $E_r=k^2 Z_0 c(\sigma/z)h Y_\ell^m$ and $E_\phi=-k^2 Z_0 c(u'/z)\,\jj Y_\theta$. The product $E_r E_\phi^*$ is $Z_0^2$ times the electric-type $H_r H_\phi^*$, and $\varepsilon_0 Z_0^2=\mu_0$ restores the same $\dot J_{z,\mathrm{em,out}}$. The map $(\E,\H)_{\mathrm{mag}}=(Z_0\mathbf H_{\mathrm{el}},-\E_{\mathrm{el}}/Z_0)$ sends the electric-type pair onto the magnetic-type pair and leaves both $\Re(S_r)$ and $T_{r\phi}$ unchanged.

\medskip\noindent\emph{Same-index superposition.}
Write $a=a_{\ell m}$, $b=b_{\ell m}$ for the electric- and magnetic-type amplitudes and keep $m$ for the azimuthal index of $Y_\ell^m$. With $M_r=0$, $N_r=-(\sigma/z)hY_\ell^m$, $E_r=-k^2 Z_0 b N_r$, and $H_r=-k^2 a N_r$. The $\phi$-components are $E_\phi=-k^2 Z_0(a M_\phi+b N_\phi)$ and $H_\phi=-k^2(a N_\phi-b M_\phi)$. With $T_{r\phi}=\tfrac12\varepsilon_0\mathrm{Re}(E_r E_\phi^*)+\tfrac12\mu_0\mathrm{Re}(H_r H_\phi^*)$,
\begin{align*}
\mathrm{Re}(E_r E_\phi^*)
&=k^4 Z_0^2|b|^2\,\frac{\sigma}{z}\,|h Y_\ell^m|^2\,\Re\Bigl(\frac{u'}{z}\,\jj Y_\theta\Bigr)
+k^4 Z_0^2\,\Re\bigl(b a^* N_r M_\phi^*+b b^* N_r N_\phi^*\bigr),\\
\mathrm{Re}(H_r H_\phi^*)
&=k^4|a|^2\,\frac{\sigma}{z}\,|h Y_\ell^m|^2\,\Re\Bigl(\frac{u'}{z}\,\jj Y_\theta\Bigr)
+k^4\,\Re\bigl(a b^* N_r M_\phi^*-a a^* N_r N_\phi^*\bigr),
\end{align*}
using $N_\phi=(\jj u'/z)Y_\theta$ and $M_\phi=hY_\phi$. The four cross products $N_r M_\phi^*$, $N_r N_\phi^*$, $N_r M_\phi^*$ in the $ab^*$ and $a^*b$ slots combine so that every cross contribution to $\mathrm{Re}(T_{r\phi})$ is proportional to $\Re\{(\mathbf M\times\mathbf M^*)_r\}$, $\Re\{(\mathbf N\times\mathbf N^*)_r\}$, or $\Re\{(\mathbf M\times\mathbf N^*)_r+\ (\mathbf N\times\mathbf M^*)_r\}$; each of these real parts vanishes pointwise as in the Poynting proof ($\Re\{(\mathbf M\times\mathbf M^*)_r\}=0$ and $\Re\{(\mathbf M\times\mathbf N^*)_r+(\mathbf N\times\mathbf M^*)_r\}=0$ because $\Im(h v^*+v h^*)=0$ for $v=[u']/z$). The surviving $|a|^2$ and $|b|^2$ terms add to
\[
T_{r\phi}=\tfrac12\mu_0 k^2\bigl(|a|^2+|b|^2\bigr)\,\Re\Bigl(H_r H_\phi^*\Big|_{a_{\ell m}\ \text{only}}\Bigr)
\]
up to the electric--magnetic duality replacement $E_r E_\phi^*\leftrightarrow Z_0^2 H_r H_\phi^*$, producing the single-type density $\propto m\,\Im(h[u']^*)\,|Y_\ell^m|^2/\sin\theta$. With $\dot J_{z,\mathrm{em,out}}=-R^3\iint T_{r\phi}\sin^2\theta\,d\phi\,d\theta$ and $\int_{4\pi}|Y_\ell^m|^2\,d\Omega=1$,
\[
\dot J_{z,\mathrm{em,out}}=\tfrac12\mu_0 k^2\bigl(|a_{\ell m}|^2+|b_{\ell m}|^2\bigr)\,m R\,\Im(h[u']^*),
\]
which matches the $|a|^2+|b|^2$ Poynting integral and yields \eqref{eq:ch1.ex61-single-mode-transport} via $kZ_0=\mu_0\omega$. This step does not infer the stress identity from the power identity: the $ab^*$ and $a^*b$ stress terms are removed by the real-part cross-product identities above, not by their Poynting counterparts alone.
\end{proof}

\begin{remark}[Scope of the Hansen flux laws]
\label{rem:ch1.single-multipole-scope}
Propositions~\ref{prop:ch1.single-hansen-power} and~\ref{prop:ch1.single-hansen-transport} apply to the \emph{full} same-index outgoing pair \eqref{eq:ch1.ex61-hansen-outgoing} with arbitrary complex $m_{\ell m}$ and $n_{\ell m}$. They do \emph{not} extend to sums over distinct $(\ell,m)$, to full-vector $\mathcal{J}_z$ eigenfields, to the equatorial ring-source coefficient map, or to the separable $(m,\sigma)$ bookkeeping symbols below unless separately proved; Jackson's shell-content identity \eqref{eq:ch1.ex61-jackson-shell} remains distinct from the Maxwell-stress flux \eqref{eq:ch1.ex61-flux-from-stress}.
\end{remark}

A classical radiation-zone result of a different character is recorded by Jackson \cite{jackson1999}. In SI units, for outgoing multipole fields in the zone $kr\gg 1$, the time-averaged energy $U$ and angular momentum $\mathbf M$ contained in a spherical shell are obtained from the field energy density and from the moment of the field momentum density, not from the Maxwell-stress flux \eqref{eq:ch1.ex61-flux-from-stress}. For a pure electric multipole of order $(\ell,m)$, comparison of Jackson's shell energy \cite[Eq.~(9.136)]{jackson1999} with the $z$ component of shell angular momentum \cite[Eqs.~(9.143)--(9.144)]{jackson1999} gives
\begin{equation}
\label{eq:ch1.ex61-jackson-shell}
\frac{\mathrm{d}M_z}{\mathrm{d}r}=\frac{m}{\omega}\,\frac{\mathrm{d}U}{\mathrm{d}r},
\end{equation}
independent of radius in that zone. Jackson represents the same real field by a phasor $\mathbf E_{\mathrm{J}}$ with time factor $\exp(-\mathrm{j}\omega t)$ and outgoing radial function $h_\ell^{(1)}(kr)$, whose far-zone factor is $\exp(+\mathrm{j}kr)$. The manuscript uses $\mathbf E$ with $\exp(+\jj\omega t)$ and $h_\ell^{(2)}(kr)$, whose far-zone factor is $\exp(-\jj kr)$. For real $kr$, $h_\ell^{(2)}(kr)=\bigl[h_\ell^{(1)}(kr)\bigr]^*$, and the two phasors are complex conjugates of one another when they encode the same $\mathbf E_{\mathrm{phys}}(\mathbf r,t)$. The azimuthal index $m$ in $Y_\ell^m$ is unchanged, and the ratio $m/\omega$ is therefore the same number in both conventions. Equation~\eqref{eq:ch1.ex61-jackson-shell} is Jackson's radiation-zone shell-content identity, distinct from the Maxwell-stress flux \eqref{eq:ch1.ex61-single-mode-transport}. Identifying $\mathrm{d}U/\mathrm{d}r$ or $\mathrm{d}M_z/\mathrm{d}r$ with the manuscript surface fluxes $P_{\mathrm{out}}$ and $\dot J_{z,\mathrm{em,out}}$ would require a separate argument relating shell content to Maxwell-stress and Poynting surface integrals; that identification is not made here. Jackson's Problem~9.8 is an electric-dipole flux exercise at infinity and notes that angular momentum is not carried by the transverse radiation-zone fields alone \cite{jackson1999}.

\begin{remark}[Vector generator and mixed multipoles]
\label{rem:ch1.am-power-target}
Proposition~\ref{prop:ch1.single-hansen-transport} is proved for the full same-index pair \eqref{eq:ch1.ex61-hansen-outgoing}. A stronger statement that an arbitrary field belonging to an eigenspace of the full vector generator $\mathcal{J}_z$ with eigenvalue $M$ would satisfy
\begin{equation}
\label{eq:ch1.ex61-total-M-transport}
\dot J_{z,\mathrm{em,out}}=\frac{M}{\omega}\,P_{\mathrm{out}}
\end{equation}
is not proved. Distinct $(\ell,m)$ contributions and electric--magnetic cross terms must be controlled separately. In Jackson's radiation-zone shell calculation the mixed electric--magnetic superposition includes interference between multipoles whose $\ell$ values differ by one \cite[Eq.~(9.145)]{jackson1999}. The separable template of Eq.~\eqref{eq:ch1.ex61-separable-mode} is not a substitute for \eqref{eq:ch1.ex61-hansen-outgoing}.
\end{remark}

\paragraph{Equatorial ring source (separate construction).}
Let $I_m$ denote the $m$th azimuthal Fourier component (amperes) of a thin equatorial ring of radius $a$ at $\theta=\pi/2$, with line-current phasor $I_m e^{\jj m\phi}$. A single $m\neq0$ harmonic is one term of a periodic feed $I(\phi,t)=\mathrm{Re}\{\sum_m I_m e^{\jj m\phi}e^{\jj\omega t}\}$; surface charge on the ring is fixed by $\nabla\cdot\mathbf K+\jj\omega\sigma=0$ on the support when the source is modeled as a surface current $\mathbf K$ \cite{stratton1941}. Projecting the filament onto the regular spherical vector waves is a source-expansion step. The retarded dyadic Green function that would convert $I_m$ into a unique pair of outgoing Hansen coefficients is not constructed in this chapter; that mapping is deferred to the source-expansion development of later chapters. Propositions~\ref{prop:ch1.single-hansen-power} and~\ref{prop:ch1.single-hansen-transport} apply to an already identified exterior Hansen pair, and do not themselves supply the coefficient map from $I_m$.
\begin{corollary}[Fixed-axis spin plus azimuth only when $M=m+\sigma_z$ is proved]
\label{cor:ch1.m-plus-sigma}
With $\hat{\mathbf{e}}_\sigma=(\hat{\mathbf{x}}+\jj\sigma\hat{\mathbf{y}})/\sqrt{2}$ and $\mathcal{J}_z=-\jj\partial_\phi+\mathcal{S}_z$, $\mathcal{S}_z\hat{\mathbf{e}}_\sigma=\sigma\hat{\mathbf{e}}_\sigma$, the relation $\dot J_{z,\mathrm{em,out}}=(m+\sigma)P/\omega$ would follow from Eq.~\eqref{eq:ch1.ex61-total-M-transport} only when the \emph{full} far-zone vector mode is an exact $\mathcal{J}_z$ eigenmode with $M=m+\sigma$ under the declared conventions---not from the single-multipole identities of Propositions~\ref{prop:ch1.single-hansen-power} and~\ref{prop:ch1.single-hansen-transport}, and not for the literal $\phi$-only template of Eq.~\eqref{eq:ch1.ex61-separable-mode} without a Maxwell-valid $\theta$ envelope \cite{allen1992,cerjan2011,monteiro2009}.
\end{corollary}

Restrict attention to a \emph{single} outgoing separable mode on a predominantly radiative shell: propagation axis $\hat{\mathbf{z}}$, azimuthal dependence $e^{\jj m\phi}$ in the transverse phasor, and transverse polarization fixed in the circular basis $\hat{\mathbf{e}}_\pm=(\hat{\mathbf{x}}\pm\jj\hat{\mathbf{y}})/\sqrt{2}$ with helicity label $\sigma=\pm1$ matching the $\hat{\mathbf{e}}_\sigma$ used in the mode \cite{stratton1941,allen1992}. Let the modal electric field on $S_R$ be
\begin{equation}
\label{eq:ch1.ex61-separable-mode}
\mathbf{E}_m=\frac{E_0}{R}\,e^{\jj m\phi}\,\hat{\mathbf{e}}_\sigma,
\qquad
\mathbf{H}_m=\frac{1}{\eta_0}\,\hat{\mathbf{z}}\times\mathbf{E}_m,
\end{equation}
with $E_0$ independent of $\phi$ and $R$ the sphere radius in the far zone where $|\E|,|\H|\propto 1/R$ and $\hat{\mathbf{z}}\cdot\E=0$ \cite{stratton1941}. The pair \eqref{eq:ch1.ex61-separable-mode} is a \emph{modal} template on $S_R$: polarization is fixed in the laboratory circular basis, and the modal power scale $P_{\mathrm{modal}}$ is \emph{defined} in Eq.~\eqref{eq:ch1.ex61-modal-power}, not by reading a phase portrait alone. A full vector Maxwell solution on $4\pi$ generally requires a $\theta$-dependent transverse envelope (for example, a vector-harmonic or declared outgoing sheet) so that $\E\cdot\hat{\mathbf{r}}=0$ and $\curl\E=-\jj\omega\mu_0\H$ hold together; Eqs.~\eqref{eq:ch1.ex61-orbital-from-stress}--\eqref{eq:ch1.ex61-helicity-transport-law} are evaluated in the separable $(m,\sigma)$ class under that normalization \cite{stratton1941,allen1992}. Modal radiated power carried by this mode alone is
\begin{equation}
\label{eq:ch1.ex61-modal-power}
P_{\mathrm{modal}}:=\frac{|E_0|^2}{2\eta_0},
\end{equation}
a \emph{declared} power scale for the separable $(m,\sigma)$ class---not the signed full-sphere Poynting integral of the literal template. For Eq.~\eqref{eq:ch1.ex61-separable-mode} with $\mathbf{H}_m=\eta_0^{-1}\hat{\mathbf{z}}\times\mathbf{E}_m$, $\oint_{S_R}\Re\{\mathbf{S}_\mathrm{c}\cdot\hat{\mathbf{r}}\}\,dS$ can vanish by symmetry while $\mathbf{E}\cdot\hat{\mathbf{r}}\neq0$ on much of $4\pi$; the same template yields $\dot J_{z,\mathrm{em,out}}\not\approx(m+\sigma)P_{\mathrm{modal}}/\omega$ under direct stress integration \cite{stratton1941}.

\paragraph{Plane-wave stress countercheck (symmetric open disk).}
For a uniform outgoing circularly polarized plane wave $\E=(E_0/\sqrt{2})(\hat{\mathbf{x}}+\jj\sigma\hat{\mathbf{y}})e^{-\jj kz}$, $\H=\eta_0^{-1}\hat{\mathbf{z}}\times\E$ (Eq.~\eqref{eq:ch1.helicity-plane-wave} with real $E_0$), the Maxwell stress tensor is spatially constant. On a finite transverse disk centered on the beam axis, $\oint[\mathbf{r}\times(\Tavg\cdot\hat{\mathbf{z}})]_z\,dS=0$ by azimuthal symmetry: power through the disk is $P=|E_0|^2 A/(2\eta_0)$, but the symmetric stress-moment integral does \emph{not} yield $\sigma P/\omega$ \cite{stratton1941,beth1936,stewart2005}. Finite-beam edge contributions can restore nonzero stress moments in more general geometries \cite{stewart2005}; this countercheck is scoped to the stated uniform field and symmetric disk only. Beth \cite{beth1936} torque on an \emph{anisotropic} load is therefore not reproduced by this uniform-beam disk formula; the modal symbol $\dot J_{z,\mathrm{spin,out}}\stackrel{\mathrm{modal\ class}}{=}\sigma P_{\mathrm{modal}}/\omega$ is bookkeeping, not a consequence of Eq.~\eqref{eq:ch1.ex61-flux-from-stress} for the infinite uniform plane wave on a centered disk.

\FloatBarrier
\paragraph{Modal index bookkeeping (not a stress theorem).}
Because $\mathbf{E}_m\propto e^{\jj m\phi}$ with $\phi$-independent envelope on the shell, $\partial_\phi\mathbf{E}_m=\jj m\,\mathbf{E}_m$ and therefore
\begin{equation}
\label{eq:ch1.ex61-phi-identity}
\Im\{\mathbf{E}_m^*\cdot\partial_\phi\mathbf{E}_m\}=m\,|\mathbf{E}_m|^2
\end{equation}
for real $m$. Orthogonality of azimuthal modes on $S_R$ is
\begin{equation}
\label{eq:ch1.ex61-phi-orthogonality}
\int_0^{2\pi} e^{\jj(m-m')\phi}\,d\phi=2\pi\,\delta_{mm'},
\end{equation}
so distinct $(m,\sigma)$ modes do not mix in signed \emph{net} flux when they are orthogonal on $S_R$ in the usual far-zone inner product \cite{stratton1941,allen1992}. Physical outward flux $\dot J_{z,\mathrm{em,out}}$ is defined only by the Maxwell-stress integral Eq.~\eqref{eq:ch1.ex61-torque-flux} (equivalently Eq.~\eqref{eq:ch1.outward-am-flux}). The following \emph{modal bookkeeping} symbols $\dot J_{z,\mathrm{modal,bk}}$ pair azimuthal index $m$ and helicity label $\sigma$ with the declared scale $P_{\mathrm{modal}}$; they are \textbf{not} consequences of Eq.~\eqref{eq:ch1.ex61-flux-from-stress} for the literal template of Eq.~\eqref{eq:ch1.ex61-separable-mode} and must not be substituted for $\dot J_{z,\mathrm{em,out}}$ without a separate Maxwell-valid mode proof \cite{allen1992,stratton1941}:
\begin{equation}
\label{eq:ch1.ex61-orbital-from-stress}
\dot J_{z,\mathrm{orb,bk}}\;\stackrel{\mathrm{modal\ bookkeeping}}{=}\;\frac{m}{\omega}\,P_{\mathrm{modal}},
\qquad
\dot J_{z,\mathrm{spin,bk}}\;\stackrel{\mathrm{modal\ bookkeeping}}{=}\;\frac{\sigma}{\omega}\,P_{\mathrm{modal}},
\end{equation}
\begin{equation}
\label{eq:ch1.ex61-helicity-transport-law}
\dot J_{z,\mathrm{modal,bk}}\;\stackrel{\mathrm{modal\ bookkeeping}}{=}\;\frac{m+\sigma}{\omega}\,P_{\mathrm{modal}}.
\end{equation}
Paraxial and beam models in Chapter~\ref{ch:02} may justify the same $(m,\sigma)$ pairing under explicit approximations \cite{allen1992}. Beth \cite{beth1936} measured torque on an anisotropic wave plate; that experiment supports load-torque sign conventions (Eq.~\eqref{eq:ch1.ex61-torque-flux-preview}), not a universal closed-sphere or uniform-plane-wave stress identity for Eq.~\eqref{eq:ch1.ex61-separable-mode} \cite{beth1936,stratton1941}. For general full-vector superpositions, $\dot J_{z,\mathrm{em,out}}$ and $P_{\mathrm{rad}}=\oint_{S_R}\Re(\mathbf{S}_\mathrm{c}\cdot\hat{\mathbf{n}})\,dS$ must be computed from $(\E,\H)$ on the declared surface \cite{stratton1941,beth1936}.

\FloatBarrier
In steady harmonic operation, torque on the source support and outward field angular-momentum transport are related by opposite signs:
\begin{equation}
\label{eq:ch1.ex61-torque-flux-preview}
\boldsymbol{\tau}_{z,\mathrm{source}}=-\dot J_{z,\mathrm{em,out}},
\qquad
\boldsymbol{\tau}_{z,\mathrm{load}}=+\dot J_{z,\mathrm{em,out}}.
\end{equation}
Eq.~\eqref{eq:ch1.ex61-torque-flux-preview} is the rotational face of Eqs.~\eqref{eq:ch1.angular-balance-integral} and~\eqref{eq:ch1.outward-am-flux}. Beth \cite{beth1936} made the conservation law visible as mechanical torque on a quartz wave plate under circularly polarized light: the plate changes the polarization state of the transmitted beam, so the load torque follows from the net angular-momentum flux into the plate's control surface (incident minus outgoing), not from the source's isolated outward flux alone \cite{beth1936,stratton1941}. The torque is not inferred from the appearance of the polarization ellipse; it follows from the stress-moment balance at the anisotropic interface \cite{beth1936,stratton1941}.

The radiation example is a localized $z$-directed electric dipole on finite support
enclosed by a far-zone sphere $S_R$; Figure~\ref{fig:ch1.ex61-beth-helical-sketch} is a torque-balance schematic and does not, by itself, specify a rotating dipole moment or a quantitative closed-sphere stress proof for Eq.~\eqref{eq:ch1.ex61-helicity-transport-law}. On $S_R$ in the source-free exterior,
$P_{\mathrm{rad}}=\oint_{S_R}\Re\{\mathbf{S}_\mathrm{c}\cdot\hat{\mathbf{n}}\}\,dS$ is
finite and independent of $R$ on any enclosing sphere in the lossless source-free exterior
\cite{stratton1941,beth1936}. A predominantly radiative shell makes its interpretation as outward radiation especially transparent. The dipole example uses $P_{\mathrm{rad}}$ from stress/Poynting integration on $S_R$; it does \emph{not} identify $P_{\mathrm{rad}}$ with $P_{\mathrm{modal}}$ or assert Eq.~\eqref{eq:ch1.ex61-helicity-transport-law} without a separate Maxwell-valid mode construction \cite{stratton1941}.

\begin{figure}[b]
\centering
\ChOneFigBlock{%
\begin{tikzpicture}[ch1geo]
 \draw[ch1shell, dashed, line width=0.45pt] (0,0) circle (0.95);
 \node[ch1oncanvas, font=\scriptsize] at (0,1.12) {$S_R$};
 \draw[ch1axis, line width=0.5pt] (0,-0.15) -- (0,1.35);
 \node[ch1oncanvas, font=\scriptsize] at (0.12,1.35) {$\hat{\mathbf{z}}$};
 \draw[-{Latex[length=1.6mm,width=1mm]}, line width=0.5pt]
   (0.85,0.25) arc (15:75:0.85);
 \node[ch1oncanvas, font=\scriptsize] at (1.05,0.55) {$\mathbf{S}_\mathrm{c}$};
 \draw[line width=0.5pt] (-0.95,0) -- (-0.55,0);
 \draw[-{Latex[length=1.6mm,width=1mm]}, line width=0.55pt]
   (-0.75,0) -- (-0.75,0.4);
 \node[ch1oncanvas, font=\scriptsize] at (-1.05,0.5) {$\tau_z$};
 \node[ch1oncanvas, font=\scriptsize, align=center] at (0,-1.15)
 {torque on matter vs.\ flux through $S_R$};
\end{tikzpicture}%
}
\caption{Circularly polarized radiation from a localized dipole enclosed by far-zone sphere $S_R$: outward Poynting flux balances finite radiated power; stress integration gives $\boldsymbol{\mathcal{T}}_J$ and outward transport $\boldsymbol{\Phi}_{J,\mathrm{out}}=-\boldsymbol{\mathcal{T}}_J$; in steady operation the source support feels $\boldsymbol{\tau}_{\mathrm{source}}=-\boldsymbol{\Phi}_{J,\mathrm{out}}$ while an absorbing load feels $\boldsymbol{\tau}_{\mathrm{load}}=+\boldsymbol{\Phi}_{J,\mathrm{out}}$ \cite{beth1936,stratton1941}.}
\label{fig:ch1.ex61-beth-helical-sketch}
\end{figure}

Figure~\ref{fig:ch1.ex61-beth-helical-sketch} shows the localized $z$-directed dipole radiator in circularly polarized far-zone operation, enclosed by sphere $S_R$ at stand-off \(R=5\lambda\), where
\(|P_{\mathrm{reac}}(R)|/P_{\mathrm{rad}}\ll 1\) in the sense of Subsection~\ref{sec:ch01.2.4}---net reactive power no longer dominates the normal flux pattern on the shell---so the stress integral names radiation rather than
near-zone oscillation \cite{stratton1941}.

On a predominantly radiative shell, fields scale as \(|\E|\propto 1/R\) with the azimuthal phase of the dominant mode; Maxwell stress Eq.~\eqref{eq:ch1.maxwell-stress} with outward normal \(\hat{\mathbf{n}}=\hat{\mathbf{r}}\) supplies the integrand \([\mathbf{r}\times(\Tavg\cdot\hat{\mathbf{n}})]\cdot\hat{\mathbf{z}}\). When multiple modes contribute, Eq.~\eqref{eq:ch1.ex61-multimode-torque} sums separable \((m,\sigma)\) contributions from Eqs.~\eqref{eq:ch1.radiation-intensity-Ur}--\eqref{eq:ch1.directional-am-flux}; each mode carries a definite \((m+\sigma)/\omega\) factor only under the stated modal assumptions \cite{stratton1941}. Evaluating the same balance on \(R=\lambda\) without stating \(|P_{\mathrm{reac}}(r)|/P_{\mathrm{rad}}\) can attribute reactive-dominated storage exchange to far-zone transport \cite{stratton1941}.

For steady harmonic operation, storage variation over one cycle-average window is
small. The inherited angular-momentum balance therefore reduces to
\begin{equation}
\label{eq:ch1.ex61-torque-flux}
\dot J_{z,\mathrm{em,out}}
=
-\oint_{S_R}
\big[\mathbf{r}\times(\Tavg\cdot\hat{\mathbf{n}})\big]\cdot\hat{\mathbf{z}}\,dS
=-\mathcal{T}_{J,z},
\end{equation}
where \(\mathcal{T}_{J,z}\) is the $z$-component of the Maxwell-stress moment on \(S_R\). Outward transport \(\dot J_{z,\mathrm{em,out}}=\Phi_{J,z,\mathrm{out}}\) is the negative of the stress-moment integral under the convention of Eq.~\eqref{eq:ch1.outward-am-flux}.

For an absorbing load that intercepts the declared net angular-momentum flux through its enclosing control surface, \(\tau_{z,\mathrm{load}}=\dot J_{z,\mathrm{em,out}}\); for the enclosed source support, \(\tau_{z,\mathrm{source}}=-\dot J_{z,\mathrm{em,out}}=\mathcal{T}_{J,z}\) in steady operation \cite{beth1936}. A partially transmitting or polarization-changing load requires the full incoming and outgoing balance at that body’s surface; the source’s outward flux is not automatically equal to the load torque \cite{beth1936,stratton1941}.

When a declared modal-class evaluator (Eq.~\eqref{eq:ch1.ex61-helicity-transport-law}) is adopted for a particular beam or aperture model, Eq.~\eqref{eq:ch1.ex61-torque-flux} gives the corresponding torque bookkeeping
\begin{equation}
\label{eq:ch1.worked-helical-torque}
\tau_{z,\mathrm{source}}
\;\stackrel{\mathrm{modal\ class}}{=}\;
-\frac{m+\sigma}{\omega}\,P_{\mathrm{modal}},
\qquad
\tau_{z,\mathrm{load}}
\;\stackrel{\mathrm{modal\ class}}{=}\;
+\frac{m+\sigma}{\omega}\,P_{\mathrm{modal}},
\end{equation}
\addtocounter{equation}{-1}\refstepcounter{equation}\label{eq:ch1.ex61-torque-power}%
Eq.~\eqref{eq:ch1.ex61-torque-power} is not a new conservation theorem: it
states modal-class torque bookkeeping consistent with Eq.~\eqref{eq:ch1.ex61-torque-flux-preview} when Eq.~\eqref{eq:ch1.ex61-helicity-transport-law} is \emph{assumed} for the beam model. Subsection~\ref{sec:ch01.1.5} supplies the stress-moment balance for $\dot J_{z,\mathrm{em,out}}$ on a closed surface; it does not, without a separate Maxwell-valid mode proof, identify that flux with $(m+\sigma)P_{\mathrm{modal}}/\omega$.

For multimode radiation under the same separable assumptions,
\begin{equation}
\label{eq:ch1.ex61-multimode-torque}
\tau_{z,\mathrm{source}}
=
-\sum_{m,\sigma}\frac{m+\sigma}{\omega}\,P_{m,\sigma},
\end{equation}
where \(P_{m,\sigma}\) is radiated power carried by mode \((m,\sigma)\). Eq.~\eqref{eq:ch1.ex61-multimode-torque} is a modal specialization of Eq.~\eqref{eq:ch1.ex61-torque-flux}, not a universal identity for arbitrary full-vector fields \cite{stratton1941}.

When several \((m,\sigma)\) modes carry comparable power, the single-mode law
Eq.~\eqref{eq:ch1.ex61-torque-power} is insufficient: each mode contributes its own
\((m+\sigma)/\omega\) factor in the stated description \cite{stratton1941}.

The electromagnetic torque and the angular-momentum flux are two manifestations of the same conservation law when the material boundary is the surface across which the exchange occurs \cite{beth1936}. Beth's experiment demonstrated that circularly polarized radiation carries angular momentum whose mechanical counterpart is the torque exerted on matter through Maxwell stress at the interface \cite{beth1936}.

\begin{figure}[t]
\centering
\VolIPhysFig{\VolIPhysBethSchematic}
\caption{Beth \cite{beth1936}: circularly polarized light incident on a wave plate on a torsion suspension; the plate receives torque through Maxwell stress at the interface.}
\label{fig:ch1.beth-schematic}
\end{figure}

\begin{example}[Electromagnetic torque scale]
\label{ex:ch1.beth-torque}
Under the modal-class evaluator of Eq.~\eqref{eq:ch1.ex61-helicity-transport-law} with $m=0$, $\sigma=+1$, and $P_{\mathrm{modal}}=P=1\,\mathrm{W}$, the bookkeeping symbol gives $\dot J_z\stackrel{\mathrm{modal\ class}}{=}P/\omega$. At $f=10\,\mathrm{GHz}$, $\omega=2\pi f$ and $P/\omega\approx1.59\times10^{-11}\,\mathrm{N\,m}$. This numeric illustration is \emph{not} a Maxwell-stress integral over a uniform plane wave on a symmetric disk (Subsection~\ref{sec:ch01.6.1}); Beth-type loads require the anisotropic boundary model of the original experiment \cite{beth1936,stratton1941}.
\end{example}

On a predominantly radiative shell, outward angular-momentum flux and source reaction torque are opposite in steady harmonic operation (Eqs.~\eqref{eq:ch1.ex61-torque-flux-preview} and~\eqref{eq:ch1.outward-am-flux}) \cite{beth1936}. The limitation is equally physical: evaluating the same balance on
\(R=\lambda\) without stating \(|P_{\mathrm{reac}}(r)|/P_{\mathrm{rad}}\) as defined in Eq.~\eqref{eq:ch1.ex64-power-functionals} can
mask reactive-dominated storage exchange as if it were far-zone radiation \cite{stratton1941}.

\subsection{Phase structure without angular-momentum transport}
\label{sec:ch01.6.2}

A field may possess nontrivial angular phase structure while its net angular-momentum transport vanishes. Consider two coherent components with azimuthal orders \(+m\) and \(-m\) on an open transverse observation plane \(S_o\). The forward power and forward angular-momentum flux associated with the outgoing modal components crossing \(S_o\) are forward-propagating modal statements on that aperture; the exact global conservation law remains the closed-surface Maxwell-stress balance on an enclosing surface \(\Gamma_R\) that bounds a control volume containing the source \cite{stratton1941,belinfante1940}. Their signed modal fluxes have opposite signs; if their powers are equal, the signed orbital transport cancels even though intensity and phase remain strongly structured \cite{allen1992}. Signed orbital transport is proportional to the intensity difference \(I_+-I_-\), not to the visibility of azimuthal structure alone \cite{stratton1941}.

\begin{figure}[t]
\centering
\VolIPhysFig{\VolIPhysOpenPlaneVsClosedSphere}
\caption{Open transverse plane $S_o$ (modal flux and power through the aperture) versus closed enclosing sphere $\Gamma_R$ (exact Maxwell stress-moment conservation on a control surface) \cite{stratton1941,beth1936}.}
\label{fig:ch1.open-plane-vs-closed}
\end{figure}

With common radial envelope \(A(\rho)\), the transverse field on \(S_o\) exhibits strong azimuthal variation in both amplitude and phase. Allen \emph{et al.} \cite{allen1992} showed that a \emph{single} helical mode radiates definite orbital flux; equal counter-rotating weights remove net signed angular-momentum transport while retaining visible structure and substantial radiated power \cite{poynting1884}.

\begin{figure}[t]
\centering
\VolIPhysFig{\VolIPhysCounterRotatingAperture}
\caption{Counter-rotating components with $m=\pm2$ on open aperture $S_o$ (schematic): with matched amplitudes, $|E_t|^2\propto\cos^2(2\phi)$ shows four intensity lobes on the plane, while equal modal powers $P_+=P_-$ cancel signed forward orbital angular-momentum flux through $S_o$ ($\dot J_z=0$ on the open-surface modal bookkeeping, distinct from a closed-sphere stress integral) \cite{allen1992,poynting1884}.}
\label{fig:ch1.ex62-counter-rotating-aperture}
\end{figure}

Visible azimuthal structure does not determine net angular-momentum transport: when \(I_+=I_-\), signed flux through $S_o$ vanishes while $P_{\mathrm{tot}}=P_++P_-$ may remain large (Figure~\ref{fig:ch1.phase-vs-transport-sketch}) \cite{beth1936,stratton1941}. Component intensities follow from time-averaged Poynting flux:
\begin{equation}
\label{eq:ch1.ex62-intensity-definition}
I_\pm(\rho)=\tfrac12\Re\!\left\{\mathbf{E}_\pm\times\mathbf{H}_\pm^*\right\}\cdot\hat{\mathbf{z}},
\end{equation}
with $\mathbf{H}_\pm=(1/\eta_0)\,\hat{\mathbf{z}}\times\mathbf{E}_\pm$ on an \emph{outgoing paraxial/modal} branch (exact for a uniform plane wave; for spatially varying $\mathbf{E}_\pm$ it is a template, not a full-vector Maxwell solution unless longitudinal components and $\curl$ residuals are verified) and $\mathbf{S}_{\mathrm{c},\pm}=\tfrac12\mathbf{E}_\pm\times\mathbf{H}_\pm^*$ \cite{stratton1941,allen1992}. For
$\mathbf{E}_\pm=(A(\rho)/\sqrt{2})e^{\pm\jj m\phi}\hat{\mathbf{e}}_t$, each component carries
$I_\pm(\rho)=|A(\rho)|^2/(4\eta_0)$; matched amplitudes give equal power while signed outward fluxes cancel.

Let a monochromatic transverse field on observation surface \(S_o\) be formed by
superposing two coherent counter-rotating orbital components with common radial
envelope \(A(\rho)\), cylindrical coordinates \((\rho,\phi)\), and integer order
\(m>0\). Fix a linear transverse polarization unit vector \(\hat{\mathbf{e}}_t\). The
\(\pm m\) components are
\begin{equation}
\label{eq:ch1.ex62-counter-rotating-components}
\mathbf{E}_\pm(\rho,\phi)=\frac{A_\pm(\rho)}{\sqrt{2}}\,e^{\pm\jj m\phi}\,\hat{\mathbf{e}}_t,
\end{equation}
and their sum is
\begin{equation}
\label{eq:ch1.ex62-counter-rotating-superposition}
\mathbf{E}_t(\rho,\phi)
=
\mathbf{E}_+(\rho,\phi)+\mathbf{E}_-(\rho,\phi)
=
\sqrt{2}\,A(\rho)\cos(m\phi)\,\hat{\mathbf{e}}_t
\end{equation}
when \(A_+=A_-=A\). With \(\mathbf{H}_\pm=(1/\eta_0)\,\hat{\mathbf{z}}\times\mathbf{E}_\pm\), Eq.~\eqref{eq:ch1.ex62-intensity-definition} gives \(I_\pm=|A|^2/(4\eta_0)\) when amplitudes match. For each orbital component, the time-averaged angular-momentum flux through \(S_o\)
scales as \((\pm m/\omega)\) times component power
\(P_\pm = \int_{S_o} I_\pm\,dS\) \cite{stratton1941}. Summing components yields the open-aperture modal bookkeeping (distinct from the closed-sphere identity of Proposition~\ref{prop:ch1.single-hansen-transport})
\begin{equation}
\label{eq:ch1.ex62-channel-sum}
\dot J_{z,\mathrm{modal},S_o}
\;\stackrel{\mathrm{modal\ bookkeeping}}{=}\;
\frac{m}{\omega}(P_+ - P_-)
=
\frac{m}{\omega}\int_{S_o}(I_+-I_-)\,dS,
\end{equation}
which is Eq.~\eqref{eq:ch1.ex62-net-am-transport}. When \(I_+=I_-\), both \(P_+\) and
\(P_-\) are equal and the bookkeeping symbol vanishes even though the coherent intensity
\(|E_t|^2=|E_+|^2+|E_-|^2+2\Re(\mathbf{E}_+\cdot\mathbf{E}_-^*)\) exhibits a standing azimuthal pattern through the cross term; modal power budgets use \(P_\pm=\int I_\pm\,dS\), not \(|E_t|^2\) alone \cite{stratton1941,allen1992}.

\begin{example}[Equal $+m$ and $-m$ components]
\label{ex:ch1.counter-rotating-numeric}
Take $m=2$ with equal component powers $P_+=P_-$. Then the open-aperture bookkeeping of Eq.~\eqref{eq:ch1.ex62-net-am-transport} vanishes. A one-percent imbalance, $P_+=1.01P_-$, gives a nonzero bookkeeping symbol $(m/\omega)(0.01P_-)$: a small power asymmetry produces signed modal transport while the standing azimuthal pattern may remain visually similar \cite{allen1992}. This is not a closed-sphere Maxwell-stress evaluation.
\end{example}

Define complex amplitudes for the imbalance calculation:
\begin{equation}
\label{eq:ch1.ex62-channel-amplitudes}
E_{+}(\rho,\phi)=A_+(\rho)e^{\jj m\phi},\qquad
E_{-}(\rho,\phi)=A_-(\rho)e^{-\jj m\phi},
\end{equation}
with intensities $I_\pm(\rho)=\tfrac12\Re\{\mathbf{E}_\pm\times\mathbf{H}_\pm^*\}\cdot\hat{\mathbf{z}}$
from Eq.~\eqref{eq:ch1.ex62-intensity-definition}, or equivalently
$I_\pm(\rho)=|A_\pm(\rho)|^2/(4\eta_0)$ on a uniform outgoing branch when $|\mathbf{E}_\pm|=|A_\pm|/\sqrt{2}$. Let
\(\dot J_z\) be net angular-momentum transport rate through \(S_o\). The inherited
balance law gives
\begin{equation}
\label{eq:ch1.ex62-net-am-transport}
\dot J_z
=
\frac{m}{\omega}
\int_{S_o}\!\big(I_+-I_-\big)\,dS.
\end{equation}
Here \(m/\omega\) is the orbital angular-momentum-per-energy factor and \(I_+-I_-\) is
signed transport imbalance. Net transport is controlled by \(I_+-I_-\), not by \(|E_t|^2\) or its azimuthal derivative \cite{stratton1941}. The equal-amplitude counterexample is
\begin{equation}
\label{eq:ch1.ex62-equal-channel-condition}
I_+(\rho)=I_-(\rho)\quad\forall \rho\in S_o.
\end{equation}
Substitution into Eq.~\eqref{eq:ch1.ex62-net-am-transport} yields
\begin{equation}
\label{eq:ch1.ex62-zero-transport-result}
\dot J_z=0
\qquad\text{while}\qquad
E_t(\rho,\phi)=\sqrt{2}\,A(\rho)\cos(m\phi)\neq 0.
\end{equation}
If \(P_+\neq P_-\), a finite signed flux remains. In steady operation, section~\ref{sec:ch01.5} gives
\begin{equation}
\label{eq:ch1.ex62-zero-torque-prediction}
\tau_{z,\mathrm{source}}
=
-\dot J_z
=
0
\quad\text{under Eq.~\eqref{eq:ch1.ex62-equal-channel-condition}}.
\end{equation}
Let
\(I_+(\rho)=I_0(\rho)+\delta I(\rho)\) and
\(I_-(\rho)=I_0(\rho)-\delta I(\rho)\). Then
\begin{equation}
\label{eq:ch1.ex62-imbalance-sensitivity}
\dot J_z
=
\frac{2m}{\omega}\int_{S_o}\delta I(\rho)\,dS.
\end{equation}
\phantomsection\label{prop:ch1.ex62-phase-insufficiency}
Under Eq.~\eqref{eq:ch1.ex62-equal-channel-condition}, Eqs.~\eqref{eq:ch1.ex62-zero-transport-result} and~\eqref{eq:ch1.ex62-zero-torque-prediction} give \(\dot J_z=0\) and \(\tau_{z,\mathrm{source}}=0\) even when \(|E_t|\) and \(\partial_\phi E_t\) are large on \(S_o\) \cite{stratton1941,beth1936}. A finite signed flux requires \(I_+\neq I_-\), as in Eq.~\eqref{eq:ch1.ex62-imbalance-sensitivity}. Unequal modal powers, polarization-selective loss, incomplete modal orthogonality, or aperture truncation can break exact cancellation; gauge relabeling cannot restore transport removed by balanced superposition \cite{poynting1884,harrington2001}.
\subsection{Gauge dependence of local angular-momentum densities}
\label{sec:ch01.6.3}

Under \(\mathbf{A}'=\mathbf{A}+\nabla\chi\) and \(\phi'=\phi-\jj\omega\chi\), the fields \(\E=-\nabla\phi-\jj\omega\mathbf{A}\) and \(\mathbf{B}=\nabla\times\mathbf{A}\) are unchanged \cite{maxwell1873}. Two potential descriptions can therefore represent exactly the same electromagnetic field. Any local density constructed from \(\phi\) or \(\mathbf{A}\) alone inherits gauge sensitivity; stress traction, Poynting flux, and total angular-momentum density built from \((\E,\H)\) do not \cite{belinfante1940,stratton1941}.

An ideal oscillating electric point dipole inside a spherical control volume $V_g$ of radius $R=2\lambda$ illustrates the point. Coulomb and Lorenz gauges yield the same $(\E,\H)$ on $\partial V_g$ yet different $\nabla\phi$ at interior points, so any potential-linked local map differs pointwise between them. Mechanical torque $\oint_{\partial V_g}\mathbf{r}\times(\Tavg\cdot\hat{\mathbf{n}})\,dS$ and outward angular-momentum flux built from $(\E,\H)$ remain unchanged \cite{belinfante1940,stratton1941}.

\begin{figure}[t]
\centering
\ChOneFigBlock{%
\begin{tikzpicture}[ch1geo]
 \draw[ch1shell, dashed, line width=0.45pt] (0,0) circle (0.9);
 \node[ch1oncanvas, font=\scriptsize] at (0,1.05) {$\partial V_g$};
 \draw[gray!60!black, line width=0.55pt] (0,0.15) -- (0,0.75);
 \draw[gray!60!black, line width=0.55pt] (-0.08,0.45) -- (0.08,0.45);
 \node[ch1oncanvas, font=\scriptsize] at (0.25,0.55) {dipole};
 \draw[-{Latex[length=1.6mm,width=1mm]}, gray!40!black, line width=0.5pt]
   (0.65,0.35) arc (30:120:0.65);
 \node[ch1oncanvas, font=\scriptsize] at (0.85,0.55) {$\mathcal{L}_C$};
 \draw[-{Latex[length=1.6mm,width=1mm]}, gray!45!black, line width=0.5pt, dashed]
   (0.55,-0.25) arc (-30:-120:0.55);
 \node[ch1oncanvas, font=\scriptsize] at (0.75,-0.45) {$\mathcal{L}_L$};
 \draw[-{Latex[length=1.4mm,width=0.9mm]}, gray!50!black, line width=0.5pt]
   (0.9,0) -- (1.15,0);
 \node[ch1oncanvas, font=\scriptsize] at (1.25,0.12) {$\mathbf{M}\!\cdot\!\hat{\mathbf{n}}$};
 \node[ch1oncanvas, font=\scriptsize, align=center] at (0,-1.15)
 {same $\E$, $\H$; different local maps};
\end{tikzpicture}%
}
\caption{Ideal oscillating electric point dipole inside gauge sphere $V_g$: Coulomb and Lorenz gauges yield different gauge-dependent local maps at interior points while identical $(\E,\H)$ on $\partial V_g$ fixes the same stress-flux integral \cite{belinfante1940,stratton1941}.}
\label{fig:ch1.ex63-gauge-volume-sketch}
\end{figure}

Gauge relabeling shifts potential-based local maps; it does not shift Maxwell fields or gauge-invariant surface integrals built from \((\E,\H)\) \cite{maxwell1873,belinfante1940}. Subsection~\ref{sec:ch01.6.4} holds the same dipole fixed while moving the observation shell from reactive to radiative radius \cite{stratton1941,chu1948}.
\subsection{Radial structure of a localized radiator}
\label{sec:ch01.6.4}

What changes in the electromagnetic field as one moves away from a localized source? The source itself is fixed; only the observation radius increases. Maxwell's outgoing solution then reorganizes its radial content: terms that decay as $1/r^3$ and $1/r^2$ dominate near the source, while the $1/r$ radiative term eventually sets the transverse far-zone form \cite{stratton1941,poynting1884}. Consider a $z$-directed electric dipole moment phasor $\mathbf{p}=p_0\hat{\mathbf{z}}$ at the origin with $\exp(+\jj\omega t)$ and outgoing factor $e^{-\jj kr}$. In spherical coordinates $(r,\theta,\phi)$ the exact source-free fields are \cite{stratton1941,harrington2001}
\begin{subequations}
\label{eq:ch1.dipole-fields-explicit}
\begin{align}
E_r&=\frac{p_0\cos\theta}{4\pi\varepsilon_0}\,e^{-\jj kr}
\left(\frac{2}{r^3}+\frac{2\jj k}{r^2}\right),
\label{eq:ch1.dipole-Er}\\[0.25em]
E_\theta&=\frac{p_0\sin\theta}{4\pi\varepsilon_0}\,e^{-\jj kr}
\left(\frac{1}{r^3}+\frac{\jj k}{r^2}-\frac{k^2}{r}\right),
\label{eq:ch1.dipole-Etheta}\\[0.25em]
H_\phi&=\frac{\jj k p_0\sin\theta}{4\pi\varepsilon_0\eta_0}\,e^{-\jj kr}
\left(\frac{1}{r^2}+\frac{\jj k}{r}\right),
\label{eq:ch1.dipole-Hphi}
\end{align}
\end{subequations}
with $k=\omega/c$. For a $z$-directed dipole moment $\mathbf{p}=p_0\hat{\mathbf{z}}$, the retarded Lorenz-gauge vector potential in the source-free exterior ($r>0$) is
\begin{equation}
\label{eq:ch1.dipole-A-retarded}
\mathbf{A}(\mathbf{r})=\jj\omega\mu_0\,\frac{p_0\hat{\mathbf{z}}\,e^{-\jj kr}}{4\pi r},
\qquad r>0,
\end{equation}
consistent with $\mathbf{A}(\mathbf{r})=\mu_0\int \Jimp(\mathbf{r}')\,G(\mathbf{r}-\mathbf{r}')\,d^{3}r'$ and the chapter Green function \cite{stratton1941,harrington2001,orfanidis2014}. Because $\E=-\jj\omega\mathbf{A}-\nabla\phi$ and $\H=(1/\mu_0)\curl\mathbf{A}$, each spatial derivative of $e^{-\jj kr}/r$ lowers the power of $r$ and, when acting on the oscillatory factor, introduces powers of $\jj k$. Collecting terms produces the $r^{-3}$, $r^{-2}$, and $r^{-1}$ hierarchy displayed above: the $r^{-3}$ part matches the static dipole field; the $r^{-2}$ terms are the induction fields; the $-k^2/r$ term in $E_\theta$ together with $H_\phi\propto k/r$ is the radiative branch \cite{stratton1941}. At fixed $kr$, several orders can be comparable, so ``near,'' ``transition,'' and ``far'' describe dominant contributions rather than non-overlapping shells \cite{stratton1941,poynting1884}.

\begin{figure}[t]
\centering
\VolIPhysFig{\VolIPhysDipoleThreeZones}
\caption{Concentric observation shells: near zone ($kr\ll1$, $r^{-3}$ and $r^{-2}$), transition ($kr\sim1$), and far zone ($kr\gg1$, $r^{-1}$ radiation) \cite{stratton1941,poynting1884}.}
\label{fig:ch1.dipole-three-zones}
\end{figure}

\begin{example}[Reading the three regions of a dipole field]
\label{ex:ch1.dipole-kr}
At $f=300\,\mathrm{MHz}$, $k=2\pi/\lambda\approx6.3\,\mathrm{rad/m}$. At $r=0.5\,\mathrm{m}$, $kr\approx3.1$ (transition); at $r=5\,\mathrm{m}$, $kr\approx31\gg1$ and the $r^{-1}$ terms in Eqs.~\eqref{eq:ch1.dipole-Etheta} and~\eqref{eq:ch1.dipole-Hphi} dominate, giving transverse radiation with $U(\theta)\propto\sin^2\theta$, $P_{\mathrm{rad}}$ from Eq.~\eqref{eq:ch1.ex64-power-invariance}, and $D(\theta)=\tfrac32\sin^2\theta$ \cite{stratton1941,harrington2001}.
\end{example}

For $r>0$, $\nabla\cdot\E=0$ follows from the spherical-coordinate divergence condition
\[
\partial_r E_r+\frac{2}{r}E_r+\frac{1}{r}\partial_\theta E_\theta+\frac{\cot\theta}{r}E_\theta=0
\]
when Eqs.~\eqref{eq:ch1.dipole-Er} and~\eqref{eq:ch1.dipole-Etheta} are substituted \cite{stratton1941}. For $r>0$, the displayed triad satisfies $\nabla\cdot\E=0$ and $\nabla\cdot\H=0$ with $\H=H_\phi\hat{\boldsymbol\phi}$ independent of $\phi$, so only $\partial_\phi H_\phi$ enters the spherical divergence and vanishes \cite{stratton1941,harrington2001}. The curl identities equivalent to $\curl\E=-\jj\omega\mu_0\H$ and $\curl\H=+\jj\omega\varepsilon_0\E$ hold with $\Jimp=\mathbf{0}$; the radial curl component $(\curl\H)_r=(r\sin\theta)^{-1}\partial_\theta(\sin\theta H_\phi)$ balances $\jj\omega\varepsilon_0 E_r$ even though $H_r=H_\theta=0$ \cite{stratton1941,harrington2001}. At $r=0$ the ideal point dipole requires distributional sources. For a $z$-directed electric dipole moment phasor $\mathbf{p}=p_0\hat{\mathbf{z}}$ ($p_0$ in C\,m), the charge density in the distributional sense is
\begin{equation}
\label{eq:ch1.dipole-rho-distributional}
\rho(\mathbf{r})=-\mathbf{p}\cdot\nabla\delta^{3}(\mathbf{r})
=-p_0\,\partial_z\delta^{3}(\mathbf{r}),
\end{equation}
so that $\int f(\mathbf{r})\,\rho(\mathbf{r})\,d^{3}r=+p_0\,\partial_z f|_{\mathbf{r}=\mathbf{0}}$ for smooth compactly supported test functions $f$ (integration by parts on $\partial_z\delta^{3}$) \cite{jackson1999,stratton1941}. Gauss's law $\divg\E=\rho/\varepsilon_0$ then requires $\divg\E=-(p_0/\varepsilon_0)\,\partial_z\delta^{3}$ at the origin. The compatible impressed current is
\begin{equation}
\label{eq:ch1.dipole-Jimp-distributional}
\Jimp(\mathbf{r})=\jj\omega p_0\,\hat{\mathbf{z}}\,\delta^{3}(\mathbf{r}),
\end{equation}
because $\divg\Jimp+\jj\omega\rho=0$ with $\divg\big(\hat{\mathbf{z}}\,\delta^{3}\big)=\partial_z\delta^{3}$ in $\mathcal{D}'(\mathbb{R}^{3})$ \cite{jackson1999,stratton1941}. With outgoing $G(\mathbf{r})=e^{-\jj kr}/(4\pi r)$ and $(\nabla^2+k^2)G=-\delta^{3}(\mathbf{r})$ in $\mathcal{D}'(\mathbb{R}^{3})$, the Lorenz potentials $\mathbf{A}=\mu_0\int \Jimp G\,d^{3}r'$ and $\phi=-\epsilon_0^{-1}\int\rho G\,d^{3}r'$ yield the Hertzian field pair
\begin{equation}
\label{eq:ch1.dipole-E-dyadic-distributional}
\E=\epsilon_0^{-1}\left[k^{2}\mathbf{p}\,G+(\nabla\nabla G)\cdot\mathbf{p}\right],
\qquad
\H=\jj\omega\,\nabla G\times\mathbf{p},
\end{equation}
which reduce to Eqs.~\eqref{eq:ch1.dipole-Er}--\eqref{eq:ch1.dipole-Hphi} for $r>0$ \cite{stratton1941,harrington2001,yaghjian1982}. The isotropic contact term in the distributional Hessian of $1/(4\pi r)$,
\[
\partial_i\partial_j\frac{1}{4\pi r}
=\operatorname{PV}\!\frac{3\hat r_i\hat r_j-\delta_{ij}}{4\pi r^{3}}
-\tfrac13\delta_{ij}\,\delta^{3}(\mathbf{r}),
\]
under the spherical principal-volume prescription used here, combines with the $k^{2}G\mathbf{p}$ term so that no spurious $\delta^{3}$ remains in $\divg\E$ beyond Eq.~\eqref{eq:ch1.dipole-rho-distributional} \cite{yaghjian1980dyadic,yaghjian1982}. For smooth compactly supported $\psi$, distributional Gauss's law follows from $(\nabla^2+k^2)G=-\delta^{3}$:
\[
\divg\E=\epsilon_0^{-1}\mathbf{p}\!\cdot\!\nabla\big[(\nabla^2+k^2)G\big]
=-\epsilon_0^{-1}\mathbf{p}\!\cdot\!\nabla\delta^{3},
\]
so
\[
\langle\divg\E,\psi\rangle
=-\epsilon_0^{-1}\sum_i p_i\langle\partial_i\delta^{3},\psi\rangle
=+\epsilon_0^{-1}\mathbf{p}\!\cdot\!\nabla\psi(\mathbf{0})
=\langle\rho/\epsilon_0,\psi\rangle,
\]
using $\langle\partial_i\delta^{3},\psi\rangle=-\partial_i\psi(\mathbf{0})$ and Eq.~\eqref{eq:ch1.dipole-rho-distributional} \cite{yaghjian1982,stratton1941}. Magnetic Gauss's law follows from $\H\propto\nabla G\times\mathbf{p}$ and $\divg(\nabla G\times\mathbf{p})=0$ distributionally; Faraday's law from $\curl$ of a gradient term plus $k^2/\epsilon_0=\omega^2\mu_0$; Amp\`ere--Maxwell uses the constant-$\mathbf{p}$ identity $\curl(\nabla G\times\mathbf{p})=-\mathbf{p}(\nabla^2 G)+(\mathbf{p}\!\cdot\!\nabla)\nabla G$, then $\nabla^2 G=-k^2G-\delta^{3}$, giving $\curl\H=\Jimp+\jj\omega\epsilon_0\E$ with exactly the delta current of Eq.~\eqref{eq:ch1.dipole-Jimp-distributional} \cite{yaghjian1982,harrington2001}. Eq.~\eqref{eq:ch1.dipole-maxwell-distributional-closure} collects this closure; quadratic Maxwell stress at $\mathbf{r}=\mathbf{0}$ is not asserted well-defined---only the linear Maxwell system in $\mathcal{D}'(\mathbb{R}^{3})$ \cite{yaghjian1980dyadic,stratton1941}.
\begin{equation}
\label{eq:ch1.dipole-maxwell-distributional-closure}
\divg(\epsilon_0\E)=\rho,\quad \divg\mathbf{B}=0,\quad \curl\E=-\jj\omega\mathbf{B},\quad \curl\H=\Jimp+\jj\omega\epsilon_0\E
\quad\text{in }\mathcal{D}'(\mathbb{R}^{3})
\end{equation}
for the source pair Eqs.~\eqref{eq:ch1.dipole-rho-distributional}--\eqref{eq:ch1.dipole-Jimp-distributional} and fields Eq.~\eqref{eq:ch1.dipole-E-dyadic-distributional}.
Finite-width regularizations of $\mathbf{p}$ approach the same exterior fields when $k\ell\ll1$ for source scale $\ell$ \cite{stratton1941,harrington2001,yaghjian1982}. The $E_r$--$E_\theta$--$H_\phi$ split is the outgoing Hertzian template with the published coefficients \cite{stratton1941,harrington2001}. The same ideal point dipole at
$f=1\,\mathrm{GHz}$ used in Subsection~\ref{sec:ch01.6.3} illustrates shell dependence on
two concentric spheres $S_{r_1}$ and $S_{r_2}$ with $kr_1=0.5$ and $kr_2=5$, as in Figure~\ref{fig:ch1.radial-field-zones}.

In the standard multipole description of an outgoing field,
\begin{equation}
\label{eq:ch1.ex64-multipole-radial}
\mathbf{E}(\mathbf{r})=\mathbf{E}_3(\theta,\phi)\,r^{-3}+\mathbf{E}_2(\theta,\phi)\,r^{-2}+\mathbf{E}_1(\theta,\phi)\,r^{-1}+\cdots,
\qquad
\mathbf{H}(\mathbf{r})=\mathbf{H}_2(\theta,\phi)\,r^{-2}+\mathbf{H}_1(\theta,\phi)\,r^{-1}+\cdots,
\end{equation}
where the angular factors $\mathbf{E}_n,\mathbf{H}_n$ depend on direction but not on $r$ in the source-free exterior \cite{stratton1941}. Near the source the $r^{-3}$ and $r^{-2}$ terms may be comparable with, or larger than, the $r^{-1}$ radiation field. Those terms describe stored and induction fields: electric and magnetic energy oscillate locally, and the normal Poynting vector can be largely imaginary even though the source already radiates a definite outward power. At increasing distance the $r^{-1}$ terms dominate, $\E$ and $\H$ approach the transverse relation of an outgoing wave, and the real Poynting flux carries energy away from the source \cite{poynting1884}.

The radial intensity of that flux follows from the field itself. With
$\mathbf{S}_\mathrm{c}=\tfrac12\E\times\H^*$ and $H_\phi$ from
Eq.~\eqref{eq:ch1.dipole-Hphi} (within Eq.~\eqref{eq:ch1.dipole-fields-explicit}), the radial component is
\begin{equation}
\label{eq:ch1.dipole-sc-radial}
\Re\!\left\{(\mathbf{S}_\mathrm{c}\cdot\hat{\mathbf{r}})\right\}
=\tfrac12\Re\!\left\{E_\theta H_\phi^*\right\},
\end{equation}
since $\H$ has only a $\phi$-component \cite{stratton1941}.
On a sphere of radius $r$, the radiative part of $\Re(S_{\mathrm{c},r})$ scales as $1/r^2$ because both $\E$ and $\H$ carry the $1/r$ factor. The element of area grows as $r^2$, so
\begin{equation}
\label{eq:ch1.ex64-power-invariance}
P_{\mathrm{rad}}=\oint_{S_r}\Re\!\left\{\mathbf{S}_\mathrm{c}\cdot\hat{\mathbf{r}}\right\}dS
=\frac{\omega^4|p_0|^2}{12\pi\varepsilon_0 c^3},
\end{equation}
remains constant in a lossless source-free region between enclosing surfaces \cite{stratton1941}. The source has radiated the same total power all along; what changes with $r$ is the local composition of the field through which that conserved power is carried.

\begin{figure}[t]
\centering
\VolIPhysFig{\VolIPhysDipoleRadiationPattern}
\caption{Far-zone radiation intensity $U(\theta,\phi)\propto\sin^2\theta$ for an electric dipole aligned with $\hat{\mathbf{z}}$: intensity vanishes on axis and peaks in the equatorial plane. With Eq.~\eqref{eq:ch1.ex64-power-invariance}, $D(\theta)=\tfrac32\sin^2\theta$ and $D_0=\tfrac32$ \cite{stratton1941,harrington2001}.}
\label{fig:ch1.dipole-radiation-pattern}
\end{figure}

Figure~\ref{fig:ch1.dipole-radiation-pattern} is the angular face of Eq.~\eqref{eq:ch1.ex64-power-invariance}: $P_{\mathrm{rad}}$ is fixed, but its solid-angle distribution is strongly non-uniform. Axisymmetric radiation carries finite power yet may carry zero net $z$-angular-momentum flux when the stress integrand is symmetric about $\hat{\mathbf{z}}$ \cite{beth1936}.

The complex normal flux separates reactive circulation from outward radiation. With Eq.~\eqref{eq:ch1.ex64-power-functionals},
\begin{equation}
\label{eq:ch1.ex64-complex-power}
\int_{S_r}\mathbf{S}_\mathrm{c}\cdot\hat{\mathbf{r}}\,dS=P_{\mathrm{rad}}+\jj P_{\mathrm{reac}}(r),
\end{equation}
where $P_{\mathrm{rad}}$ is the time-averaged outward power and $P_{\mathrm{reac}}(r)=\int_{S_r}\Im\{\mathbf{S}_\mathrm{c}\cdot\hat{\mathbf{r}}\}\,dS$ is reactive power through the shell \cite{poynting1884}. Because $P_{\mathrm{rad}}$ is fixed in a lossless exterior while $|P_{\mathrm{reac}}(r)|$ varies with shell radius, the ratio $|P_{\mathrm{reac}}(r)|/P_{\mathrm{rad}}$ marks reactive dominance when it is much greater than unity and radiation dominance when it is much less than unity \cite{stratton1941}.

On $S_{r_1}$ with $kr_1=0.5$, reactive circulation dominates the normal flux pattern even though $P_{\mathrm{rad}}$ is unchanged. Integrating Eqs.~\eqref{eq:ch1.dipole-Etheta} and \eqref{eq:ch1.dipole-Hphi} on $S_{r_1}$ via Eq.~\eqref{eq:ch1.ex64-power-functionals} with $|p_0|=1\,\mathrm{C\,m}$ gives $|P_{\mathrm{reac}}(r_1)|/P_{\mathrm{rad}}\approx 8$. On $S_{r_2}$ with $kr_2=5$, the same integration gives $|P_{\mathrm{reac}}(r_2)|/P_{\mathrm{rad}}\approx 8\times 10^{-3}$, and $\E$--$\H$ approach the far-zone transverse relation
under the outgoing branch and radiation condition stated in section~\ref{sec:ch01.1} \cite{stratton1941}.

Angular-momentum transport must be read with the same distinction. Outward field angular-momentum flux through $S_r$ is
\begin{equation}
\label{eq:ch1.ex64-am-flux-integral}
\dot J_{z,\mathrm{em,out}}
=-\oint_{S_r}\big[\mathbf{r}\times(\Tavg\cdot\hat{\mathbf{r}})\big]\cdot\hat{\mathbf{z}}\,dS
\end{equation}
by Eqs.~\eqref{eq:ch1.outward-am-flux} and~\eqref{eq:ch1.stress-moment-integral}, and is conserved on nested source-free spheres in steady state when the same outgoing Maxwell branch is followed \cite{stratton1941}. Moving the enclosing surface outward does not create angular momentum; it changes the local stress distribution through which the same signed flux is carried. By contrast, the integral of the magnitude $\int_{S_r}|[\mathbf{r}\times(\Tavg\cdot\hat{\mathbf{r}})]\cdot\hat{\mathbf{z}}|\,dS$ may vary with $r$ because near-zone reactive structure contributes differently to the stress moment than the far-zone radiative field---without implying growth of the conserved net flux \cite{stratton1941,beth1936}.

Chu's quality factor for a control volume $V$ enclosing the source is
\begin{equation}
\label{eq:ch1.chu-q-definition}
Q_{\mathrm{Chu}}(V)=\frac{2\omega\max(W_e,W_m)}{P_{\mathrm{rad}}},
\end{equation}
with cycle-averaged stored energies $W_e$, $W_m$ and outward radiated power $P_{\mathrm{rad}}$ \cite{chu1948}. At resonance, $Q_{\mathrm{Chu}}=\omega(W_e+W_m)/P_{\mathrm{rad}}$. The shell integral $P_{\mathrm{reac}}(r)$ is reactive power through $S_r$, not $Q_{\mathrm{Chu}}$ \cite{chu1948,stratton1941}. Eq.~\eqref{eq:ch1.chu-q-lower-bound} applies to Chu's idealized minimum-$Q$ spherical-wave radiator of radius $a$ in free space ($ka\ll 1$ in the electrically small limit); it is not a universal lower bound on every antenna geometry without restating the enclosing volume, stored-energy definition, and mode content \cite{chu1948}. Under that model,
\begin{equation}
\label{eq:ch1.chu-q-lower-bound}
Q_{\mathrm{LB}}=\frac{1}{(ka)^3}+\frac{1}{ka}.
\end{equation}
Local $r^{-3},r^{-2},r^{-1}$ weighting changes with radius; $P_{\mathrm{rad}}$ and nested-sphere $\boldsymbol{\Phi}_{J,\mathrm{out}}$ do not \cite{poynting1884,stratton1941}. Chapter~\ref{ch:02} develops angular-spectrum mappings for bounded sources \cite{harrington2001}.

\section{Summary}
\label{sec:ch01.7}

Within the harmonic Maxwell class developed here, compatible impressed sources, interface and boundary data, and outgoing selection define the field problem whose flux integrals are physically meaningful; uniqueness and existence require the hypotheses of Subsection~\ref{sec:ch01.1.3} \cite{maxwell1873,stratton1941}. From that closure follow Poynting's theorem, the Maxwell stress tensor, and angular-momentum transport $\boldsymbol{\Phi}_{J,\mathrm{out}}=-\boldsymbol{\mathcal{T}}_J$ under the sign convention fixed in Subsection~\ref{sec:ch01.1.5} \cite{poynting1884,beth1936}. Chapter~\ref{ch:02} develops angular-spectrum and dyadic-Green operators on the same outgoing branch \cite{harrington2001,tai1994}.

\phantomsection\label{sec:ch01.7.1}
\phantomsection\label{eq:ch1.continuity-admissibility}%
\phantomsection\label{eq:ch1.continuity-charge-continuous}%
\phantomsection\label{eq:ch1.master-admissibility-chain}%
\phantomsection\label{eq:ch1.master-charge-continuous-chain}%
\phantomsection\label{eq:ch1.report-class-lock}%
\phantomsection\label{sec:ch01.7.2}
\phantomsection\label{sec:ch01.7.3}
\phantomsection\label{eq:ch1.bridge-compatibility-criteria}%
\phantomsection\label{eq:ch1.compatibility-criteria}%

